\documentclass[11pt,a4paper]{article}
\usepackage[margin=26.5mm,includefoot]{geometry}
\usepackage{graphicx,amsmath,booktabs,xurl,textcomp}
\usepackage{fontspec}
\usepackage[mathrm=sym]{unicode-math}
\setmainfont{NotoSans-Regular.ttf}[Path=../fonts/,BoldFont=NotoSans-Bold.ttf,ItalicFont=NotoSans-Italic.ttf,BoldItalicFont=NotoSans-BoldItalic.ttf]
\setsansfont{NotoSans-Regular.ttf}[Path=../fonts/,BoldFont=NotoSans-Bold.ttf,ItalicFont=NotoSans-Italic.ttf,BoldItalicFont=NotoSans-BoldItalic.ttf]
\setmonofont{NotoSansMono-Regular.ttf}[Path=../fonts/,BoldFont=NotoSansMono-Bold.ttf]
\setmathfont{FiraMath-Regular.otf}[Path=../fonts/]
\usepackage[numbers,sort&compress]{natbib}
\setlength{\bibsep}{7pt}
\usepackage[colorlinks=true,linkcolor=blue,citecolor=blue,urlcolor=blue]{hyperref}
\emergencystretch=3em
\setlength{\parskip}{4pt}
\setlength{\parindent}{0pt}
\setcounter{tocdepth}{1}
\hypersetup{pdftitle={NeuroTranslation: A New Paradigm In Brain-Body Mapping.},pdfauthor={@eliminatenoise},pdfsubject={Unpublished Working Draft; source and rendered-document review pending}}
% Original sans-serif master retained; new isolated source additions are unrendered.
\usepackage{ragged2e}
\usepackage[justification=raggedright,singlelinecheck=false]{caption}
\AtBeginDocument{\RaggedRight}
% Keep paragraph starts/ends together; no font, content or scientific change.
\makeatletter
\AtBeginDocument{\clubpenalty=10000\@clubpenalty=10000\widowpenalty=10000\displaywidowpenalty=10000}
\makeatother
\renewcommand{\bibfont}{\RaggedRight}
\AddToHook{cmd/LARGE/after}{\hyphenpenalty=10000\exhyphenpenalty=10000}
\renewcommand{\_}{\char95\relax}
\begin{document}
\title{NeuroTranslation: A New Paradigm In Brain-Body Mapping.\\[.4em]\large Simple English Edition}
\author{@eliminatenoise}
\date{Unpublished Working Draft\\Original manuscript date: October 13, 2026\\Source assembly: November 5, 2026\\Typeset \today}
\maketitle

\begin{abstract}
Both larger models improved the average. Both failed the test. The latest LINK study uses six recording sessions from one macaque to predict the positions of two finger groups, expressed on the study's normalized scale. Four sessions supply the training examples; two supply the development-validation examples used to compare the models. The simpler reference models combine neural samples using weights---numbers fitted from training data. The two larger candidates read a sequence while updating a software memory of what came before. That numerical memory is not a measurement of biological memory. The comparison uses three fixed random starting conditions. Across 18 runs and 720 saved epochs, the two history-reading candidates lower average error. An epoch is one pass through all the training examples allowed for that run.

The original requirements must hold together. No evaluated output may get worse than the selected weighted-sum reference: a worse recording-day/finger-group remains a failure, however good the average looks. Every group on every validation day must also have a defined, positive ordinary $R^2$ score. That means doing better, under the stated squared-error score, than always predicting that group/day's observed mean. Each score must also exceed its matching association control, which challenges the pairing between neural observations and movement. Those checks give the requirement for positive neural information its meaning here. All five added benefit claims fail too. Neither model is admitted by these development results. The models share a fitting update rule---the procedure that changes their adjustable numbers during training. These results do not independently verify that rule or show how the models perform on the reserved confirmation recordings.

NeuroTranslation asks what we can establish by connecting a stated neural recording to a stated body target, then checking the prediction and each further claim against the evidence it needs. Its software, NeuroTranslator, works on selected, recording-specific problems using models, replay and body interfaces. The larger NeuroTranslation Network remains proposed. Earlier reaching, sensory, attempted-finger and simulated-hand tests answer different questions and remain separate from LINK. A predicted movement is not a learned skill. A simulated hand can carry out the issued command even when the neural decision behind it is wrong; a goal-driven hand can complete a phrase without recorded brain activity driving that policy.

The explanation follows what the sources mean, their units and clocks, model selection and uncertainty, missing pathways, anatomy, simulation, interface limits and the evidence needed for learning that lasts. The proposed next comparison would pair brain and muscle observations with independently measured movement. Physical calibration, clocks, contact identities and participant registration still need resolving. It would test the declared signal-to-movement mapping against that independent reference, then distinguish help during performance from learning that remains after the help ends.

Anatomical and functional comparisons, final-checkpoint and update verification, and complete GUI and offline checks remain unfinished. Predicting movement, delivering assistance and retaining a biological skill each require their own evidence. Biological teaching and universal skill transfer remain unestablished; neither instant biological mastery nor a complete nervous-system map follows from the reported results.
\end{abstract}
\tableofcontents\clearpage
\section{Neural Signals and Physical Skills}

A file can store a model. What would make it store a skill? Knowing what a hand
should do, predicting its movement and making it move are already different
problems. A machine might put the hand on the right path while doing the decisive
work itself. A file offered as another person's skill therefore owes us an
explanation: what does it contain, what would give those contents meaning in a
different body, and what can the trainee still do when the help ends?

Those questions motivate NeuroTranslation. The proposed field maps nerve signals
and patterns that describe them between clearly stated situations: a particular
recording, a particular body and a particular task. The implemented software,
NeuroTranslator, joins selected recordings, fitted prediction models, playback
and computer-body checks. The NeuroTranslation Network is a proposed plan for
joining more components while keeping the origin and meaning of each value
clear. Brain2Body, with the numeral meaning ``to'', names the larger programme.
A field, a piece of working software and a proposed design are three different
things. Their names do not turn the current work into a universal decoder.

The ambition is faster physical learning that remains after teaching stops.
The current work is smaller and more testable. It connects real recorded neural
activity to specific targets, exposes errors and tries bounded commands in
body simulations. Its latest central neural comparison has a negative outcome.
Two models lowered average error on development recordings, but both failed
the original requirements. The average improved while declared parts of the
prediction got worse. That failure tells us where the next explanation has to
begin; the older results cannot answer for it.

The first question is concrete: can specified neural observations predict an
independently measured body target? The title does not mean that an established
scientific paradigm has been replaced or that we invented the familiar models
and simulators used here. Every claimed link must answer
to a measurement: identify what was recorded, what was calculated, which body
it refers to and what happened, then state which comparison could defeat the
claim. These conditions make the link testable.

This Simple English Edition explains the concepts from the beginning and
stands on its own. The companion Brain2Body article separately introduces the
larger biological-interface proposal; neither manuscript treats a presentation
or demonstration as scientific authority.\cite{brain2bodyworking} Both remain
unpublished working manuscripts. The historical figures retain their original
evidence scopes: relabelling a figure cannot turn it into a new experiment.

The \hyperref[programme-definitions]{definition supplement} gives this edition's
shared registry version and fingerprint; older studies retain the versions they
used. The manuscript's scientific and publication limits are stated in its edition note;
none of the results reported here establishes retained biological learning.

\section{Brain Structure and Brain Activity}

A road map tells you where roads connect. It does not tell you which roads are
busy right now, which driver will turn, how fast a car can stop or whether a
bridge is closed. A neural connection map has a similar limit. It records
structure through a particular measurement and reconstruction method. Structure
constrains what can happen, but does not by itself determine all the activity,
chemical influences, changing cell states or learning processes that produce
behaviour.

FlyWire is a major anatomical reconstruction of an adult female fly brain.
Its measured connections are useful evidence, with acknowledged reconstruction
and coverage limits. The mapped brain does not include the complete ventral
nerve cord, the fly's major lower nerve structure. The dataset is not itself a
behaving fly simulator.\cite{flywire2024} A developer can use it to constrain
software, but still has to choose or measure other ingredients. If those
choices include a body controller, the controller's contribution must be
visible. It cannot be credited to the wiring diagram alone.

This distinction also applies to our own software. A human-arm simulation is
not made into the recorded macaque by placing a brain beside it. A model can
predict a measured finger-group position without recovering every muscle and
nerve that produced that position. A shared display is useful for comparison,
but does not establish a shared biological cause.

An electrode measures activity under particular acquisition conditions.
A contact can combine contributions from more than one cell, and its identifier
belongs to that recorded inventory. Contact 12 on one date does not establish
that the same cells supplied contact 12 on another date. Our current 192 input
features come from two measurements at 96 physical contacts. Calling them 192
neurons would change the claim without adding evidence. The recording gives us
an observation of the system, with limits set by how it was acquired.

Other instruments observe other quantities. EEG means electrical activity
recorded from the scalp. ESG means electrical activity recorded at the spinal
surface. ENG concerns nerve electrical activity. EMG concerns muscle electrical
activity. These labels help distinguish the observations. EMG is not a direct
peripheral nerve recording. Even when a signal helps predict movement, the
recording method still limits what has been observed.

Functional magnetic-resonance imaging, usually shortened to fMRI, measures an
indirect signal associated with blood oxygen. It does not photograph thought.
A classic study recording electrical activity and this imaging signal together
found important dependencies on the observed electrical quantity, animal and
site.\cite{logothetis2001} A colourful brain image needs an explanation of what
the colours represent, the time support and the measurement limits before it
can support a functional claim.

Four meanings of ``large'' or ``small'' must remain separate. Physical scale
means measured distances or sizes, with units. Organizational level means a
whole body, organ, local circuit or cell. Computational detail means what one
software calculation represents. Statistical averaging means how we combine
numbers. Averaging two finger groups does not measure a larger piece of brain.
A small software model does not mean that it observes individual cells. More
features do not automatically mean more physical sites.

The useful question is therefore specific: given this recorded quantity from
this subject, task and interval, can this method predict this target under
these checks? That question can produce evidence. Asking whether a colourful
model ``understands the brain'' without naming a target makes both success
and failure much harder to judge.

\section{Measure Movement}

A person experiences a sensation. Another observer can ask about it, but a
report remains a report given by that person in a particular context. The
body's position, path, timing, contact and force can be checked with independent
instruments. Those measurements have uncertainty too, but they supply observable
targets against which a model can be wrong.

The sensory-action relationship gives this programme a practical starting
point. Sensory input informs action. Action changes the body and its environment.
The resulting feedback informs the next action. A useful test can compare what
was intended or predicted with what was measured, then ask whether a correction
improves the next result. Established feedback-control work describes how
movement can be adjusted relative to task goals under uncertainty.\cite{todorov2002}
Our programme uses that rationale; it has not thereby demonstrated every
biological loop it proposes.

Consider reaching for a cup. A target position, measured hand path and contact
outcome could make the task checkable. A decoder might predict the path while a
controller changes the muscle commands. The path could improve because the
controller is better, because the decoder is better, because an easier target
was selected or because an answer was accidentally fed into the controller.
A fair comparison keeps those explanations separate. Independent measurements,
matched targets and explicit information access make the interpretation possible.

A saved recording is not live feedback. It can show what happened in an earlier
experiment, and an offline model can predict part of it. A computer body can
provide simulated feedback. A living body can provide measured feedback only
through an actual synchronized acquisition. Playing all three on one screen
does not give them one clock or one nervous system.

Language offers a useful analogy. The meaning of a statement can depend on
who said it, the setting and what happened before it. Likewise, a neural pattern
needs its measurement and task context. Published semantic reconstruction from
brain measurements depended on the participant and cooperation; it did not
establish unique access to any person's private experience.\cite{tang2023}
This does not prove that useful shared representations are impossible. It
means their preserved information and ambiguity need to be tested rather than
assumed away.

Quantum measurement offers a much narrower analogy: observation involves an
interaction and its conditions matter. That analogy is not a mechanism for
ordinary brain measurements. Quantum backaction means that a measurement changes the state of the quantum system being measured; macroscopic neural recordings, which combine signals over much larger physical
scales, have different explanations. We do not
claim that quarks (subatomic particles), or quantum disturbance, explain our decoding limits.\cite{quantum2012}
An analogy should help pose a question, not smuggle in a scientific explanation
that the experiment never tested.

The programme borrows a discipline from code and mathematics too: state the
question, specify what would meet it and check the answer. Biology remains an
empirical subject. A test cannot replace a measurement, and a measurement does
not become error-free because a program accepts it. A one-to-one decoding of
private subjective experience is not a result here.

\section{Test Predictions Against Measurements}

Suppose a model predicts the right endpoint once. That is a useful observation,
but it leaves several questions open. Did it use the relevant input? Would it
work on another trial? Was the endpoint easy to guess? Does it describe the
mechanism, or only one consequence? Is the error acceptable in the task's units?
A single agreement cannot answer all of those questions.

Accuracy means agreement with an independent reference under stated conditions.
Precision means agreement across repeated measurements under stated conditions.
A model that repeatedly places a hand too far left can be repeatable and still
wrong. Printing more decimal places does not correct it. Neither does displaying
a finer mesh or shortening the simulator's numerical step.\cite{nist1297}

Ground truth here means an independently acquired empirical reference with
its source, calibration, sampling limits and uncertainty exposed. It is not
claimed to be a perfect view of biology. A finger sensor, a camera-derived
joint estimate and a force transducer have different meanings and limits.
When a target comes from a body-analysis model, it must be labelled as derived
rather than a directly measured living muscle quantity.

A disclosed approximation is a hypothesis that can be useful if it passes the
right tests. Biological virtualization would mean a digital model independently
shown to reproduce specified aspects of a living system's behaviour and causal
dynamics within stated error bounds. That is a demanding evidence claim.
A convincing animation does not earn it. Invented biological content presented
as faithful reproduction is what this programme calls hallucination.

The dice analogy concerns unknown reliability. If a biological claim lacks the
necessary tests, a matching answer leaves us unsure whether the model is reliable
or merely agreeable in that case. This does not say every prediction is random,
that all answers have equal odds or that all approximations are false. Learned
patterns can be informative. Independent error measurement, failure reporting,
held-out tests and assessed uncertainty tell us how much reliance is justified.

Generated examples can help fit a model or test an implementation. They cannot
create independent empirical validation for an absent biological measurement.
If a simulator generates the targets and the model matches them, that tests the
model relative to the simulator. The next biological claim still needs a real
reference. A larger pile of generated agreement cannot supply the missing
measurement by counting the same assumptions many times.


\section{Earlier Results}

The earlier studies answer different questions. Their recordings, targets and tests
stay separate; these accounts add no new confirmation. The current LINK cohort
and its completed negative comparison follow.

The prototype executes a declared chain: recorded signals enter models matched to each recording, and their outputs pass through explicit body and replay interfaces. That is the idea tested at this bounded scale. The inputs come from named public experiments; saved model states and reported outputs identify what a reproduction should check. The body simulations use named numerical models. These are computer experiments, with failures preserved in the evidence record, rather than physical experiments on an organism.

Consider the simplest chain. A public experiment recorded electrical events in a macaque's brain while it reached. We count events in completed short time intervals, called bins, preserve the original recording-channel order, and give a short history to a trained model. Such a model is sometimes called a decoder because it turns a recorded signal into a prediction about another quantity. This model returns two fingertip velocity components, meaning signed rates of movement along the two recorded coordinate directions, in meters per second. The replay view places those predictions beside the appropriate reference quantities. A separate translation rule turns velocity into a deliberately limited shoulder-and-elbow command for a human arm model. That last translation is an engineering choice. It is not a discovery that the macaque's recorded neurons correspond to particular human muscles. The reaching archive supplies the recorded neural-behavior relationship; the body experiment supplies a numerical test after a stated conversion between bodies.\cite{odoherty2020,nlb2021,myosuite2022,myosim2026}

Other chains answer different questions. One model predicts processed summaries of muscle electrical activity from a separate brain recording. Another identifies which finger movement a human participant was cued to attempt. A model using a touch-related brain recording reconstructs earlier movement and handle-force quantities from that same experiment. In the opposite mathematical direction, another model predicts expected neural counts from mechanics matched to its source. That opposite-direction model is sometimes called an encoder. These are useful connections, but their outputs are different kinds of things. A finger-class score is not a finger's path through space. An expected count is not a newly measured neural event.

The earlier application snapshot put these chains beside one another without claiming that their recordings are simultaneous. They belong to different subjects, sessions and tasks. They can run concurrently as software, just as several independent videos can play on one screen. Their coexistence does not turn them into one observed nervous system. NeuroTranslator's shared execution protocol is a way to keep their differences explicit while making them usable through the same application.

The proof also includes negative results. A shared model was tested against retained specialists and rejected. Transfer with a fixed model core from one recording from brain tissue to another failed its admission rule against stronger controls. Native macaque muscle controllers failed reserved source cases. These outcomes matter because they constrain the proposed architecture. A system that quietly substituted a successful-looking replay for a failed controller would offer a weaker proof of concept, even if its screen were more impressive.

The proposed connecting diagram can be read as three panels. The first contains an actual recording and a model matched to that recording. The second contains a bounded conversion to a declared computational body. The third contains a separate future biological teaching and retention experiment. An explicit evidence boundary belongs before the third panel. Measured body feedback should return to a recording-based model only where an actual synchronized signal supports that connection. Drawing an arrow cannot supply that missing evidence.

We use three evidence levels throughout. An observation is something present in the acquired recording, within the limits of its measurement method. A model-derived quantity is computed from observations and assumptions, such as joint angles estimated through tracking and muscle-and-skeleton analysis. A simulation is produced by an explicitly chosen numerical body and controller. A prediction can concern any of these levels, but its label must tell the reader which one. The original touch-related source study itself distinguishes recorded neural signals and body-model calculations supported by movement tracking.\cite{chowdhury2020}

That historical prototype established executable connections and several local numerical results. It does not establish the missing causal chain from a teaching signal to changes in a living nervous system to independent retained performance. Those learning-related biological changes are often called plasticity. That chain is the purpose of the proposed research program. Keeping it separate lets the current software be useful without asking it to certify an experiment that never happened.

\begin{center}
\begin{minipage}{\linewidth}
\centering
\makebox[\textwidth][c]{\includegraphics[width=160mm]{figures/architecture.pdf}}
\captionof{figure}{The working chain connects recording-matched predictions to declared body representations. Biological teaching and retained learning belong to a separate future experiment. The boundary between them is an evidence boundary.}
\end{minipage}
\end{center}

\section{Recording Access, Context and Stable Signals}

Three practical questions organize the proposed work: can we observe useful signals, connect them to a measured body task, and keep the connection useful when conditions change? Each question identifies a different way an interface might fail. Each needs a different measurement. They are research questions, not outcomes already achieved by this project.

``Access it'' asks whether useful signals can be observed without overwhelming the system being studied. More channels are not automatically better. A recording can be rich in numbers yet miss the task variable that matters. A proposed interface can promise access while imposing a burden that changes movement, development or tissue condition. The current software begins with already published recordings, so it does not demonstrate a new implant, a lower-burden sensor, or safe lifelong access. Its contribution is to preserve what those recordings actually contain and to refuse to invent the channels they lack.

``Ground it'' asks whether the signals can be connected to a measured task and a particular body. A count from a brain recording channel only becomes useful for this application when we know the source, timing, target definition and evaluation. A predicted joint angle only becomes useful for a body test when we know the coordinate convention and mechanical assumptions. This is where the interactive prototype does most of its work: source checks, explicit quantity labels, time windows, prediction comparisons, anatomy models and numerical tests of movement control.

``Keep it'' asks whether the relationship remains useful as recording conditions, bodies and tasks change. In this prototype, the separate-day motor experiment and the failed external finger pilot illustrate the difficulty. A model that performs well within a session is not automatically stable across days. A model matched to one recording may remain useful after calibration, while transfer with a fixed model core fails to beat fresh training. Longer biological stability, device durability and retained learning are future experiments, not consequences of these local comparisons.

A future experiment should begin with ordinary behavior and clearly defined measurements. If it studies a human-related biological feature, it should say exactly which feature was introduced. If it studies a tissue connection, it should compare the contribution of that connection. An experienced partner and a computer-based trainer are different possible sources of assistance. Disease-specific questions need their own appropriate comparisons. None of those biological experiments was performed by this computational work.

A proposed task sequence should keep ordinary intact behavior visible before studying a defined interruption and replay condition. Simple routes with right-angle turns should be compared separately from routes with varied headings and continuous changes in direction. An animation can test assumptions in its numerical model. It cannot by itself validate a biological injury model, a tissue connection or a disease treatment.

The central learning distinction is equally explicit. Assistance changes performance while help is supplied. Retained learning changes what the trainee can do after trainer commands are removed. Regeneration concerns restoration of tissue or pathways. These can coexist in a future experiment, but none can stand in for another. If an interrupted pathway still requires a compensatory bridge, the bridge must remain available when trainer commands are withdrawn. Otherwise the experiment removes the trainee's means of expressing movement and mistakes that loss for a test of learning.

The roles of the software and a future biological experiment remain separate. The software builds and checks parts of a measurement and execution chain. The biological experiment would ask what an organism can still do after help stops. A reader should always be able to tell whether a result is an actual recorded observation, a trained prediction, a numerical simulation or an untested proposal.

\section{Brain, Spinal Cord and Peripheral Nerves}

It is tempting to imagine a skill as a command that leaves the brain, travels down one cable, and arrives at the muscles. That picture is useful only as a starting question. The recordings in this project observe different parts of movement-related and sensory activity through different instruments. They do not reveal every relevant signal or turn the rest into a known blank.

The movement sources include recordings from the brain's outer tissue and observations of behavior. The human finger source includes movement-related and posterior-parietal brain recordings associated with attempted movement; ``posterior parietal'' identifies the source region toward the back of the relevant brain tissue. The touch-related source observes a region named Area~2 during reaching and externally disturbed movement. A separate human source records surface electrical signals at several levels during stimulation. Each experiment provides a particular window on the system. Its observation method and task determine which conclusions are justified.\cite{odoherty2020,human_fingers2023,chowdhury2020,spinal_dataset2023}

The movement brain area is formally called primary motor cortex. The touch brain area is formally called primary somatosensory cortex, with Area~2 being the specific source region used here. These names help locate recordings. They do not make a movement-area channel a pure motor instruction or a touch-area channel a complete sensory description. In the implemented models, the practical question concerns prediction: given the recorded history from this source, how accurately can the specified target be predicted on the stated evaluation samples?

The proposed body-wide instrument also needs to preserve direction. Signals travelling from the body toward the brain and spinal cord are called afferent signals; signals travelling from those central structures toward the body are called efferent signals. In plain terms, these are body-to-center and center-to-body directions. The spinal cord and peripheral nerve networks can also process information locally. A diagram with an arrow from brain to muscle should not be read as proof that every step is a passive wire or that the current files measure all of that processing.

Proprioception means sensing body position and movement. Interoception means sensing conditions inside the body. They are related, but they do not name the same information. The numerical controller receives simulated position and movement feedback. It does not thereby observe all internal sensations or reconstruct a complete subjective bodily experience. Each proposed body-wide measurement needs an explicit target and reference.

The spinal cord and peripheral nerves are important to the larger proposal, but their motor activity was not recorded alongside the macaque motor sessions used by the decoders. The registry therefore marks those observations unavailable. It does not insert zeros and call them quiet neurons. Zero would be an actual numerical assertion about activity. Unavailable means the experiment provides no observation of that quantity.

The separate human multilevel preview is valuable because its recording types share one experiment. It contains electrical recordings from the scalp, spinal surface, peripheral nerves and muscles. These combine activity from multiple contributors rather than identifying individual spinal and peripheral neurons. In the acquired preview, electrical stimulation causes the responses being studied. That supports an observation view and source inspection. It does not train the current application to reconstruct voluntary walking or voluntary brain-to-muscle control.\cite{spinal_dataset2023,spinal_companion2023}

The difference between anatomy and activity matters too. A structural map can describe possible connections, while a time series describes what an instrument observed over time. Even a very detailed connectome cannot supply the missing event sequence for a particular movement. A future architecture could use anatomical information as a stated prior, meaning a constraint or starting assumption. It would still need evidence for its time-dependent predictions and would still label unobserved channels as unobserved.

A body-wide model is a measurement challenge. More levels could help prediction or control if their information is relevant and aligned. They could also add noise, delay, burden or repeated information. A fair comparison must allow a smaller brain-only or behavior-only model to win. Otherwise, the larger diagram becomes an idea that its own experiment cannot reject.

NeuroTranslator makes the absence manageable. It can keep recorded brain electrical events, surface signals combining multiple contributors, processed muscle electrical signals, mechanical variables and simulations in separate views with explicit quantity labels. That is a useful foundation for adding matched recordings later. It is not a substitute for obtaining those recordings, and it does not establish that everything important about a skill is electrical or already represented in the current files.

\section{Neural Event Counts}

The models using brain recordings often receive spike counts. In ordinary terms, the input asks how many detected events occurred on each declared recording channel during each short time interval. The original event times are grouped into bins, and the bins are supplied in their original channel order. A model sees a history of these counts rather than a single instantaneous value.

``Unit'' is the source's identifier for a recorded neural feature. Depending on the dataset's processing, that feature comes from sorting electrical events measured outside cells or from a specified recorded-feature inventory. Sorting groups detected events according to the source's rules. The identifier belongs to that source. Unit 12 on one day is not proved to be the same cell as unit 12 on another day. Even equal channel counts do not create identity. The reaching and benchmark sources give NeuroTranslator useful activity from groups of neurons, while their definitions matched to each recording limit interpretation across recordings.\cite{odoherty2020,nlb2021,falcon2024}

Counts are integers because an observed number of events in a completed bin cannot be half an event. This sounds obvious, but a numerical interface can lose the distinction. A model can predict an expected count of 2.4 events. That is a statement about a conditional mean, not a claim that 2.4 events were physically observed. If expected values were silently fed into a port that promises recorded integer counts, the next stage would inherit a false description of its input.

The request validator checks the original numeric values before converting them to the model's lower-precision floating-point representation. The order is necessary because a large fractional number can round to an integer during conversion. The portable tests include that case. A request with 16,777,216.25 in a recorded-count field is rejected as fractional even though a lower-precision representation can conceal the fraction. This guard concerns the signal contract, not whether such a count would be physiologically plausible in the source window.

Electrical muscle targets need a different explanation. The brain-to-muscle-signal model predicts 16 source-processed electrical summaries. An envelope summarizes aspects of the recorded electrical signal after the source's filtering and numerical rescaling. In this release, the targets are rescaled values with no remaining physical amplitude unit. A stale millivolt attribute in the source file does not convert them back into electrical amplitude verified in a physical unit. The model's output is therefore neither force in newtons nor a verified muscle excitation command. The source preprocessing, rather than a convenient file label, determines the quantity.\cite{falcon2024}

The human finger classifier answers a still different question: which of ten attempted-finger classes is most consistent with the recorded window? Its class probabilities are classifier scores without established probability calibration. The task cue provides the class label. It does not supply an executed trajectory, joint torque, grip force or object manipulation sequence. A correct answer to ``which finger was attempted?'' is useful evidence for that classifier. It does not prove that the model can make the finger move.\cite{human_fingers2023}

The sensory encoder moves from mechanics to positive predicted average counts. These means can be scored against recorded counts using a probability-based test of how well a proposed average explains the observed counts. They do not constitute a validated generator of complete biological spike trains. Limits on how quickly a neuron can produce consecutive events, coordinated timing across neurons and a complete account of biological count variability would require additional models and tests. NeuroTranslator keeps the mean-count prediction useful while labeling it as inferred activity rather than measurement or stimulation.

Reversing the prediction does not reverse time or recover a unique cause. Several
neural patterns can accompany a similar movement, and a similar recorded pattern
can occur in different task conditions. A mechanics-to-neural model asks what
recorded activity is expected given the supplied mechanics and context. That is
a conditional estimate: its answer depends on what it was given. It cannot
identify the one true intention merely because the forward model predicts the
movement well. Each direction needs its own withheld reference, uncertainty
and simpler comparisons. A predicted pattern also ends at a computer output;
delivering a sensation would require a separately qualified physical writer.
The current spinal scope remains observation only.

Understanding these distinctions makes the screen readable. Similar-looking traces can represent events, rescaled electrical summaries, classifier scores, measured mechanics or predicted expectations. Their labels carry the scientific meaning. The plot's visual smoothness does not establish that two quantities are interchangeable.

\section{The Two-Finger Recording Study}

The current LINK comparison uses six sessions from one rhesus macaque performing
a centre-out finger task. A centre-out task asks movement away from a central
position toward a target. The pinned public source is DANDI 001201, version
0.251023.2336, Monkey N.\cite{linkarchive} One animal across six sessions does
not mean six animals or a population test.

Four sessions are training: January 27 and 30, February 4 and 5, 2020. Training
changes the model's fitted numbers. Two later sessions, February 6 and 11,
are development validation: they help choose a model and describe its behaviour
during development. Two further February sessions and two June 2023 recordings
remain reserved and closed. Those untouched recordings would answer a later
question under an independently accepted plan. They supply no result in this
paper, and the paper does not imply permission to open them.

Each development session contains 375 trials. Ninety-six physical contacts in
hand primary motor cortex provide two feature types. Threshold-crossing counts
summarize detected electrical events crossing a selected threshold. Spiking-band
power summarizes power in the source's selected high-frequency band. That makes
96 features of each type, or 192 features together. It does not identify 192
cells, two separate electrode arrays or a complete measurement of motor cortex.
The contact order and source identity remain part of every example.

The count features undergo the numerical transformation log of one plus the
count. This retains zero and compresses large values. The power features retain
their archived scale. Both are then scaled using means and spread calculated
only from training data. The spread is the population standard deviation: take
each fitting value's distance from its mean, square those distances, average
them over all the fitting values and take the square root. The saved conversion
subtracts the training mean and divides by that training spread. This conversion
is the scaler. It puts features on a declared numerical scale without claiming
that the contacts measure the same biological quantity.

If the scaler used validation values, the checking recording would already
have helped set the model's input scale. That is information leaking from the
check into the procedure being checked. In the accepted preparation, validation
supplies no scaler statistics.

The scaler uses 107,605 unique training frames owned by their source trials.
Unique means a frame is counted once for those statistics even if several
prediction windows contain it. Overlapping windows can otherwise repeat the
same recorded sample many times. Fitting a scaler to all window entries would
weight samples differently; the declared unique-frame rule is part of the
study, not a cosmetic implementation choice.

The two output targets are normalized sensor positions. One concerns the index
finger. The other concerns the middle, ring and small fingers as a combined
group. Their archive names do not give three individual finger trajectories.
``Normalized'' means the source expresses them on a rescaled sensor scale.
A value on that scale is not an angle in radians, a length in meters or force
in newtons. The study has no simultaneous muscle electrical, peripheral nerve
or spinal recordings. It also lacks the individually measured joint positions
and forces needed to turn two numbers into a fully calibrated body state.

A prediction endpoint needs eleven neural samples inside one source trial and
a following neural frame for the extra learning target where used. The movement
label sits at the last input timestamp. A history that crosses into another
trial is not accepted. The eleven timestamps span about 199 to 201 milliseconds,
roughly one fifth of a second. The prepared support contains 91,105 training
endpoints and 45,757 development endpoints.

Those counts tell us how many usable recorded endpoints exist, not how many
independent experiments occurred. Adjacent windows share samples. The trials
belong to the same animal and dates. Statistical claims must reflect that
structure. In each session, 4,125 frames are not prediction endpoints; some
still provide declared history, the extra neural label or training-scaler support.
Excluding an endpoint does not mean every sample around it disappears.

\begin{center}
\begin{minipage}{\linewidth}
\centering\small
\begin{tabular}{lrrrr}
\toprule Date & Role & Trials & Source frames & Usable endpoints\\\midrule
2020-01-27 & Train &375&27,649&23,524\\
2020-01-30 & Train &375&26,944&22,819\\
2020-02-04 & Train &375&26,441&22,316\\
2020-02-05 & Train &375&26,571&22,446\\
2020-02-06 & Development &375&26,957&22,832\\
2020-02-11 & Development &375&27,050&22,925\\\bottomrule
\end{tabular}
\captionof{table}{Six recording sessions from one animal. Each usable endpoint obeys
the same history and trial rules. These are exact support counts, not a number
of independent people or a claim about confidence. The reserved later and
drift recordings are excluded.}
\end{minipage}
\end{center}

An independent source audit reproduced the declared prepared values. Its saved
scaler comparison differed by at most $4.54\times10^{-13}$ under the fixed
criterion. That is a narrow check of preparation arithmetic. It does not say
the model predicts movement well or that a live controller has that accuracy.
The manuscript reuses the report; no preparation or calculation was repeated
for this source edition.

Time has a further limit. All the series in each file share the recorded clock,
but archive bin anchoring and filtering remain uncertain. The displayed span
is not verified physiological input support or measured online delay. The
archive also contains a voltage label that conflicts with the publisher's
threshold-count description. Count interpretation remains count interpretation;
absolute power calibration remains unverified. These missing facts must be
settled before a stronger physical or online claim.

The separate human EEG/ESG/ENG/EMG preview remains another experiment, involving
evoked responses to stimulation and its own clock. It does not drive the LINK
models or the older macaque-to-human-arm simulation. Matching picture colours,
starting videos together or placing plots in the same dashboard cannot establish
a simultaneously measured voluntary brain--spine--nerve--muscle chain.


\section{The earlier recording families and their distinct targets}

The earlier models use five recording families. These are a selected working set rather than a survey of all neuroscience data. The reviewed source catalog contains 37 resource records, including data, anatomy and simulation resources, but a catalog record does not mean that its files were acquired or used to train a model. Some entries remain metadata-only. The distinction prevents a broad inventory from being mistaken for broad experimental coverage.\footnote{Local evidence: \texttt{paper/evidence\_matrix.json}, claim \texttt{catalog\_scope}; source inventory \texttt{\mbox{research}/\allowbreak{}\mbox{dataset}\_\allowbreak{}\mbox{catalog}.\allowbreak{}\mbox{json}}.}

The original movement-velocity specialist uses an Indy recording from February 2 with 130 sorted units. Its input bins last 20 milliseconds, and its history covers 20 bins, or 400 milliseconds. Its targets are the supplied two-component fingertip velocity labels. This is a particular reaching experiment and a particular preprocessing path. The public reaching archive and the NLB-related release identify the source; the project uses its own local split and does not claim an official benchmark score.\cite{odoherty2020,nlb2021}

The model for normalized muscle electrical signals uses a separate macaque source from the Few-shot Algorithms for Consistent Neural Decoding benchmark, called FALCON, a comparison framework for neural predictors. Its movement-area brain recording has 64 spike features and 16 muscle electrical channels. Its bins and context have the same stated durations as the velocity specialist, but its recording identity, channel meanings and target preprocessing are different. Matching the input's shape or bin duration would not authorize swapping the sources. The source task concerns reach-to-grasp; that does not turn the current software into an object-contact manipulation learner.\cite{falcon2024}

The human attempted-finger model uses 298 recorded units across the supplied movement-related and posterior-parietal source configuration. It receives 18 bins of 50 milliseconds each, a 900-millisecond window, and classifies ten cued attempted-finger labels. Its selected method calculates class scores directly from a summarized window, rather than using a changing internal memory through the window. The task labels are meaningful exactly as attempted-intent labels, including the limits of the experimental context. They should not be relabeled as actual finger motion or general dexterity.\cite{human_fingers2023}

The Area~2 source supplies 65 units and paired movement and handle-force quantities. The decoder outputs 17 quantities spanning hand and handle variables and tracking/model-derived joint angles. Seven angles and 39 modeled muscle-tendon lengths also support the native-body analysis. The upstream study explains how whole-arm movement and biomechanical analysis produce those quantities; they are not direct observations of isolated muscle fibers or peripheral motor firing.\cite{chowdhury2020}

The historical separately calibrated velocity model uses a genuinely different Indy date, January 31, from the original MATLAB numerical-data archive. It contains 171 nonempty sorted features. Residual or hash cells, source entries outside the sorted-unit inventory used here, are excluded. The file was selected by a fixed eligibility rule before inspecting its raw signals or model outcomes. Its velocity labels are derived from recorded fingertip positions by taking differences between successive available positions and dividing each difference by the elapsed time between those samples, without reading future positions. The new target recording occurred before the pretrained source recording, so the comparison looks backward across archived dates even though file choice was fixed before outcome inspection.\cite{odoherty2020}

The separate human multilevel preview supplies surface recordings during stimulation. Its main acquired view includes scalp, spinal, nerve and muscle electrical signals; the companion file has its own channel inventory and omits muscle electrical signals. A broader repository description cannot override that file-level fact. Neither preview was used to train a voluntary whole-body movement predictor.\cite{spinal_dataset2023,spinal_companion2023}

Licenses travel with these sources. The used brain-recording resources in the public DANDI repository are listed with the Creative Commons Attribution 4.0 license, whose identifying code in the catalog is CC-BY-4.0. The human surface sources have public-domain dedication entries identified as CC0-1.0. These are exact source-license labels; their source records remain the place to inspect the terms. Body software and other catalog items have separate terms. Open access to an upstream file does not automatically license every new combined artifact, and this local project has not verified a public software deposition or settled its own release license. Keeping source versions and attribution visible is part of reproduction, not an optional decorative appendix.

\section{Physical Units and Time Alignment}

A number without a quantity is incomplete. The value 0.2 could mean meters per second, an angle in radians, a rescaled muscle electrical summary or an expected event count. A radian is an angle unit; it is not a distance or a speed. A movement pipeline can run successfully while confusing these meanings. NeuroTranslator therefore asks for the units and conventions matched to the source that make the number usable.

The Area~2 handle variables provide a concrete example. Six columns combine pushes and pulls with twisting effects, formally called forces and moments. Forces use newtons, while moments use newton-meters. The combined six-column quantity is formally called a wrench. An incomplete single unit label on the group does not give every column the same unit. The acquisition and group description determine the six quantities. Similarly, source joint angles supplied in degrees are explicitly converted to radians before being used by the appropriate model. A request to the joint port must provide the recorded names, order, radian units and source coordinate frame.\cite{chowdhury2020}

A coordinate frame tells the system what the axes and angles refer to. The same seven numeric angles could describe different joint conventions in different models. A human simulator cannot declare itself interchangeable with the macaque source by using seven numbers. A verified frame or calibration adapter would be required. The current request guard rejects a different frame rather than silently reinterpreting the values.

Time is just as specific. An input window ends at the last complete bin. If a 20-bin history uses 20-millisecond bins, its entries must follow that cadence when timestamps are declared. Missing bins or an inconsistent endpoint are rejected. If the caller omits timestamps, the response says that expected time spacing was assumed. That is weaker evidence than checking the actual supplied times, and the response preserves the difference.

For the sensory decoder, the selected model reconstructs mechanics 40 milliseconds earlier than the neural input endpoint. This follows a comparison of specified alignments on the same set of samples and validation selection. It does not show that sensory activity is a future motor command. It also does not measure a biological conduction delay in isolation, because upstream processing and tracking delays are not fully identified. In the encoding direction, the source mechanics history is associated with expected activity at the corresponding later model time.\footnote{Local evidence: \texttt{artifacts/sensory\_metrics.json}; evidence-matrix claims \texttt{sensory\_decoder} and \texttt{sensory\_expected\_activity}.}

The January 31 velocity target illustrates a different offset. Fingertip positions are recorded at 250 samples per second. For each 20-millisecond neural bin, the processing selects the last available position sample, then takes the position difference between successive selected samples and divides by their elapsed time to obtain velocity. The actual position-sample coverage ends four milliseconds before the neural-bin endpoint. The velocity port therefore reports that earlier output time. Four milliseconds here describes sample timing; it is not an estimated neural transmission time.

Fingertip and cursor positions are also kept separate. In the audited new source, their discrepancy reaches 11.9 millimeters. A cursor may be the displayed task variable, but substituting it for the recorded fingertip would change the target. The project verifies the source transform from centimeters to meters and keeps the quantity explicit. A small coordinate or timing substitution can change a scientific interpretation even when the picture still looks plausible.

These rules define computational sample coverage. They say which recorded samples a calculation may use and where its output belongs on that source clock. They do not remove any filtering that already used later samples in the source, certify the delay in a living nervous system, or authenticate a caller's declared measurement. NeuroTranslator can check the packet's consistency; it cannot prove that arbitrary supplied values came from a real electrode. Its response states that limitation.

For an online prediction, the decisive time is when information becomes
available. A sample with an early timestamp could still arrive late, or have
been filtered using a later observation. Calling an input causal means that
the issued prediction uses only information actually available by its declared
deadline. It does not mean the prediction has proved a biological cause. Future
neural samples can serve as training targets for the extra learning task;
they must not become inputs to a prediction issued before they are available.

\section{Earlier Implementation Defects}

One of the most useful findings in this project was a data-processing failure. Earlier muscle-signal and shared-model results looked encouraging, and software tests had passed. An independent check against the raw source then showed that the event streams were not all globally ordered in time. The binning method had relied on that ordering.

The mistake is easy to understand without knowing the code. Suppose we want to count arrivals between two times and use a fast search that assumes an arrival list is sorted. If later arrivals sometimes appear before earlier ones in the list, the search boundaries no longer identify the right interval. The counting routine can return numbers of the expected shape, and the rest of the model can train on them. Successful execution does not make those counts faithful to the recording.

The audit found out-of-time-order streams in 33 of the 64 neural features in the muscle-signal experiment. Applying the ordering-dependent method before sorting had corrupted 866,594 unit-bin cells. Those are counts for particular channels in particular time bins, so the defect altered the actual numbers supplied to the models rather than a cosmetic display. Earlier affected muscle-signal, shared and uncertainty results were withdrawn.\footnote{Local evidence: \texttt{artifacts/source\_audit.json}; evidence-matrix claim \texttt{spike\_\allowbreak{}order\_\allowbreak{}withdrawal\_\allowbreak{}counts}.}

The correction sorts the event times under the documented binning rule and compares derived counts with independent whole-source histograms. Models were rebuilt using unchanged random starting values, data divisions and selection rules. Keeping those choices fixed matters: otherwise a repair could quietly become an opportunity to search for better-looking outcomes. The new results are the corrected results. The earlier positive numbers remain in a withdrawal record and must not be combined with the corrected evidence.

This failure also explains why independence between checks matters. A test that repeats the same faulty search can agree perfectly with its implementation. A check based on the raw event arrays and an independently constructed histogram asks a different question: do the resulting bins match the actual source events? Agreement between distinct routes is stronger evidence than repeated execution of one route.

The project found other numerical defects in native-body handling. Borrowed rotation storage could change after export rather than preserving the intended rigid transform. Named control functions obtained before model initialization could become stale references. Those problems required their own corrected exports and bindings. A finite array, a well-formed JSON file or a green test result can conceal the wrong relationship to a native object.

A later result stands on the check appropriate to its own failure mode. Event order needs source-level verification; coordinate transforms need geometric verification; a controller needs a physical tracking comparison. Biological learning would need a biological retention comparison. That gives each result a reason to be trusted or questioned, instead of treating successful software execution as an answer to all of them.

The withdrawal record shows how the prototype's scientific claims changed when an independent audit found a concrete defect. A future release should preserve that history so readers can distinguish valid and superseded results and see why the change was made. Keeping the corrections inspectable is part of the proof of concept and of a trustworthy demonstration.

\section{Simple-Model Baselines}

The comparison begins with the training-target mean: a model that always returns
the same two numbers calculated from training targets. That simple reference
asks how much is gained by looking at neural input at all. A second model takes
a fitted weighted sum of the last neural sample. A third takes a fitted weighted
sum of the whole eleven-sample history. A weight is an adjustable number: multiply
an input by its weight, add those contributions, then add a constant offset.
Fitting chooses the weights and offset from training examples so that the
predictions approach the targets. A weight's sign sets the direction of its
contribution; its magnitude sets how strongly the input affects the sum. This
does not establish a stronger biological connection. Their formal fitting method
is ridge regression, which discourages
excessively large weights through a numerical penalty. That penalty makes large
weights costly even when they reduce error on the examples used for fitting.

The penalty choices were fixed in advance: 0.01, 0.1, 1, 10 and 100. Both linear
models select 100 on the development dates. The constant term is not penalized.
This is a controlled small search, not a promise that ridge is the best model
for every neural task. It supplies a genuine alternative that the more complex
candidate must beat under the same examples and scoring rule.

The two larger candidates use an established gated recurrent unit. That long
name describes a software memory updated as it reads a sequence. At each sample,
the model combines the new input with a numerical summary of what it has read
so far. Fitted calculations called gates control how much of that summary is
retained and how much is changed. The next sample is read using the updated
summary, so the same current input can produce a different estimate after a
different history. Whether that capacity helps is still an experimental question.
It is not a recreated nervous system and its numerical memory is not a measured
biological memory. The same architecture receives 192 features at each sample,
has one layer of 128 internal values and returns two position estimates.

Two kinds of change must stay separate. The summary changes while a fitted model
reads an input window. The weights specify how it changes and stay fixed during
that prediction. Training adjusts the weights using examples. A changing trace
on a replay therefore does not show that the model is retraining, still less
that an animal is learning.

Both candidates keep the same total 169,090 adjustable numbers. The position-only
version actively uses 144,322 of them in its objective. The extra output for
predicting the following neural frame remains present so the total parameter
count is disclosed fairly. The second candidate trains that extra output too.
Its loss is position mean squared error plus 0.1 times next-neural mean squared
error. A loss is a numerical measure the fitting procedure tries to reduce.
The extra task is called auxiliary because it is intended to support the main
position task; whether it actually helps is a question for the comparison.
Both error terms square each individual prediction difference before averaging.
The added term gives fitting another reason to retain information in the shared
numerical summary. It does not provide a measured new pathway, and success at
predicting the next neural frame cannot replace success at predicting position.

An epoch is one pass through all eligible training windows. Each run saves 40
epochs. Batches contain 256 windows, and the final smaller batch is retained.
A batch is the subset used for one fitting update; an epoch therefore contains
many updates. The model predicts the batch targets, the objective measures its
errors, and a gradient describes how changes in the adjustable numbers would
change that objective. Backpropagation calculates those changes through the
successive model operations. It is a software calculation, distinct from sensory
feedback travelling through a living body.

The software uses the established Adam update rule with learning rate 0.001,
internal averaging values 0.9 and 0.999 and numerical stabilizer $10^{-8}$.
Adam keeps running summaries of the gradients and their squared sizes, then
uses them to scale the change to each adjustable number. The learning rate
sets the scale of those changes; the stabilizer helps the division remain
numerically defined. These settings describe the fitting procedure, not a
biological learning rate or a promise that each update improves validation.
Weight decay would add a penalty that pulls model weights toward zero. Dropout would randomly leave out internal model units while training. Gradient clipping would limit the size of the calculated update-direction vector. Target clipping would replace target values outside chosen bounds with those bounds. None of these four controls is used here.
Together, these choices specify the saved computation. Altering them would
change the comparison.

The reported configuration uses four mathematical worker threads and one
framework scheduling thread. A thread is a stream of processor work; these
settings describe the saved calculation.

Three fixed seeds, 20261007, 20261008 and 20261009, identify the random starting
conditions. The two candidates receive matched starting values and training
orders. Keeping a seed means keeping its outcome even if another seed performs
better. Each seed chooses its lowest development-error epoch; an exact tie
chooses the first. The final prediction is the equally weighted average of
all three selected models. That average is called an ensemble. Selecting only
the best-looking seed would be a different experiment.

The measured condition uses the declared source relationship. An association
control challenges the pairing that should carry useful neural information.
During control fitting, a neural history and its next-neural target move together
to another training example within the same recording day, while the position
answer stays in place. The validation observations stay unchanged. This asks
how much the result depends on learning the intended training association.
It does not remove every possible association: task structure and the held
validation context still exist. The frozen methods below give the exact rule.
A temporal-order control changes the order of the supplied history while keeping
the current bin at the end. It challenges the particular declared ordering;
it does not prove that all time information has disappeared. A history might
be useful because it averages noise without making use of sequence order;
those are different claims. Both controls have their own matched
runs, and neither enters the competition for the best model. Two candidates,
three seeds and three conditions produce 18 runs, each with 40 epochs, for 720
saved epochs altogether.

The primary score is mean squared error. Subtract the measured target from each
prediction, square the difference and average those squares. The result is in
squared source-normalized position units. Larger mistakes contribute more heavily. The study first forms
one score for each finger group on each validation day, then averages those
four scores equally. A day with more usable frames does not take more weight
in the final four-cell average merely because it has more frames.

A date/group cell is one of those four score entries. Requiring no harmed cell
makes an important choice explicit: improvement should not come by sacrificing
one of the declared outputs or days. An average can improve while a real part
of the support gets worse. The study's acceptance rule does not allow that
tradeoff. Changing it after seeing the result would make a different claim.

The uncertainty comparison resamples whole trials within each date. A bootstrap
is a repeated calculation in which samples are drawn again from the observed
support to examine variation in a statistic. Here each draw keeps every supported
frame of a trial together. The same trial draws are used for both compared models
and both output groups, so their difference is paired. There are 1,000 saved
repetitions, and the interval uses the appropriate percentiles of their scores. Percentiles are cut points below which specified fractions of the scores lie.

This does not make the selected validation data into untouched confirmation.
The models were chosen using those dates. The 95-percent interval describes
variation conditional on the selected models and those two dates from one
animal. It is not a population guarantee, a test on another species or a
prediction about every future day. Treating each overlapping frame as an
independent animal would make the uncertainty claim much too strong.

The complete development requirement has four parts. The candidate must have
lower average error than the selected ridge. It must have no higher error in
any date/group. The upper end of the paired candidate-minus-ridge interval must
be below zero. Finally, every cell must have a defined positive ordinary
$R^2$ that exceeds the corresponding association control. To obtain that score,
compare the model's summed squared errors with the summed squared errors of
always predicting this evaluation cell's observed mean, then subtract their
ratio from one. That observed mean is a scoring reference, not an answer supplied
to the deployed predictor. If the mean-reference error is zero, the ordinary score is
undefined and cannot satisfy the requirement. A positive value means beating
that mean reference; beating the association control is an additional check.
Neither comparison alone establishes a biological causal mechanism. All four
must pass.
Separate history, order-use and extra-task benefit claims keep their own
comparisons. A falling training loss does not replace any of these conditions.


\section{How the earlier small models read recent history}

The earlier NeuroTranslator implementation used small established models that read recent signal history. One such model reads a sequence while maintaining an internal numerical summary that changes with each entry. Learned controls regulate how much of the previous summary is retained and how much is updated. In this application, the sequence consists of completed neural or mechanical bins. The model's final summary and simple recent averages feed an output calculation for the specified target. The formal name of this model type is a gated recurrent unit: ``gated'' refers to those controls, and ``recurrent'' means that the updated summary returns for the next entry. It is a mathematical tool for predicting from history, not a model of every biological process involved in learning.

In those earlier models, before the history-reading model sees spike counts, the pipeline applies a square-root transformation and numerical scaling calculated from training samples. The transformation and numerical scaling are parts of the saved model contract. They must be reproduced during execution, which means using the trained model to predict from a new valid input. Recomputing that scaling from the evaluation recording would change that contract and risk using information outside the declared training samples.

The training sequence includes initial training to predict neural counts in the next bin for the relevant history-reading models. Initial training asks the model to learn patterns in the available training neural history before adjusting it on paired outputs. That adjustment then fits the model for a particular target such as velocity or processed muscle electrical signals. Neither step teaches a biological nervous system. Both change numbers in a software model. When this document says the model learned, it means that numerical optimization changed its fitted numbers.

The output calculation matters as much as the history-reading core. One simple competitor combines recent time-scale summaries through fitted weights, with a penalty discouraging weights from becoming too large; its formal name is ridge regression. Another calculates fitted class scores and turns them into positive values that sum to one; its formal name is logistic classification, and it supplies the selected human finger-intent method. These simpler methods are genuine competitors. The project does not require a complicated model to win merely because the subject matter is neural recording.

The original velocity history model has 232,388 fitted numbers, and the muscle-signal history model has 195,216. The sensory decoder and full-mechanics encoder have 38,802 and 21,266, respectively. The joint-only encoder has 18,696. The new-day velocity model trained without borrowing another model core has 256,045. These counts describe particular saved models. They do not establish that this size is optimal, that larger models cannot help, or that every task can be solved with the same architecture.\footnote{Local evidence: \texttt{paper/evidence\_matrix.json}, claim \texttt{compact\_architecture\_counts}, with paths to the corresponding metric artifacts.}

A saved checkpoint contains the fitted state and associated processing information. Loading it later should reproduce the selected model's behavior. The package also exports plain numerical arrays and architecture descriptions for execution through NumPy, a numerical-array library. That makes the learned state inspectable without relying only on a training framework. It does not make weights automatically executable inside every chat service, and it does not make their predictions independent of the source and preprocessing that gave them meaning.

The portable runtime mirrors the accepted prediction calculations and uses the same request validation as the existing protocol. Checks of numerical agreement compare its values with both saved examples and native execution. The value of those checks is practical: another compatible program can execute the same mathematical mapping. Agreement between runtimes is evidence about implementation fidelity. It does not create a new evaluation dataset or strengthen the model's biological claims.

\section{Why the earlier implementation began with specialists}

There is already substantial research on predicting from groups of neural signals and learning useful numerical descriptions. The Neural Latents Benchmark supplies common tasks for evaluating models using groups of recorded neurons, and FALCON emphasizes robustness when recording conditions change and recalibration is limited. Those efforts provide useful evaluation ideas and existing baselines. The present local results are not official scores on their entire benchmark protocols.\cite{nlb2021,falcon2024}

Earlier published neural-data models study training across multiple recording contexts and handling populations with different recorded channels. One learns from activity recorded within brain tissue across contexts; another represents neural events and uses a selective information-combining stage to accept differing populations. These are relevant precedents for a broader shared representation. They do not imply that every mixture of recordings will improve every task, and they do not remove the need to identify whose channels are being used.\cite{ndt2_2023,poyo2023}

Other model families observe different processes. BrainLM is a published model for activity represented in functional magnetic-resonance brain scans. Neuro-GPT is a published model that combines scalp electrical-signal processing with history modeling and evaluates classification during imagined movement. Those recording types have different instruments, sampling characteristics and interpretive limits from sorted electrical events recorded within brain tissue. Their use of learned representations is a reason to study methods suited to each observation type, rather than to treat all recordings as interchangeable text for one language model.\cite{brainlm2023,neurogpt2024}

The earlier implementation therefore began with explicit specialists. The current LINK study tests a further narrowly defined comparison and keeps its failed improvement requirements visible. A model matched to one recording has a narrow source and target: this recording's recent events to this recording's defined velocity, or this source's mechanics to its expected sensory counts. Its limited scope makes failures easier to locate. If a shared model later improves meaningful comparisons on samples kept outside fitting, it can be admitted with its own evidence. The first shared two-task candidate did not meet that requirement and was rejected against the retained specialists.

Local hardware also makes a bounded first implementation useful. The detected system has an AMD Ryzen AI Max+395 processor and about 125 gibibytes of memory, where the unit uses binary-sized rather than decimal-sized gigabytes. Those historical training runs used four processor threads, meaning four streams of processor work under the chosen setting. Those recorded PyTorch numerical-learning runs used the central processor; graphics acceleration was not verified for their workloads. The current LINK protocol also uses its specified CPU route. A separately reported tiny prefix took 4.565 seconds end to end on GPU versus 2.427 seconds on CPU and remained unpromoted. GPU means graphics processor. The prefix does not qualify full training, energy savings or subscription-cost savings. A large memory capacity is helpful for holding arrays and running local tools, but it is not evidence that this machine trained a large general model or accelerated work through unverified graphics hardware.\footnote{Local environment receipt: \texttt{paper/evidence\_matrix.json}, field \texttt{observed\_local\_environment}.}

This hardware choice makes the prototype reproducible within a clearly described environment. It also limits what can be inferred about cost and performance elsewhere. A different processor, numerical library or thread configuration may produce different timings. The package's physics and native-model components have separate runtime requirements. A useful handoff tells the recipient which parts are lightweight model execution, which parts are cached replay, and which parts need a larger or different environment.

Model size remains a choice to test. The broader design may need a larger shared core if it improves adaptation, robustness or estimates for recording types that are missing. It may instead hide a simple source mismatch or perform worse than a model matched to one recording. Complexity has to earn its place in that comparison, which must allow either outcome. NeuroTranslation is not defined by keeping every model small or by making every model larger.

\section{How the earlier studies separated fitting and evaluation}

When a model is fitted to examples, its predictions on those same examples do not by themselves show that it will work elsewhere. The project separates data used to fit its numbers, data used to select a model, and data used to describe its performance after selection. These sets are conventionally called training, validation and test samples. Their meanings depend on the actual data-division rules, not just the names printed in a file.

The history-reading prediction procedures divide samples in time order, with roughly 70 percent for training and 15 percent each for validation and local testing, using trial rules matched to each source. A one-second gap at the boundaries prevents an input history or its next prediction target from casually crossing the edge between sets. Whole trials are retained where the source and analysis require it. Area~2 splits retain whole successful trials, and the sensory timing comparison uses a common set of times before validation chooses among the specified offsets.

The boundary gap matters because adjacent windows overlap and because a model history can include samples from before its endpoint. A test endpoint just after a training boundary could otherwise contain much of the same input history. The gap and sample-coverage rules address that particular leakage of fitting information into evaluation. They do not remove every possible dependence in the underlying experiment or every source-side filtering choice.

Validation has a legitimate but limited job. It can choose among predefined candidate decoders, stopping points or alignments. Once choices are made, test results describe the chosen procedure on the declared test samples. Repeatedly looking at test outcomes and redesigning the procedure turns that support into development evidence. A later independent evaluation needs fresh support appropriate to the new claim.

Several results in this project have that development status. The human finger evaluation was inspected during an earlier pilot. Earlier test outcomes had also been inspected before the later shared-model comparison. Both transfer conditions and all four methods use the same local motor test windows. These facts do not make the computations useless. They change how strongly the results can support claims about new recordings and reliably measured performance differences.

The models use numerical rescaling fitted only on training samples in the audited project calculations. However, earlier source processing can still include filtering, estimating between observations or rescaling done before the files reached this project. Those source operations can use samples beyond a live prediction's endpoint. A completed-window rule in the project cannot undo that earlier use of later samples. For that reason, the paper separates the samples used by its own calculations from timing in a living nervous system. A model can use only its declared recent bins while its target still carries processing from a broader source interval.

Shuffled-input controls ask whether the association being used depends on the relevant alignment rather than a simpler artifact. Linear comparisons ask whether a complicated temporal model contributes beyond a simpler alternative. They answer different questions, so neither substitutes for the other. A future biological learning experiment would need additional controls on practice, assistance, task exposure and retention that these software splits do not supply.

Samples kept out of fitting within one session still evaluate that session. Once a sample has been inspected, changing its filename cannot make it a fresh test. A fixed file-selection rule documents how an archived recording was chosen; it does not turn that recording into a test of future biological deployment. Report that history with the result so the reader can judge its scope without guessing how often it was revisited.

\section{Results: Average Gains Failed the Required Tests}

The study completed every planned run and saved all 720 epochs. The result is
not missing because fitting stopped early. It is negative because the completed
candidates did not meet their declared requirements. An independent audit supports
the reported negative calculations after fitting; that is the meaning of
``postfit''. It checks the report without establishing that the candidate models
are scientifically effective.

\begin{center}
\begin{minipage}{\linewidth}
\centering\small
\begin{tabular}{p{.47\textwidth}rl}
\toprule Model & Average squared error & Result\\\midrule
Constant training-target mean &0.0317577871&Reference\\
Linear, last sample &0.0314982382&Selected reference\\
Linear, eleven samples &0.0383765124&History did not help\\
Position history model, three-model average &0.0312304247&FAIL\\
Position plus next-neural task, three-model average &0.0299609175&FAIL\\\bottomrule
\end{tabular}
\captionof{table}{Saved development summaries from one macaque and two dates. Lower
average error is one requirement among several. Both history-reading models
fail the full requirement. No values were recomputed for this manuscript.}
\end{minipage}
\end{center}

The longer linear history performs worse than the last-sample linear model.
That is a result for these fitted linear references; it is not proof that
neural history contains no useful information. The two recurrent candidates
produce lower average error than the selected last-sample ridge. That tells
us one favourable fact about their aggregate score, but does not close the
comparison.

For the position-only candidate, the candidate-minus-reference average error
is $-0.0002678134$. Its conditional 95-percent interval runs from
$-0.0014148839$ to $0.0009468841$. Negative means lower candidate error; positive
means higher candidate error. This interval crosses zero, so it does not satisfy
the requirement that its upper end be negative.

For the candidate with the additional neural-prediction task, the average
difference is $-0.0015373207$ and the interval runs from $-0.0027586737$ to
$-0.0003470099$. That interval does stay below zero. It is a genuine favourable
aggregate description under the declared development scope. It still cannot
override the other required checks.

On February 11, the position-only model is worse than ridge for both measured
groups: its squared-error excess is 0.0052471986 for index and 0.0061405635 for
the combined group. The extra-task model is worse than ridge for index by
0.0043574767. These are differences on the source-normalized sensor scale.
They are not radians of joint error or newtons of force, and a conversion to
those quantities is not silently supplied.

Ordinary $R^2$ asks how squared prediction error compares with always predicting
the evaluation cell's observed mean. One is exact agreement under that score.
Zero matches that mean reference. A negative value is worse than the reference.
The position-only model's second-day values are about $-0.099742$ and
$-0.353799$. The extra-task model's values are about $-0.071162$ and $-0.116390$.
Those negative values fail the requirement for positive information in each
cell. A model can do better than a weak shuffled comparison and still fail
that positive requirement.

The training-mean skill score is separately labelled because it uses a different
constant reference. It is not another spelling of ordinary $R^2$. Likewise,
an error score is not a fraction of a skill learned. Keeping these definitions
plain prevents a favourable-looking number from changing meaning as it moves
from a report to a figure caption.

All five further benefit claims fail their complete requirements: history for
the position-only model; sequence-order use for that model; history for the
extra-task model; sequence-order use for it; and benefit of the extra task over
the position-only model. Some subchecks pass. The complete claims do not. Those
outcomes remain recorded with the failures rather than being replaced by a
single success label.

The negative result can be checked without certifying the entire learning
calculation. Both learning routes share Torch's automatic differentiation and
Adam: the former calculates how the objective changes with the fitted numbers;
the latter uses those changes to update the state. If both routes inherit a
mistake there, they could agree for the same wrong reason. Their agreement
cannot independently verify that shared machinery.

A separate numerical route has to check more than a final score. The forward
calculation makes a prediction. Derivatives describe how that calculation changes
with its adjustable numbers. The update rule uses those derivatives and saved
running summaries to choose the next numbers. A checkpoint saves fitted state
at an epoch; the random-generator state also determines subsequent random
choices. Restoring only part of that state could reproduce a prediction while
failing to reproduce the next update. Complete independent numerical values,
random-state continuity, reload behaviour, final-checkpoint predictions,
derivatives, update state and the whole learning trajectory remain unverified.
Checking them would establish the specified computation, not make an ineffective
model effective or turn its software memory into biological memory.

Those missing checks prevent a verified continuation. The final-checkpoint and
update requirements remain unmet, no new fit is currently admitted, all reserved
dates stay closed, and neither a decoder nor a GPU route has been promoted.
The report supports these negative computational findings, not biological
teaching, a universal neural language or completion of the programme.

\begingroup
\looseness=1
The later day's index output poses a specific next question: why did it get
worse? Keeping that failure visible makes the comparison useful even when its
average improves. A further study would have to explain the failure under an
independently admitted plan, without feeding the checking answers back into
the predictor.
\par
\endgroup

\section{Failure Cases}

The saved diagnosis asks what shape the later-day index failure has. The model
with the extra neural task predicts, on average, 0.058645 below the measured
target in the source-normalized sensor units. About 78 percent of its excess
mean squared error relative to ridge is the difference in squared mean error;
about 22 percent is the difference in the spread of errors around their mean.
Those are descriptive components of an already observed failure.

A mean offset is not a cause. Different trial composition, a changing recorded
relationship and a calibration change could all affect such a summary. The
analysis leaves those explanations competing. It does not show which one
occurred, fit a correction or give permission to add the validation answers
to a prediction. Moving the predictions until their validation mean matches
the target would use the checking data to adjust the model.

A separate analysis looks only at the four training dates. It defines context
as the average of the ten preceding neural bins in each existing eleven-bin
window. That context is not an observed biological state; it is a specified
summary of recorded features. The existing recurrent model already sees those
bins. Showing that the summary varies cannot prove that a new branch will
extract useful information from it.

The training support contains 1,500 trials and 91,105 endpoints. The saved
analysis divides summed feature variance into variation within a trial,
variation between trials on the same day and variation between days. Variance
summarizes squared deviations from a stated mean. For threshold-count features,
the three shares are 71.56, 18.85 and 9.59 percent. For archived power features,
they are 51.28, 13.55 and 35.17 percent. These shares give greater weight to
features with greater variance. They are not percentages of the brain doing
different jobs.

An alternative calculation gives equal weight to dates, trials and frames
within trials and produces similar fractions. That is another view of the
same training support, not another cohort. These summaries have no population
confidence interval and do not establish neural drift, a particular cell state
or a benefit for decoding. The typical within-trial feature-to-target
associations are weak. A proposed interpretation should remain proportional
to those observations.

The analysis also checks whether prior sensor labels are actually supported.
Exact joins of source index, time, trial and ownership retain 89,605 endpoints
with one supported prior sensor sample and 88,105 with two. When a beginning
of a trial lacks a label, it stays missing. No value is interpolated and no
sample is borrowed from a neighbouring trial. Sixteen combinations of the two
prior-position quartiles have support on every training date. A quartile is
one of four divisions of an ordered distribution; combining two four-way
divisions gives sixteen cells.

That support is useful for descriptive checks. It is not permission to feed
body-position labels to a predictor meant to use neural input only. Nor are
the correlated frame counts independent observations. The question for a new
neural branch remains whether its available neural information improves the
specified prediction under all the original controls.

Even descriptive arithmetic can fail. The first analysis version missed the
fixed tolerance for two skewness values, which measure distribution asymmetry.
The absolute and relative tolerances stayed at $10^{-12}$. A corrected version
changed only how the per-date skewness was accumulated. Its 768 checked values
passed the named independent comparisons while all 119,848 other numerical
leaves remained unchanged. The rejected report was retained. This accepts a
limited calculation, not a model or a new fitting procedure.

The proposed next branch would let past neural context modify the existing
output. It would add 780 parameters and compress that context into four learned
summaries that modulate the prediction. This is the proposed rank-four branch:
rank restricts the numerical representation; it does not name four measured
brain regions or an anatomical scale. The comparison would have to test constant
context, broken association and altered order. Choosing a model strictly within
training also requires a smaller fitting-and-checking division inside that
training support. Whole trials must stay together in those inner divisions,
and the divisions must be independently verified. Each inner fitting set must
supply its own scale, counting its unique samples once. Scaling the already
rounded, globally training-normalized packet again cannot establish that
isolation: the inner checking samples may already have contributed to the
earlier scale. Another scaling step cannot prove that their influence was absent.

The proposal remains unexecuted and unadmitted for current fitting. It adds no
body labels, date identities or validation offset corrections to the neural-only
predictor. A new fitting stage would have to pass its source, interface, random
state (the exact state of its random generators), derivative, update, selection and uncertainty requirements before any
new result could be claimed. The four reserved dates remain closed. Proposed
body feedback, reward-driven refinement and broader extensions also remain
proposed; the description is not an execution report.


\section{Earlier Model Results}

The velocity model reports a prediction score of about 0.729 on the local samples kept out of fitting, compared with about 0.559 for the simple weighted-sum competitor. The score's formal name is $R^2$. It compares squared prediction error with the variation left by a prediction that always uses that output's mean reference value in the evaluation set, under the stated scoring definition. A value closer to one indicates a closer fit on those samples. It is not a percentage of biological understanding, and 0.729 does not mean that 72.9 percent of a skill has been decoded.\footnote{Local evidence: \texttt{artifacts/metrics.json}; evidence-matrix claim \texttt{mc\_rtt\_velocity}.}

The corrected model for normalized muscle electrical signals reports about 0.704 versus 0.595 for the simple weighted-sum method. This average value hides considerable variation across channels. The weakest channel is about 0.088, while other channels are much stronger. A reader considering a particular muscle-related output must inspect its own performance. The average cannot certify every channel equally. These are local comparisons on the processed targets, not an official FALCON result or an identified force-control model.\footnote{Local evidence: \texttt{artifacts/emg\_metrics.json}; evidence-matrix claims \texttt{falcon\_emg} and \texttt{emg\_channel\_heterogeneity}.}

The sensory decoder reports a prediction score averaged over physical output groups of about 0.710. The scoring averages five groups of physical outputs with equal group weight, rather than allowing a group with many columns to dominate simply because it has more columns. The distinction matters when comparing this value with a mean taken over individual outputs. Its meaning is reconstruction of the stated earlier source mechanics from this source's touch-area history.\footnote{Local evidence: \texttt{artifacts/sensory\_metrics.json}; evidence-matrix claim \texttt{sensory\_decoder}.}

The full-mechanics sensory encoder gains about 0.177 bits per recorded spike over a constant count calculated from training samples. The joint-only encoder gains about 0.119. In plain terms, their predicted count means explain the recorded counts better under the chosen count-probability test than a constant training-derived mean does. The bits-per-spike unit expresses the difference between count-probability scores on a base-two information scale, divided by observed spike count. It is not the number of bits required to upload a movement or evidence that generated spikes would behave biologically.\footnote{Local evidence: \texttt{artifacts/sensory\_metrics.json} and \texttt{artifacts/\allowbreak{}sensory\_\allowbreak{}joint\_\allowbreak{}metrics.json}; evidence-matrix claims \texttt{sensory\_expected\_activity} and \texttt{joint\_only\_activity}.}

The selected human finger classifier correctly classifies 45 of 50 evaluation trials, or 90 percent, compared with 68 percent for its recent-history model alternative. The evaluation is within-session development evidence because the run had been inspected earlier. A separate uncalibrated external-day pilot achieves 10 percent, matching the ten-class chance level. The favorable within-session result and unfavorable external result belong together: classification matched to one recording worked in one setting, while the broader uncalibrated claim did not.\footnote{Local evidence: \texttt{artifacts/finger\_specialist\_metrics.json} and \texttt{\mbox{artifacts}/\allowbreak{}\mbox{finger}\_\allowbreak{}\mbox{metrics}.\allowbreak{}\mbox{json}}; evidence-matrix claims \texttt{human\_finger\_intent} and \texttt{\mbox{human}\_\allowbreak{}\mbox{finger}\_\allowbreak{}\mbox{external}\_\allowbreak{}\mbox{failure}}.}

These scores answer several useful questions without answering the final learning question. They show that specified quantities can be predicted from particular recordings, sometimes better than simpler alternatives. They expose differences in performance between channels and cross-day failure. They do not show how a new person would perform, whether a different task would work, or whether the predictions can teach a living trainee. Keeping each score attached to its target, baseline and support is how the proof stays proportionate to the evidence.

\section{Fixed-Model Transfer Tests}

Transfer learning asks whether a model's earlier training helps with a different setting. For a neural interface, the appealing hope is that a useful representation from one recording can reduce the work needed for another. The separate Indy day provides a real opportunity to test part of that hope. It changes the recording and target processing, and the experiment does not assume that individual cells can be matched across days.

The original January 31 MATLAB file is 208,792,806 bytes. The publisher's file fingerprint and a locally computed fingerprint identify the acquired bytes. The source's residual/hash entries are excluded from the sorted-feature inventory, and the file does not supply a per-unit record of when each unit was observable. These details matter because a transfer result depends on the actual target source, not merely a date label or an array of convenient width.\footnote{Local source and experiment evidence: \texttt{artifacts/motor\_transfer\_metrics.json}; evidence-matrix claims \texttt{motor\_transfer\_controls} and \texttt{motor\_transfer\_budget}.}

\begin{center}
\begin{minipage}{\linewidth}
\centering
\makebox[\textwidth][c]{\includegraphics[width=160mm]{figures/recorded_replay.pdf}}
\captionof{figure}{Recorded events from the separate Indy day alongside fingertip-velocity predictions from the freshly trained recording-matched model. Reference velocity is calculated from recorded fingertip positions. The channels identify this source's sorted features, without matching cells across days.}
\end{minipage}
\end{center}

Four methods are compared. The transferred model learns a translation layer into a previously trained model core and output calculation whose existing fitted numbers remain fixed. A control uses a translation layer of the same size and a fixed core started with random numbers. A fresh model is trained for the new source without borrowing the previously trained core. The simple weighted-sum method supplies the fourth comparison. The random-core method is particularly useful because it asks whether the transferred core's prior learning adds anything beyond using a fixed structure whose output is more than a weighted sum of its inputs.

Two label conditions were fixed before inspecting outcomes: the first five percent of eligible new training labels, and the full eligible training set. The small condition supplies 1,425 adaptation labels, corresponding to 28.5 seconds at the stated bin duration. It also uses 6,000 validation labels, or 120 seconds, and has 570.14 seconds of neural samples available for training. The previously trained source model previously used 22,598 training and 4,723 validation label windows. These additional sources of information must stay visible when discussing efficiency.

``Five percent'' therefore refers to one adaptation-label budget. It does not mean five percent of all supervision used by the entire procedure. ``28.5 seconds'' does not mean that a person learned the skill in that time, that the method ran online for that duration, or that no other data was required. It is the duration represented by a particular set of paired adaptation windows. A fair future learning-speed claim must account for validation, prior training, training samples of neural signals without paired movement labels and actual deployment time.

In the small condition, transfer using the fixed core has validation prediction score 0.393 and local test prediction score 0.425. The control with a randomly initialized fixed core has 0.312 and 0.331. The freshly trained model has 0.533 and 0.520. The simple weighted-sum method has 0.276 and 0.220. The previously trained core thus improves over the random-core comparison, but fresh training performs better. The result supports the narrow statement that prior core training carries useful information relative to this random control. It does not support admission over the strongest competing method.

In the full condition, the transferred method has validation prediction score 0.683 and local test prediction score 0.704. The control with a randomly initialized fixed core has 0.631 and 0.646. The freshly trained model has 0.724 and 0.747. The simple weighted-sum method has 0.550 and 0.548. Again, fresh training wins. These eight method-condition rows use one fixed configuration and the same 6,050 local test windows; they are single reported values without uncertainty intervals from repeated runs with different random starting values.

The admission rule required transferred validation performance to exceed every control by at least 0.01. Neither transferred condition qualifies. The freshly trained model using the full label set becomes the usable new model matched to this recording, while transfer with a fixed core remains rejected. The 0.747 local result cannot be read as a direct improvement over the original source's 0.729 because the recordings, labels and processing differ. What worked was local calibration on a distinct recording; transfer without new calibration and verified cross-day cell correspondence were not established.

\begin{center}
\begin{minipage}{\linewidth}
\centering
\begin{tabular}{llrr}
\toprule New examples & Method & Selection score & Local test score \\
\midrule
First 5\% & Use the old fixed model & 0.393 & 0.425 \\
First 5\% & Use a random fixed model & 0.312 & 0.331 \\
First 5\% & Train a new model & 0.533 & 0.520 \\
First 5\% & Simple fitted formula & 0.276 & 0.220 \\
Full & Use the old fixed model & 0.683 & 0.704 \\
Full & Use a random fixed model & 0.631 & 0.646 \\
Full & Train a new model & 0.724 & 0.747 \\
Full & Simple fitted formula & 0.550 & 0.548 \\
\bottomrule\end{tabular}

\captionof{table}{All eight fixed method-and-label conditions. Scores compare prediction error with a mean-reference prediction; closer to one means a closer fit on these samples. Adaptation labels exclude additional validation and prior-source supervision. All rows reuse the same local test windows.}
\end{minipage}
\end{center}

\begin{center}
\begin{minipage}{\linewidth}
\centering
\makebox[\textwidth][c]{\includegraphics[width=160mm]{figures/transfer_comparison.pdf}}
\captionof{figure}{The previously trained fixed core helps relative to the random-core control, but fresh training performs better in both label conditions. Neither transferred model passes admission. The bars are single-run results, without repeated-run uncertainty intervals.}
\end{minipage}
\end{center}

The separately trained shared two-task model was also rejected on validation. Together, these outcomes make the broader-design proposal more precise. A broader shared representation remains an idea to test. It is not already NeuroTranslator's successful universal core, and it has to beat appropriate specialists before replacing them.

\section{Prediction Uncertainty}

A prediction can be accompanied by a band that looks reassuringly scientific. The question is what the band was constructed to mean. NeuroTranslator uses limits calculated from validation prediction errors to supply constant-width bands for each output. An error here is the difference between a prediction and its reference. The chosen limit summarizes the size of those observed validation errors for that output. It is a practical diagnostic, not a complete theory of uncertainty.

The nominal level is 90 percent per output. That statement concerns individual channels, not all channels at once. In the reported velocity case, mean test coverage per output is about 90.6 percent while simultaneous two-output coverage is about 84.8 percent. For the muscle electrical signals, the mean is about 90.1 percent but simultaneous coverage across all 16 outputs is about 43.5 percent. The sensory case likewise reports about 90.9 percent mean individual coverage and about 49.5 percent simultaneous coverage. These descriptive values show why a per-channel label cannot be promoted to whole-body confidence.\footnote{Local evidence: \texttt{artifacts/uncertainty\_metrics.json} and\linebreak[4]
\texttt{artifacts/sensory\_metrics.json}; evidence-matrix claim \texttt{empirical\_bands}.}

The windows overlap, and errors through time are correlated. The bands therefore do not receive a mathematical guarantee of future coverage merely because their limits were chosen from earlier errors. They do not separately identify sensor noise, imperfect labels, model uncertainty, source drift or absent observations. A broad band around a measured target is also not a way to create a missing spinal label.

Classifier scores have a related limitation. A finger output of 0.8 is a normalized classifier score. Without calibration evidence, it should not be read as an established 80 percent frequency of correctness for all cases assigned that score. A model can rank classes usefully while expressing unreliable numerical confidence. The present interface calls the scores uncalibrated rather than encouraging a stronger interpretation.

The joint encoder uses an explicit check against training ranges. It compares every angle in the input history with the training minimum and maximum for that coordinate, including a declared small margin for rounding and numerical error. An input outside those bounds is rejected unless the caller explicitly requests extrapolation. This is a useful guard against one obvious source of misuse.

Range membership has a much narrower meaning than complete familiarity. Seven individually familiar angles can occur together in a configuration the model never saw. Familiar angles can also be traversed at unfamiliar velocities or accompanied by different loads and contacts. The training-range report therefore states that coordinate-range membership does not prove joint-configuration coverage and does not validate neural generation. Allowing extrapolation leaves those limits explicit; it does not clip the input into support and pretend the request became observed training-like data.

A future system should distinguish several questions: what was measured with uncertainty, what a fitted model is unsure about, where the source has changed, which variable is missing, and whether the task permits the proposed action. Those questions may need different measurements and methods. The current prototype makes a useful start by keeping signal type, source identity, observed bands and support warnings separate. It does not claim a single confidence score that resolves all of them.

\section{Earlier Software Tests}

The historical NeuroTranslator snapshot exposed seven named operations through one request protocol. Each operation, called a port, declares its input and output. The protocol gives callers one way to request them while preserving those different meanings. Selecting the right port is part of the scientific contract: one operation must not silently take the place of another.

The ports are original motor velocity, normalized muscle electrical signals, human finger intent, sensory mechanics reconstruction, full-mechanics expected sensory activity, joint-only expected sensory activity, and calibrated-day motor velocity. Their names describe the actual operations. They do not imply seven simultaneously observed pathways or a complete map of the nervous system. Two sensory directions share a source model family, and the other operations retain their own source identities.

A request supplies the translator identifier, recording identifier, input kind and complete window. Spike requests must also supply the original ordered unit inventory. Mechanics requests must supply the recorded channel names, physical units and source frame. Timestamps can be declared for checks of time spacing, and an input-origin field distinguishes a caller's claim of recorded, source-derived, mixed, decoded or simulated input. That origin label remains a declaration rather than independent measurement authentication.

The response identifies the output's source and checkpoint, names and units, time and reference kind. It states that predictions are not observed data. Expected sensory activity remains ``model-expected activity.'' The response says that it is not a neural stimulation command and is not an inference of missing spinal activity. The joint operation also returns its training-range report. These fields matter because downstream software may otherwise treat a plausible number as permission to make a stronger claim.

The portable adapter runs an exported model for a caller. It reads the exported
catalog---the saved list of model descriptions---rather than rebuilding that
list from saved training states. It loads the plain numerical model exports
and checks that their declared training-state fingerprint matches the catalog.
A fingerprint is a computed identifier for file contents; this check compares
the declared identities. The adapter then applies the same request-checking
rules as the existing execution interface, preserving the source, unit, frame,
time-spacing and integer checks. The model can therefore run without loading
the training software. The recipient still needs a compatible execution
environment and the supplied numerical software.

The executable examples use actual included source-replay windows. Each example gives a request and an expected response. The historical live check compared all seven operations within the declared numerical tolerance of 0.000002 and required the other response metadata to match exactly. That tolerance concerns agreement between implementations, not error against a new biological recording. A wrong recording, reordered units, wrong mechanics units or frame, broken time spacing, fractional recorded count or stale model identity is a reason to reject the request. Rejection is part of the interface's functionality, because accepting the wrong contract would create misleading outputs.

Concurrent execution does not merge clocks. If the application runs motor and sensory examples together, each result stays on its own recording's clock with its own offset. The display can coordinate playback as a user-interface choice, but that coordination is not a measured physiological synchronization. A future matched multilevel recording would need its own verified common time base before NeuroTranslator could claim a joined chain.

Another AI or programmer can inspect the schemas, sources, example requests and model arrays, then execute the defined operations in a compatible environment. That is the handoff's practical portability: the same accepted computation can run with the same declared meaning. The archive should not be assumed to supply arbitrary biological skill programs.

\section{Separate attempted-finger and simulated-keyboard tests}

These October 4 studies have their own participants, sources, decisions and bodies. They are not the LINK comparison or replications of its improvement claim.

\subsection{Recorded Finger Intent}

The public recording came from one participant with tetraplegia, a condition
that affects movement in all four limbs. The experiment cued the participant to
attempt a finger movement. Its labels name those attempts. They do not supply a
measured path of the finger through space, a measured muscle command or a
recording of the participant playing an instrument.

The three recording dates used for the later test were May 16, May 20 and May
23, 2016. They are three dates from one person. The 150 later examples are not
150 people. These dates contain posterior parietal cortex recordings, from a
named region of the brain. They do not contain the matching motor-cortex,
spinal, peripheral-nerve or muscle recordings needed to observe the whole route
from brain to hand.

The software counted recorded events in eighteen completed time bins. Each bin
was 50 milliseconds long. Together, the bins covered the interval from 0.1 to
1.0 seconds after the experiment's Go signal. The decision was released only
when that window had finished. The earliest release was therefore one second
after Go, not an instantaneous response.

A new recording date also needed examples. Fifty early examples with known
attempted-finger labels were used to adjust the model to that date. This is
local calibration: using examples from where a model will be used. The later
50 examples from the same date supplied its test. Changing from twenty to fifty
early examples changes both the amount of data and how recently those examples
were recorded; it does not isolate the effect of example count alone.

The study kept each recording's unit identities separate. A unit number on one
date was not assumed to name the same cell on another date. When an observation
was missing, the software kept it unavailable. An unavailable channel did not
become a measured channel with zero events.

\subsection{Finger-Intent Model Comparison}

The study compared seventeen fixed candidates. Earlier validation dates were
used to select the model before the new test examples were opened. The selected
model was a simple classifier: software that chooses one category from ten
attempted-finger categories. It correctly chose 116 of the 150 later labels.
The three date scores were 35, 44 and 37 correct choices out of 50. A population
model given no early local examples made 11 correct choices out of 150.

This result supports a specific calibrated mapping. It does not show that the
same cells were identified across dates, that an unfamiliar person needs no
examples, or that one model works for every organism. It also does not show that
a large model would be a better use of this machine or dataset. The competitors,
including the less successful models, remain in the evidence record.

The recording contains cued attempts. The effect of seeing the cue has not been
separated from the effect of attempting the movement. A correct category choice
therefore cannot establish that the software read a general physical skill or
that it reconstructed every biological command needed to carry it out.

\subsection{Command Selection and Muscle Control}

The body controller received the actual predicted category. A wrong prediction
stayed wrong. The software did not look at the source answer and swap in the
correct finger. Four right fingers could command four keys. Left fingers and
the right thumb were retained in the record but could not command that body.

The body was an openly available, generic human-hand simulation with 39 muscles.
A previously trained policy supplied muscle excitation commands, numbers that
control the simulated muscles. The wrist, thumb and movements that spread the
fingers were held by braces. The requested motion used fixed fingertip targets
and a half-second press-and-release pulse. These targets were engineering
choices, not measurements of what the participant's muscles did.

The body was initialized once for each whole recording. After that, positions
and velocities advanced through the simulator. The controller did not overwrite
the hand's pose to make a press look successful. Key contact, travel and release
were checked in the physical model.

The body carried out all 54 supported commands it received. That sounds perfect
until we ask a different question: did those commands match the source attempt?
Only 41 of the 60 supported right-finger attempts led to the correct contact.
There were 13 unintended presses. The body could follow a wrong instruction
accurately. That is why 54/54 command execution and 41/60 source-intention
success must stay separate.

Three controls make the distinction clearer. With every muscle command set to
zero, no supported source attempt produced the right contact: 0/60. A reference
control was given the true source labels and reached 60/60. That control measures
what this body and controller can do when given the answer; it is not neural
decoding. A third control deliberately assigned neural decisions to different
trials. It reached only 4/60 source-intention contacts even though the body still
carried out all 54 supported commands. The deliberate reassignment is an offline
comparison, not a new recording of a nervous system.

All source trials stayed in these comparisons, including unsupported attempts
and mistakes. The recorded decisions were computed offline, then released on the
original source clock after their windows had completed. Watching the replay
at a normal pace does not make it live brain acquisition.

\subsection{Simulated Musical Phrases}

The muscle policy had already been trained in a separate phrase-and-chord task
before the recorded-neural connection was tested. That study asked whether a
simulated hand could carry out a requested sequence of notes through muscle
commands. The policy received note goals, timing and fixed smooth fingertip
paths. Recorded neural clips did not drive the original phrases, and during
the later held test it received no teacher actions.

The 32 episodes combine eight phrase orders, two tempos and two starting poses;
they are not 32 people or animals. The generic four-key simulated hand keeps
its wrist, thumb and finger spreading braced and uses scripted smooth fingertip
targets. These episodes test the goal-driven policy; the recorded-neural
classifier above answers the separate attempted-finger question.

The held cases used phrase orders, irregular rhythms, chords, tempos and starting
poses reserved before tuning. The selected corrected policy completed all 32
phrases. It produced 704 intended strikes and 192 chords, released the keys on
time and made no extra presses. The starting policy completed no whole phrases.
The first new fit completed 20. The named reference controller completed 24.
Zero muscle excitation completed none. Independent numerical replay checked the
saved trajectories and scores.

The original evaluator ended with exit code 143. That means it was terminated; why it stopped is unknown. It had already saved the complete closed results and final output. A separate full replay checker exited 0, meaning that checker finished successfully. These two facts must both be reported. The virtual spring-key limits are soft constraints, not rigid stops. Key travel can therefore exceed the nominal 6 millimetres. The saved poses and travel were not clamped, meaning forced into an allowed range, or invented to make the display look better. This remains a computer simulation.

Local training of that muscle-policy study took about 2 minutes 36 seconds.
The later neural-candidate training and validation took about 24.7 seconds.
Those are saved software runtimes. Recording acquisition, labeled calibration,
physical checks and independent verification require additional time. Neither
runtime says how long a living trainee would need to learn.

These are useful results in a constrained numerical body. They do not show that
the participant played those phrases, that a pianist's neural skill was uploaded,
or that a living hand learned music. They show that a saved learned policy can
carry out specified goals through simulated muscles, and that a separate
recorded-neural classifier can supply a narrower set of finger goals.



\section{Human Arm Simulation}

A predicted movement quantity can be statistically accurate yet lead to an unsuitable physical command. Muscle-and-skeleton simulation provides a separate check on that gap. The human-model branch uses an open MyoArm asset from the MyoSim/MyoSuite ecosystem and runs its dynamics through MuJoCo. These are existing community tools, not anatomy or physics invented by this project.\cite{myosim2026,myosuite2022,mujoco2012}

The model has 63 numerical muscles that produce forces, with internal responses that change over time. The experiment explicitly converts two velocity components into shoulder-and-elbow angle goals, while other independent joints are braced. This reduction is a choice about a bounded engineering test. It does not reconstruct a whole posture from two numbers or identify a macaque-to-human neural correspondence. The target velocity and human body remain distinguishable parts of the chain.

The controller applies muscle excitation commands that stay within declared nonnegative limits. Excitation is the requested drive supplied to the numerical muscle. Activation is the muscle model's changing internal response to that drive; the two are not the same value, and activation does not change instantly. The controller accounts for the linked shoulder-girdle loads through a calculation that relates the modeled forces and motions to the controlled joint angles. The simulation updates 500 times per second. It does not obtain success by overwriting joint angles with the desired pose, and it does not use recorded movement as feedback. Extra artificial joint-turning forces do not replace the muscle drive. Recorded position-derived motion is reserved for evaluating the resulting path.

The correction branch uses feedback about the simulated body's position and movement, plus speed and acceleration limits. These can change the command target itself. That creates two legitimate but different scoring questions: how closely does each branch follow its own command, and how closely does each follow a common scoring reference derived from recorded movement? Reporting only the first could exaggerate the correction's apparent value by ignoring the altered target.

In the bounded 500-bin replay on the new recording, the typical corrected tracking error relative to its own command is 0.0374 radians, compared with 1.5803 for the raw command sent without feedback correction. Against the common recorded-motion reference after conversion to this body, corrected error is 0.26302 radians and raw error is 1.49325. Here ``typical error'' uses a root-mean-square calculation: errors are squared, averaged and square-rooted, so larger deviations contribute more heavily. A radian is an angle unit, with a full turn equal to about 6.28 radians. The common-reference comparison makes the target difference visible.\footnote{Local evidence: \texttt{artifacts/motor\_transfer\_physics\_metrics.json}; evidence-matrix claim \texttt{new\_day\_muscle\_retargeting}.}

The corrected maximum anatomical-limit penetration is about 0.000406 radians, below the declared 0.001-radian numerical tolerance. The raw replay violates that tolerance in 75 of 500 frames. These are model-coordinate checks under the numerical setup. They do not prove that a real subject's tissues experience the same loads or that a decoded command is biologically correct. A simulator's constraint is only as relevant as the model and task that define it.

The improvement belongs to the bounded controller after movement has been converted between bodies under the explicit rules. It remains a genuinely muscle-driven trajectory, and the scoring reference was not secretly supplied as feedback. This separates decoding quality from physical correction: the changed stage explains the improved trajectory without claiming recovery of hidden spinal signals or learning in the nervous system.

The human branch remains a reduced adapter. Its success cannot erase the failures in the native macaque branch, and it cannot justify sewing or fingertip manipulation in a model whose tested control concerns two larger arm coordinates. Physical grounding starts with a named task and a named body, then expands only when the additional coordinates, contacts and evidence are actually present.

\section{Macaque Arm Geometry and Remaining Errors}

The separate macaque body comes from the model family used for the source's arm analysis. It is closer to the relevant anatomy than a human arm with two controlled angles. It has seven source coordinates and 39 muscle-tendon paths. A coordinate is an independently named angle or motion setting in the model. A muscle-tendon path is the modeled route from one attachment to another, including how the route passes around nearby structures. The model's hand is rigid: it has no separate finger motions.\cite{chowdhury2020,limblabmodel2026}

The retrieved revision postdates the source paper. Source-family membership is useful, but it does not prove that this revision is the exact model used for the original subject analysis. Matching joint names and units is not enough to certify matching specimen geometry. The project preserves that uncertainty. A later improvement in numerical agreement should not retroactively be described as recovery of the original subject's true anatomy.

The retrieved original model also contains zero-valued rate properties in its older muscle formulation. With those settings, it cannot calculate the required force-driven motion while keeping its numerical states finite. A separately named numerical fork restores 117 properties to defaults provided by the modeling software. A fork is a distinct saved version, not a silent edit of the original. The original geometry and the assumed default-rate version remain separately identified.\cite{opensim2018}

Using software defaults makes a numerical experiment possible, but does not measure the subject's physiology. The assumed rates are not recovered activation times, tendon properties or force capacity. This distinction is especially important when a controller fails. The failure may reveal a limitation of the numerical assumptions or control design, but it cannot be assigned to a particular biological mechanism without further evidence.

Three native-body controller candidates are retained as rejected experiments. The first concerns an earlier source case with a different way of combining errors across joints. The actual-tendon-force feedback candidate measures modeled tendon forces and uses them in correction, yet still fails its reserved source comparison of tracking across all seven joint angles. This case had already been inspected before the initial-speed protocol was corrected, so the corrected rerun is development evidence rather than a fresh uninspected test. On its reserved source trial, its typical error is about 0.05505 radians, compared with about 0.04512 for the tonic baseline, a simpler controller that maintains its defined steady muscle drive. This error calculation squares all angle errors, averages the squared values across coordinates and time, then takes the square root. It does not average seven separately calculated typical errors. The more elaborate feedback is worse on that comparison.\footnote{Local evidence: \texttt{artifacts/native\_force\_metrics.json}; evidence-matrix claim \texttt{\mbox{native}\_\allowbreak{}\mbox{tendon}\_\allowbreak{}\mbox{feedback}\_\allowbreak{}\mbox{rejection}}.}

The short-horizon candidate tries possible commands in a copied model state before choosing what to do. It improves a synthetic validation comparison, from about 0.01989 radians for the tonic baseline to about 0.01270. On the reserved recorded-source case, however, it produces about 0.05971 radians versus about 0.03263 for tonic. It is rejected. Better performance on the synthetic validation case does not permit substitution for the failed source case.\footnote{Local evidence: \texttt{artifacts/native\_horizon\_metrics.json}; evidence-matrix claim \texttt{\mbox{native}\_\allowbreak{}\mbox{horizon}\_\allowbreak{}\mbox{rejection}}.}

The copied-state planner also takes about 4.52 seconds to plan each 20-millisecond source interval. That is an offline experiment. Showing the resulting cached trajectory at normal playback speed does not turn its planning into real-time control. The checks that isolate copied state and restore model-owned control functions are useful implementation evidence, but the latency and tracking failure remain part of the result.

These failures give the architecture a concrete next problem: establish a stable, applicable numerical muscle-driven body and controller before asking it to evaluate learning rewards. The successful human-model movement-conversion branch does not answer that problem, because it uses different anatomy and a reduced task. A source-looking body is a better starting point for a source task, but its numerical behavior still has to earn its place through measurement, calibration and appropriate controls.

\section{Muscle Paths and Physiological Validation}

The native geometry audit reproduced an invalid muscle detour around a wrapping surface associated with the ulna, one of the forearm bones. A wrapping surface is a modeled obstacle that keeps a muscle path from passing straight through a bone or another structure. Here, a doughnut-shaped surface, formally called a torus, produced the detour in both the original geometry and the default-rate version. Making the muscle rates finite had therefore left the route problem unresolved.

The correction uses a rule derived from attachment topology. Topology here means which attachments and path segments are connected, rather than how an anomalous frame happens to look. The rule restricts three named wraps to the unique original segment from the humerus, the upper-arm bone, to the ulna in the forearm, using the software's numbering convention that starts at one and includes both ends of the selected range. No muscle, bone, attachment point or wrap shape is deleted to make the image nicer. The saved version names the three changed ranges and preserves the original model separately.\cite{limblabmodel2026,opensim2018}

Across seven inspected anomalous frames, the correction removes 14 contacts outside that wrapping surface while retaining all seven coordinates and 39 muscles. It is tested on training samples, on already inspected diagnostic frames, and on a newly reserved same-session geometry window. Those supports have different evidential status and are reported separately. The checks also inspect the saved rotations and positions of rigid body parts and changes in path length as the joint angles change. They establish specified numerical behavior, not an independent anatomical measurement.\footnote{Local evidence: \texttt{artifacts/monkey\_wrap\_candidate\_metrics.json}; evidence-matrix claim \texttt{\mbox{topology}\_\allowbreak{}\mbox{range}\_\allowbreak{}\mbox{candidate}}.}

A later comparison finds a broad length mismatch between this model family and the source-model-derived reference lengths. The project fits one common length factor using 128 training poses and all 39 ordered references. A pose is a set of model coordinates at a stated instant. To choose the common factor, multiply each model length by its matching reference length and add those products, then divide by the sum of the squared model lengths. That calculation chooses the single multiplier with the smallest summed squared length error on the training poses. It does not fit a separate shape to every pose or every muscle. The factor is 1.2413815, so the numerical scaling is about a 24 percent increase. It is estimated from training poses only and then kept fixed for separate validation and test poses.

On those separate poses, the typical length discrepancy falls from about 27.75 millimeters to about 0.552 millimeters. This is substantial agreement with a specific reference. The reference is itself produced by tracking and muscle-and-skeleton analysis; it is not a collection of directly measured fiber lengths. Consequently, the fit demonstrates improved agreement between modeled quantities under the stated source conventions. It does not identify individual bone dimensions, tissue mass, strength or activation rates.\footnote{Local evidence: \texttt{artifacts/source\_scale\_frozen.json} and \texttt{artifacts/\allowbreak{}source\_\allowbreak{}scale\_\allowbreak{}metrics.json}; evidence-matrix claim \texttt{source\_scale\_fit}.}

A separate saved version applies the frozen factor through the official OpenSim scaling tools. This is stronger numerical evidence than multiplying exported lengths after the fact: the model itself is scaled and its computed lengths are compared with the prediction from the single factor. Their maximum disagreement is about 6.13 micrometers. A consistency check on the modeled moment arms, the leverage distances relating joint rotation to muscle-path length, is within about 0.181 millimeters. These remain checks of the numerical model.\cite{opensim2018}

The scaled version retains the assumed original masses and applies the corresponding geometric updates to rotational resistance and muscle-length parameters. No mass, inertia or physiology was measured for the subject through this operation. Dynamics and control are blocked for this scaled geometry-only version. Its 400 displayed observed and decoded poses reuse diagnostic frames and are not a fresh evaluation of decoded-pose accuracy. The number counts displayed model states; it does not add independently recorded subjects or a new test.

\begin{center}
\begin{minipage}{\linewidth}
\centering
\makebox[\textwidth][c]{\includegraphics[width=160mm]{figures/length_fit.pdf}}
\captionof{figure}{One common length factor, fitted on training poses, reduces mismatch on separate validation and test poses. The reference lengths come from source tracking and body-model analysis. Agreement with those lengths does not establish physiological calibration.}
\end{minipage}
\end{center}

The order of events also matters for explanation. The known observed wrap anomaly occurs later than the evaluated window of the failed lookahead controller. The later geometry correction therefore does not establish that wrapping caused the earlier controller rejection. A useful repair can coexist with an unresolved failure. Treating every newly found defect as the explanation for every old result would replace testing with a convenient story.

\section{Template Anatomy and Participant Anatomy}

A template brain is a useful reference shape. It is not automatically the brain
of the person or animal whose electrical activity was recorded. Registration
means establishing the spatial relationship to that actual subject with
appropriate independent evidence. A template can illustrate where hand motor
cortex generally lies while leaving exact electrode and neuron positions
unresolved.

The current source reconstruction reproduces 122,284 template points, 244,760
triangles and 81,924 atlas values. The points are vertices: locations joined by
triangles to make a surface. An atlas supplies region labels within that
template. Three source MRI planes also match the named source. Those are useful
checks of source correspondence, rather than proof that every biological
interpretation of the image is true.

The mesh checks ask a narrow question: are these surfaces joined in the
specified way? The two mid-cortical surfaces pass the specified closed-surface
connectivity checks. Their Euler characteristic is $+2$: count the vertices,
subtract the edges and add the triangular faces. That number summarizes mesh
connectivity; it does not measure how closely the shape matches the recorded
subject. The optional region boundary still contains 432 open edges, 838
inconsistent triangle directions and six invalid point neighbourhoods. Those
defects remain even though the other surfaces passed their own checks. A
second reader using the established nibabel imaging library was unavailable.

What can this template tell us about the recorded subject? The source checks
do not register that subject, identify individual neurons, establish a region's
function or supply an absent peripheral nerve measurement. A colour needs a
declared value, source and time; brightness cannot improve the recording's
spatial evidence.

A lifelike, tangible three-dimensional model would make the anatomical
comparison inspectable. To judge it, we need a series of identifiable real-world
references covering the relevant structures, views, poses, conditions and scales.
One convenient picture cannot cover them all. Each reference must identify its
original source, organism, body part, task, measurement method, units or scale
and the exact feature it supports. Derived geometry and synthetic images retain
their labels; neither supplies independent ground truth.

The comparison has to follow a part into the assembly. A plausible forearm
could still meet a floating muscle, an implausible interface or an impossible
joint pose. Comparable pose, projection and scale let reference and model views
test those relationships. Close-ups, tilted orthogonal views along perpendicular
axes, biological identification colours, transparency and cutaways can reveal
them; the reveal still needs measured agreement.

Function adds another comparison. A displayed signal must keep its site, source
clock, measured quantity and availability, with recorded channels, predictions
and simulated muscle activation separately labelled. Body claims need held-out
movement, joint and force comparisons wherever those quantities are claimed.
Those physical measurements, full lifelike-anatomy review and independent
part-plus-assembly reference acceptance remain missing.

The four historical figures here are diagrams and numerical plots with their
original evidence scopes. They explain relationships; they supply neither
newly validated lifelike 3D biological models nor the missing anatomical and
reference-series comparisons.

\section{Application Performance}

The graphical user interface, usually called a GUI, should help a reader follow
the evidence chain. Choosing another source must change the source label and
clock rather than leaving a convincing stale value on screen. A gap in recorded
support must remain a gap. A camera reset must not change the scientific meaning
of a pose. Playback should show which values were measured, predicted or
simulated, and whether they are current for the selected source interval.

The bounded original-control checks contain evidence for source switching,
source gaps, replay and camera/restoration behaviour. Their acceptance is scoped.
It is not whole-GUI acceptance or a measurement of better performance. Full
source-clock, geometry, material, renderer-lifetime, functional-map, device,
anatomical-reference and offline criteria remain binding.

The separately checked hand replay is deliberately a saved-data component.
It automatically loads only an explicitly identified current cached replay.
The demonstrated cache has 336 frames with a 20-millisecond source interval.
Play, pause, seek, single-frame stepping and restoration have component-level
evidence. If no cache exists for a selection, the component does not quietly
request a new run. Missing values appear as dashes. An unavailable or failed
source clears the prior replay and geometry instead of presenting it under a
new identity.

Its saved development example shows 22 of 22 correct strikes, 6 of 6 chords,
22 of 22 timely releases, 34-millisecond worst timing error and no extra strike.
A chord means more than one requested key at a time. Those values describe
one constrained simulation example. They are not a replacement for the separate
32-phrase evaluation or the recorded-intention-to-contact study. Separate human
neural clips do not drive the goal-driven hand policy.

Opening the online candidate requires its whole directory. HTML is the entry
file a browser reads to assemble the page. Moving opaque resources out of it
reduces that file from 113,298,857 to 34,295,597 bytes, while the complete
candidate directory remains 93,831,812 bytes. The resources have moved within
what must be delivered; the smaller entry file does not contain the application
on its own.

The separate static offline export preserves the original source and occupies
113,581,899 bytes. Its supported static content can be accessed without the
running online service, but research views that require the application service
interface, or API, remain unavailable. Testing that original-source export and
testing the whole-directory online candidate are different tasks.

File sizes establish a static packaging result. They supply no measured change
in page-loading time, interaction speed, memory, crashes or graphics-processor
performance. The candidate's browser route was unavailable, so the candidate
remains unpromoted; an asset archive and a source check cannot fill that missing
interaction test.

The practical consequence is also scientific. An old trace under a new source
label suggests a relationship that was never observed. A zero in place of a
missing measurement asserts activity the recording cannot support. A saved
replay presented as a fresh model run hides its development status. Delivery
and playback have to preserve those distinctions before the reader can judge
the displayed result.

The human EEG/ESG/ENG/EMG preview has its own evoked-stimulation context and clock.
The macaque-derived human-arm branch has its own prediction and simulation
chain. The GUI must not combine their appearance into a claim of simultaneous
voluntary cortical, spinal, peripheral and muscle control. A hypothetical
biological teaching view must likewise stay hypothetical. The screen can
explain a missing link; it cannot measure it by drawing an arrow.


\section{Earlier Viewer}

The application is intended to let a reader follow the chain rather than accept a still image. Source controls select a recording or an observation view. Replay shows recorded inputs and appropriate references. Model outputs identify what was predicted. Body views distinguish reference geometry, decoded geometry and numerical muscle-driven motion. Their coexistence makes comparison easier, but their labels remain necessary.

A moving skeleton can come from several routes. It may show joints obtained from upstream tracking and analysis. It may show joints predicted by a neural model. It may show the outcome of a numerical controller that acts through muscles. In the first two routes, the displayed pose can be set directly from coordinates for inspection. In the third, the motion follows the simulated forces and dynamics. Only that third route tests whether the chosen controller can move the numerical plant under its rules.

Playback has another distinction. A cached trajectory has already been computed and can be displayed smoothly. Paced playback controls how quickly the user sees stored frames. It does not report the time originally needed to compute them. The native lookahead result is the clear example: a normal-speed replay can show an offline plan that required seconds for each short interval. A real-time claim would need actual end-to-end timing under the deployment conditions.

The natural-language interface translates supported user requests into structured task descriptions. A person can ask to practise a slow precision reach or replay a supported recorded task. The compiler identifies a task, parameters and whether the available implementation can train or only display it. The structured description is written as JSON, a common text format with named fields and explicit values. This makes the request inspectable before execution.

The language interface turns a supported request into a specification the application can execute. That is its demonstrated role. Arbitrary prose does not become a verified neural program through the same step. A sewing request still needs sewing observations; a puzzle-manipulation request still needs tested object contacts and individual fingers. A fluent response cannot supply capabilities absent from the files and body model.

Sewing, Rubik's Cube solving and breakdancing remain future composite tasks. Sewing would need tested needle, fabric and fine-finger interactions; puzzle manipulation would need object contacts and controlled fingers; breakdancing would need whole-body balance and changing contacts. Each would also need recordings and outcomes appropriate to its learning claim. A reach primitive can be a useful component without becoming evidence that the complete task has been learned.

The skill-practice branch adapts two simulated controller gains. A gain is a number that determines how strongly the controller responds under its chosen rule. The learner compares outcomes using verified tracking error, effort and joint feasibility, and records reward/error curves and body replay. Its competitor is the existing controller, which is itself a candidate. Improved gain values can demonstrate a better numerical parameter choice for this setup; they do not demonstrate biological plasticity, retained human skill or a general reinforcement-learning trainer.

An interface can name a supported practice task, show its numerical history and body model, and expose where the controller failed. Showing the assumptions alongside those outcomes helps the reader judge the software. Whether a living trainee retained a skill after assistance stopped requires a different test.

The application also preserves missing observations. A source view that lacks motor spinal channels should say so. A human surface preview should remain an evoked experiment on its own clock. A hypothetical organism render should remain a concept. This is not merely cautious captioning: these labels determine whether a downstream user interprets the system as replay, prediction, control or teaching. A truthful live demonstration can be ambitious while keeping those operations visibly different.

\section{Digital-Fly Validation}

A detailed neural map can support a useful prediction and still leave important
predictions wrong. Shiu and colleagues tested a fly connectome-derived model
with real perturbations and behavioural outcomes. Their tested successes
coexist with missed effects involving Phantom and Usnea and wrong predictions
for G2N-1 and Usnea in water-related silencing comparisons. Those are named
circuit examples, not proof that the whole fly is wrong in every setting.
Missing baseline activity and chemical mechanisms are among the discussed
limits, including neuropeptides: small chains of amino acids used as signals.
Those possible limitations are not uniquely established causes for each error.\cite{shiu2024}

This is the positive standard too: a useful approximation earns its credibility
through a test tied to a specified prediction. It does not need a claim of
complete biological identity. A model can help science by showing where a
structural explanation succeeds and where an additional measurement or mechanism
is required. Retaining failures makes the next question more specific.

The official Eon technical account from March 2026 describes a different
engineered integration. It combines a fly-derived brain model with body
controllers and hand-chosen interfaces. It acknowledges largely missing state
and plasticity and internal dynamics that were not yet validated against
relevant known patterns.\cite{eon2026} That account supplies dated developer
disclosures, not an independent full-animal reproduction test. A sparse motor
interface can produce an impressive body demonstration while leaving the
biological mechanism unproved.

Two other source examples show why added software must remain visible. The
pinned fly-chess code includes symbolic board inputs, a reduced mathematical
network, a trained readout and external legal-move logic and lookahead. The
pinned DoomFly code uses a different male nervous-system map, adds enemy
geometry obtained from the game engine and trains an action readout. Its
authors describe the failure of the originally intended looming-vision route.
These are inspected software sources and author reports, not measured evidence
that a living fly learned chess or Doom. Their playing strength was not
independently reproduced here.\cite{flychess,doomfly}

There is no verified primary evidence here that a FlyWire-derived brain
performed something no fly could ever do in a lifetime. The source-supported
criticism concerns the engineered inputs, controllers, mappings and training
that supply the software's capabilities. Their biological interpretation
remains unvalidated; a more dramatic impossibility claim would add no evidence.

Body simulation has its own failure modes. NeuroMechFly v2 reports that simulated
leg tips can sit slightly below the ground because of numerical errors or
instabilities. Its authors use a 0.05-millimetre margin to avoid a false detection
of excessive stretching.\cite{neuromechfly2024} That is a documented contact
artifact. It is not proof of gross tunnelling, perpetual motion or a FlyWire
brain's impossible ability. A small explicit error is more scientifically
useful than a larger invented one.

Published first-hand accounts of other simulations describe creatures exploiting
numerical integration or contact errors that supply artificial energy. A model
can then appear to move well for a reason that real physics would not permit.
The error needs a physical repair and a new check; an improved optimizer would
otherwise learn the error more effectively.\cite{surprises2018}

An objective can fail even when the physics is not the problem. In the reported
CoastRunners example, the agent kept circling reward targets rather than finishing
the race. Its score improved while the intended outcome did not.\cite{specgaming2020}
This is why a reward curve needs independent task outcomes. The number optimized
must be checked against what the experiment was supposed to accomplish.

The protein example has a similar lesson about scope. The original AlphaFold3
paper reports 62.9 percent success on its specified structural-interface benchmark,
with a 95-percent interval of 51.4 to 73.5 percent. That result averages over
65 interface clusters from 71 complexes and uses a structural agreement threshold
and confidence-ranked samples: 1,000 seeds times five generated samples per
target. The model generates these candidate structures by refining noise in
stages, a method called diffusion sampling. DockQ is the structural-interface agreement score, and success means
it exceeds 0.23. These generated samples are not independent biological
experiments. Its 37.1 percent complement is the complement
for that benchmark, not the percentage of every protein that fails to function
in real life. Structure and biological function are different endpoints.\cite{alphafold32024}

Crystallization is a real obstacle for some measurements, but it is not the
only route to a structure. Cryo-electron microscopy uses electron images of rapidly frozen samples. Nuclear magnetic resonance measures how atomic nuclei respond under controlled magnetic conditions. These provide other methods, each with limits. We do not infer that most
unknown proteins are unknown because they fail to crystallize.\cite{pdbmethods}
Missing measurements are an important problem without inventing a universal
failure percentage to describe them.
Training can also draw on generated examples without every example containing neural signals, muscle activity and movement measured together. A simulation can supply an example; a model can supply a predicted training target, called a pseudo-label. Those values remain generated or derived. Agreement with the model or musculoskeletal simulation that produced them is not an independent measurement of living biology. Generated test examples can check a specified calculation, but they cannot fill a missing measurement. Biological claims still need comparison with independent measurements in the stated domain.

There are stronger positive examples for brain--body interfaces. Lorach and
colleagues linked measured cortical signals to spinal stimulation and evaluated
actual standing and walking in one participant with chronic tetraplegia. The
movement happened in the person rather than only on a display.\cite{lorach2023}
The interventional outcome supports the stated single-participant interface;
it does not establish every possible decoder, biological skill transfer or
complete nervous-system virtualization.

Our own claim should follow the same discipline. Ask which quantity was observed,
which mechanism was assumed, which outcome was independently measured and which
comparison failed. Then state the tested domain. A realistic animation, a
correct guess and a validated biological model remain different things even
when all three look persuasive at first glance.


\section{Prediction and Reward-Based Training}

For a particular neural-to-body prediction, measured paired examples tell us which observations belong together. A supervised model receives an input and its corresponding desired output during fitting, then faces examples withheld under the stated rules. This gives the signal-to-target claim a reference before a controller changes actions to earn a reward. A movement component can also learn first from qualified movement-only data. That complementary route can supply a useful starting representation without supplying neural labels that were never recorded; its transfer still has to be tested.

The information supplied at prediction time defines the claim. An EEG-only
predictor may use scalp observations and declared available context; it cannot
secretly receive the muscle or movement target it is supposed to predict.
EMG-only, movement-only and explicitly combined predictors answer other useful
questions and need their own labels and matched comparisons. Muscle or body
measurements may help teach a model during an explicitly declared training
procedure without remaining available to the deployed EEG-only route. Such
training information is called privileged because the deployed predictor does
not have it. Independent evaluation must verify that the privilege ends where
the claim says it ends.

Physical correction comes next, with the decoder held fixed. The controller can ask whether the decoded quantity leads to a feasible bounded movement in the numerical body. Holding the decoder fixed makes the comparison easier to interpret. If corrected motion improves, the improvement belongs to the control stage rather than being confused with a newly trained neural representation. The current human-model branch supplies a limited example of that separation.

Only after the relevant body and controller are stable does a reward-driven correction rule become a candidate worth testing. Such a rule would add a bounded correction to the existing command. Reinforcement learning is the general approach of changing action choices according to evaluated outcomes. In this project, any future such refinement must start with outcomes that the simulator can actually verify and with a fixed comparison against the supervised decoder and baseline controller.

Possible rewards include tracking a specified target, respecting joint and contact constraints, controlling effort, limiting abrupt changes in muscle-drive commands and completing a defined task. These are numerical criteria. Their meanings depend on the body model, sensors and task. A tracking reward can verify movement toward the simulator's target; it cannot independently verify the unobserved spinal command that supposedly produced that movement.

A controller can appear excellent when its answer is hidden inside its input. Feeding the recorded reference into action calculation changes the question being tested. Keep hidden test targets out of deployment. In the present bounded muscle replay, recorded position-derived motion is reserved for scoring and is not supplied as feedback. Future training and evaluation must disclose the information the agent receives.

Robustness also needs more than one successful trajectory. Future tests should hold out recordings and relevant body/task conditions, and should compare changes in morphology, friction, delay and measurement quality where those variables apply. A learned correction that exploits one plant's convenient numerical quirk may fail under another plausible setting. The test should be selected for the claim being made rather than expanded merely to accumulate impressive-looking cases.

This development order allows physics to reject impossible or unsuitable commands without pretending that physics identifies the true neural interpretation. It also prevents a reward curve from being presented as proof of biological learning. The current two-gain practice learner is a working software lesson interface. A general adaptive controller and a biological trainer remain separate future achievements, each with its own controls and failure criteria.

\section{Matching future organisms to actual tasks}

The broader architecture needs body models that fit the recordings and tasks. An organism-specific adapter preserves actual anatomy rather than blending convenient features from several species. Shared functional labels such as limb role, task direction, contact, effort or sensory quantity can help organize comparisons. They do not establish identical muscles, neurons or learning mechanisms across animals.

The virtual-rodent work is a relevant precedent because it combines animal behavior, embodied modeling and analysis of recorded groups of neurons. It supports the idea that an appropriately grounded body model can help ask questions about recorded activity. It does not imply that a human arm, macaque arm, rodent body and fly body can exchange neural commands without evidence.\cite{virtualrodent2024}

Choosing a body begins with actual anatomy, available recordings and a task the model can represent. A hypothetical animal drawing is not a validated body model or a recording. A transparent render cannot supply hidden nerve measurements or show that an animal can tolerate a proposed interface.

Species selection remains a future decision. Care, development, available tools, ordinary behavior and interface burden all bear on it. A specified human biological feature or a tissue connection needs a comparison with direct measurements; the current companion supplies no biological outcome or repair result for those branches. A computational bridge must begin with the anatomy and data actually available.

A future rodent branch would require a supported rodent body and a task represented in matched recordings. A fly branch would need its own appropriate movement and neural sources and a body whose force-producing components have the correct declared interpretation. A human whole-body branch would require relevant contacts, muscles and task evidence beyond the current reduced arm. The source catalog can help locate candidates, but a candidate's presence in the list is not completion of these requirements.

Shared movement can provide a useful starting point before every source has
paired neural recordings. Homologous structures share an evolutionary anatomical
origin. Mechanically analogous structures can perform a comparable job without
that shared origin. Either can motivate transfer when the actual similarities
are stated: the joints and their movement directions, limb proportions, contact,
loads and time scales. Similar-looking limbs need not have the same force
capacity or move on the same time scale. A source-specific adapter must preserve
those differences rather than quietly treating them as equal.

A movement representation could first learn from qualified motion capture,
then be adapted using actual paired neural and body observations. Motion can
also be retargeted to a declared body under anatomical and physical constraints.
Here retargeting means calculating a corresponding movement for the new body;
the calculated movement remains derived, not a newly measured neural pairing.
Scans could constrain a specific anatomical gap after the selected model's own
provenance has been checked. A scan from another person remains an external
reference; it does not register the recorded participant or validate simulated
forces.

A separate current hand/arm proposal would test movement pretraining across
two people, with separately calibrated sensor-based movement descriptions. It
would compare target-only training, the proposed anatomical mapping and a
deliberately wrong mapping under matched resources. Its source and target
archives, sensor conventions and numerical implementation still require
admission; no new fit or successful transfer is reported here. Broader motion
collections would need their own rights, task coverage and original-recording
overlap checks before entering that route. Repeated copies of one source in
different collections cannot become independent training and evaluation data.

Comparisons across organisms should include behavior-only and within-organism controls. If a shared representation appears helpful, sessions kept outside fitting, animals and eventually species should test where its help persists. Anatomical similarity can motivate the comparison without supplying cell identity. Missing channels should remain unavailable or explicitly inferred with a stated uncertainty method, rather than being reconstructed into the observed-data column. Useful mechanical transfer would earn a claim about that tested mapping; neural equivalence, causal physiology, sensation delivery and biological teaching still need their own evidence.

The proposed middle ground is therefore functional coordination with preserved anatomy. A task may ask different bodies to move toward a target or control effort, while each body uses its own joints, muscles and sensor models. That makes the research question concrete: which representation helps under these adapters and controls? It avoids claiming a universal nervous-system language before the experiments show one.

\section{Next Experiments}

The next computational question comes from the failed comparison. Both
recurrent candidates miss the original no-harm and positive-neural-information
requirements, and all five added benefit claims fail. Lower average error
admits neither predictor. NeuroTranslation remains a proposed field;
NeuroTranslator remains recording-specific software.

A further proposal would have to connect supported observations to decoded
estimates, pass them through a declared body adapter and identify the feedback
returned. Simple baselines need the same tests. Illustrative motor replay or a
deliberately goal-driven simulation can examine that bounded computation if its
goals, assistance and simulator information are disclosed. A positive
neural-control claim still requires a predictor that passes the unchanged
criteria; the current candidates do not.

That leaves two questions for the sections that follow. \emph{Tests of Retained
Learning} asks what remains when teaching help stops. \emph{From a reach to a
testable prediction} develops the paired mechanism, and \emph{Measure both ends.
Test the translation.} specifies its observations and comparisons. Reward,
simulation and independently measured task outcomes need separate checks.
Physical calibration, synchronized recording chains, participant-specific
registration and generalization remain unverified. Better simulated control
would support that tested computation; biological teaching, retention and
repair still need their own evidence. Their absence leaves room for a clearly
bounded computational proposal, not a biological result.

\section{Tests of Retained Learning}

The biological learning ambition needs a different kind of evidence from model prediction. A trainer could provide commands, feedback or task information that improves a trainee's performance while connected. That is assistance. To claim learning, the experiment must show a change in the trainee's later independent performance under an appropriate test. To claim faster learning, it must compare the relevant practice, information and time budgets.

Biological and computer-based trainers need separate tests. A biological trainer would be an experienced partner under a specified linking condition. A computer-based trainer would supply assistance under a defined rule. Either could help a trainee while connected without producing learning that lasts. Compare their effects through the same measured task and matched conditions, rather than assuming that the identity of the trainer explains the result.

Experience itself needs a control. An experienced-to-novice pair should be compared with a novice-to-novice pair and with individual practice. A matched no-link condition keeps the relevant experience, exposure and setting while removing the proposed live signal connection. A trainee that practises alone supplies a training control; it is not the same thing as a previously guided trainee being tested alone after training. Groups need comparable task exposure before retention is assessed. Otherwise a benefit could come from extra practice, observation or assistance time rather than the proposed transfer route.

Time-shuffled or yoked comparisons can help identify whether the proposed signal alignment contributes. A yoked condition receives a matched amount of assistance according to another trial or rule, rather than assistance appropriately contingent on its own current state. The exact design belongs to the future experiment, but the principle is to make information timing and assistance budget separable. A simple controller and a behavior-only predictor are also important competitors. More neural measurements should have to show an additional relevant benefit.

Retention should be measured after trainer commands stop and again after an appropriate later interval. Transfer should ask whether improvement helps on a task condition, maze or body condition withheld from practice. Accuracy, consistency and independent control need explicit definitions. A new-looking maze image is not automatically a new scientific test. An equal target score does not prove equal difficulty across route designs or bodies. Those difficulty comparisons need their own measurements and rules. A well-executed rehearsed route can coexist with poor adaptation to a new route, and a one-time success during guidance can coexist with no later benefit.

Replay comparisons also need to identify what is being replayed. A proposed test can compare activity from the same trainee and matched session with stale or unmatched activity and with an inactive or sham comparison. A sham is a control that resembles the relevant procedure without supplying its proposed active contribution. These are conceptual study-design distinctions, not current animal results. They help test whether the matched signal matters beyond the existence of a device, a replay label or an extra source of help.

The persistent-interruption branch requires particular care. Suppose the trainee needs a standardized compensatory interface to express any movement through the interrupted route. Switching off that entire interface prevents motor expression whether or not the trainee learned. The retention comparison should remove trainer commands while retaining the necessary bridge and measurement conditions. A comparison in an intact group, or another independently justified test, may be needed to answer different learning questions. Bridge-off failure alone cannot establish absence of learning.

The corresponding positive claim also needs discipline. Success with the bridge present could still depend on active compensatory computation. The experiment should identify what remains constant, what the trainer previously supplied and what the trainee now controls. A result can demonstrate independent control through a necessary bridge without demonstrating that the interrupted tissue regenerated. Learning, assistance and regeneration are different outcomes even when they concern the same movement.

Humanized and graft comparisons add their own mechanism questions. A changed cell population or host--graft connection would need a specific definition and a justified measurement. It cannot be described as making the animal a miniature human or as supplying an unspecified universal repair. Disease-specific branches need disease-specific outcomes. These are proposed research comparisons, and the current application supplies none of their biological results.

The computational prototype contributes an instrument-building discipline to these future stages. It can specify tasks, record where sources and commands came from, preserve timing, compare predictors and expose physical failures. It cannot replace experimental responsibility for a living trainee or infer plasticity from a software reward. The claim worth pursuing is measurable improvement that remains when the trainer's contribution is removed under the right conditions. That is a harder and more useful claim than an impressive assisted performance.

\section{Tests That Could Reject the Proposal}

A proposal becomes scientifically useful when it says what would count against it. The broader NeuroTranslation idea should lose a comparison if it does not help. A simpler instrument may capture enough information. A behavior-only predictor may outperform added neural channels. A model matched to one recording may beat a shared model. Ordinary practice may produce the same retention as a proposed neural trainer under matched exposure.

Several current results already have this structure. The transferred fixed core had to beat every predefined control and did not. The shared candidate had to improve validation evidence against retained specialists and did not. Native force feedback and copied-state planning had to pass their reserved source comparisons and did not. These failures are not proof that all future variants will fail. They are evidence that these particular variants have not earned the broader claim.

For the next body experiment, a useful milestone is a stable numerical body appropriate to the source whose control comparisons survive appropriate cases kept outside fitting. If geometry, force behavior or timing cannot be calibrated sufficiently for the task, the project should report that limit rather than use a prettier animation. For the next recording experiment, a useful milestone is matched voluntary measurements with explicit time and quantity support. If the added levels do not improve a relevant comparison, their cost and burden need to be reconsidered.

For a future trainer, the decisive milestone is independent retention under matched conditions. If a trainee improves only during guidance, the result supports assistance. If retained improvement is no better than ordinary matched practice, the proposed teaching mechanism has not shown an advantage. If improvement stays on one practiced route but does not transfer, the claim should remain task-specific. Each outcome can still be useful when it is named accurately.

There is also a practical distinction between new information and more computation. A model cannot recover a verified label for an unmeasured quantity simply by becoming larger. It can make an inference, but the inference needs its own uncertainty and evaluation. If the missing variable is contact, specimen geometry, channel correspondence or upstream timing, a targeted measurement may resolve the problem more directly than further optimization on the same incomplete data.

We want the instrument to expose these limits early. That is a design judgment, supported here by the source-processing failure found in the audit and the failed native-body comparisons, rather than a measured biological result. Preserving the failures and each quantity's origin makes the next question easier to choose. Hiding them carries the same assumptions into a larger system that is harder to check.

The proposed new approach is therefore an evidence sequence: identify observations, preserve their meaning, fit a bounded model, check its predictions, test applicable mechanics, and only then test a specified teaching claim. The sequence is not guaranteed to culminate in instant skill acquisition. It is designed to reveal where the ambition fails and where an additional measurement or simpler solution actually helps.



\section{Terminology and Measurement Plan}

The shared names identify what is being tested. Interoceptomics combines
interoception and connectomics: how sensing conditions inside the body relates
to nerve-cell connections. NeuroTranslation names the proposed mapping between
defined neural situations; NeuroTranslator names the implemented software;
NeuroTranslation Network names the larger proposed design. Their roles were
introduced in \emph{Neural Signals and Physical Skills} and remain distinct.
These are working project names, not externally established standards or
evidence of a universal decoder or biological teaching. The definition
supplement retains the explanations and older technical aliases.

Four other meanings must stay separate. Physical scale concerns sizes and distances. Organizational level concerns whether we study a whole body, an organ, a circuit or cells. Computational detail concerns how much information one software calculation represents. Statistical averaging concerns how we combine scores or observations. An average across several output groups is not a measurement of a larger piece of anatomy. A small model is not evidence that it measures individual nerve cells. Calling a result macro or micro without saying which of these meanings is intended would hide the very distinction the programme needs to preserve.

Future body observation would synchronize the left and right tripod views of the hands and arms. A separate perpendicular depth-sensing view could be compared for what it reveals or misses. This software did not acquire or calibrate either arrangement. A transparent body model, cutaway or colored nerve layer explains selected relationships; it does not turn hidden structures into measured signals. Proposed tissue and interface comparisons need their own measurements over time. An unchanged replay file does not demonstrate an unchanged living connection.

\section{Reproducing the Studies}

Reproduction begins with source identity. A dataset name is useful for orientation, but exact files, versions and checksums identify the bytes that were processed. A checksum is a compact computed fingerprint of file contents. It helps a recipient detect a different or corrupted file. It does not certify that the experiment's labels are correct, that a source license permits every reuse, or that the model trained on those bytes is biologically valid.

The package should preserve original source links, acquisition receipts and the scope of each check. The full human recording-header files, whose filenames end in SET, were verified, but the acquired previews of the large signal files, whose filenames end in FDT, cover selected byte ranges rather than the complete large data files. Their locally hashed ranges must remain described as partial support. A whole-file claim would require a whole-file check. Derived replay arrays also need their own hashes and documented extraction rules.\cite{spinal_dataset2023,spinal_companion2023}

The computational recipe includes event sorting, half-open time bins, channel masks, unit conversions, coordinate transforms, numerical scaling calculated from training samples, random starting values, exact data divisions, boundary gaps, model selection and the saved model state kept fixed for evaluation. A half-open bin includes its starting boundary and excludes its ending boundary, so events on boundaries are handled consistently rather than counted twice. These details are part of the result because changing them can change the inputs and scores.

The body recipe includes the original asset, named forks, software revision, modeled constraints, controller settings and evaluation conventions. MuJoCo and OpenSim have separate environments in this project. A portable neural decoder does not require native OpenSim execution, while reproducing the native controller experiment does. The reader must not assume that one dependency list or a working dashboard recreates every mechanical result.\cite{mujoco2012,opensim2018,myosim2026,limblabmodel2026}

The evidence matrix maps the numerical claims in this document to inspectable project artifacts. Its claim groups retain the exact metric fields and interpretation limits. The long reader version explains those values; it does not generate a second experiment or alter the frozen scientific record. The source package also includes schemas and executable translation examples so that numerical predictions and their response metadata can be checked independently of the visual interface.

Previously invalid results remain withdrawn. The unsorted-event muscle-signal results, borrowed-rotation exports, stale control-function bindings and corrected initial-velocity protocol are not positive evidence to be revived by a later summary. A reproducible release needs the current valid states and the reasons earlier states were superseded. Software tests and archive checks then establish specified numerical integrity rather than a general certificate of scientific truth.

At this snapshot, the new project is a local reviewable software and manuscript effort. A verified public repository or deposition for its own code is not claimed. Responsible authorship, affiliations, an explicit code license, upstream attribution and an honest declaration of AI assistance remain release responsibilities. No submission to arXiv or a journal is implied by the formatting of the document. Public-source availability and local-project availability are different statements.

% R152 source proposal; provenance and exact prior text remain in the proposal manifest.
\section{From a reach to a testable prediction}
Start with a movement we can measure well. For the proposed snap showcase,
the arm begins extended with the palm down. The elbow and forearm swing and
turn; the wrist and hand turn upward during the snap, then move down to rest.
The model must account for that whole movement, including the preparation,
contact and return. Finger pairing, organism-specific anatomy and timing still
need a real-reference series; a human-looking animation cannot establish that
a particular monkey can perform the same gesture.

The cup reach supplies a simpler precision comparison: relate recorded signals
to measured movement before attempting the coordinated snap. Each task still
needs its own anatomy, contact, timing and independent comparison.

That comparison connects three interacting groups: task-linked brain
observations, named-muscle observations and a musculoskeletal simulation.
A learner would relate the observations to the mechanics; independent movement
and contact measurements would test its predictions. EEG records scalp
electrical activity and surface EMG records muscle electrical activity.
A separately recorded nerve or spinal signal needs its own declared channel.
The proposed apparatus uses two arm sleeves, scalp EEG and synchronized,
calibrated left and right tripod views. Vendor, channel, bandwidth and clock
qualification remain open.

Intent to action is the target. By understanding and measuring the intermediary,
we aim to make that mapping as close to one-to-one as possible within a declared
task. Specify what counts as the intended action, the participant and session,
state and context, permitted history, tolerated error and independent reference.
The test then asks whether observing the middle reduces ambiguity and improves
prediction or qualified control. It does not assume a unique hidden intention
or claim that the best possible mapping has already been achieved.

The proposed learner must explain the recorded brain and muscle activity,
where the hand moved and when. These are proposed comparisons: no paired trial,
trained model, data acquisition, fine-tuning, reinforcement-learning rollout or
accepted checkpoint is established for this EEG/EMG/movement loop. Its sources
and permissions, physical and computational behaviour, clinical interfaces,
reference comparisons and held-out tests still need independent assessment.

The completed LINK comparison supplies a warning about what would count as
success. Its 18 runs and 720 saved epochs tested a different question; both
model-admission gates and all five benefit claims failed. Neither those runs
nor a future successful prediction would show that a trainee received sensation
or produced an action. Recording, prediction, simulated movement, sensory
delivery and physical action each need evidence of their own. EEG and EMG
recorders alone cannot write sensation, and no sensory delivery or physical
actuation is established for this proposal.

\subsection{Form, function and action need a measured middle}
A structural map constrains what a system could do. A NeuroConnectome maps
specified neural connections; its organism, coverage, resolution and source
would have to be declared. The signals passing through those connections
still have to be observed or inferred, with the difference made explicit.
To predict the next action, we need the processes that turn structural
constraints into an outcome. The proposed programme follows form to function
to action, while intent to action names the task-level translation we want
to improve.

DNA to RNA and RNA to protein put a process between a starting constraint and
its downstream product. Regulatory context, RNA processing and cell state
change what that process produces; protein structure alone does not determine
the behaviour of a whole cell. The modelling jobs differ:
\href{https://www.nature.com/articles/s41586-025-10014-0}{AlphaGenome regulatory prediction}~\cite{alphagenomeRetained},
\href{https://www.nature.com/articles/s41592-024-02201-0}{scGPT cell-state modelling}~\cite{scgptRetained} and
\href{https://www.nature.com/articles/s41586-024-07487-w}{AlphaFold 3 structure prediction}~\cite{alphafold32024}
address different parts of that chain. AlphaTranscriptome names a desired
predictive capability here. Connecting the predictors would still leave system
dynamics and the environment to account for, with experimental outcomes needed
to test the combined prediction. These modelling examples do not establish an
integrated cell programmer. The comparison uses the reported modelling roles of
AlphaGenome and scGPT; their complete methods have not been independently
assessed here.
% Source-role pointers: alphagenome2026, scgpt, alphafold3; bounded owner checks.

In the biological pathway, Brain2Spine and Spine2Body name proposed links to
test; their composition is the Brain2Body aim. The spinal cord integrates
descending and sensory signals and contributes local coordination, reflex
regulation and adaptation. This is a different mechanism from RNA processing,
but it serves the same explanatory purpose: account for the context-dependent
process between an initial constraint and its outcome. The constrained example of
\href{https://pubmed.ncbi.nlm.nih.gov/38603479/}{Spinal sensorimotor adaptation}~\cite{lavaud2024}
illustrates a spinal process; supported observations would have to test our proposed pathway. Where they are absent, the internal
states remain inferred or simulated. Physical spinal work remains read-only.
That example establishes neither unrestricted skill transfer nor validation
of our proposed pathway. The comparison draws on the abstract and
relevant author-manuscript passages; equivalence to the final published edition
and the complete methods and statistics remain unverified here.

Thought2Action relates a declared intention to an independently measured action.
Action2World follows the action to its consequence in the task environment;
Thought2World names the proposed composition. Body2World drew attention to the
consequence outside the body; Body2Action focuses on body mechanics and a
specified task outcome, while Brain2Motion focuses on the movement pathway.
These working labels identify related scopes for a test. They do not establish
interchangeable or working controllers.

The engineering route connects neural activity to musculoskeletal simulation
and then to a programmable model, toward NeuroTranslation. The proposed
movement-focused learner is a large movement model. System 3 thinking names
the proposed human--AI supervisory partnership around it. The supervisory
model would compare task strategies and diagnose where neural or mechanical
predictions miss independent measurements, then propose training or calibration
comparisons for the signal mapping and virtual body. The task learner would
perform the declared prediction or control; the supervisor would help choose
what to test. Neither role replaces the biological spinal intermediary or a
physical interface. Final held-out outcomes must remain outside rewards and
model selection.
\par\noindent Shaw and Nave's SSRN preprint,
\emph{Thinking---Fast, Slow, and Artificial}, and its
\href{https://executiveeducation.wharton.upenn.edu/thought-leadership/wharton-at-work/2026/05/thinking-fast-slow-and-artificially/}{institutional discussion}
provide conceptual context for this proposed human--AI role.~\cite{shawNaveSystem3Retained}
The comparison is limited to that framing; the complete preprint and institutional
article have not been reviewed here. Neither establishes a third anatomical brain
system or validates this programme's proposed learner or supervisor. Expert footage may supply
movement and tactical examples, with rights and camera uncertainty accounted
for; it supplies no measured expert EEG, EMG, muscle force or stimulation plan.

Action changes the outcome. Its sensory consequences inform perception and
the next intention. Independently measured movement would anchor the action
part of that loop; it would not by itself decode a unique sensation or
perception. Known stimuli, measured feedback and appropriate recordings or
reports would have to test those further mappings. The return path must
therefore be measured rather than drawn into existence. The full
Interoceptomics aim includes internal-body sensing and connectivity; external
touch and proprioceptive tasks contribute evidence but do not by themselves
solve that wider problem. A prediction, assisted execution and retained skill
without assistance answer different questions. Each wider claim must earn
its own independent comparison.

\subsection{Context connects. Grounding checks}
Within a reach, context has a concrete job: relate brain activity to muscle
activity in this task and at this time. Attention offers a candidate computation for those relationships, following
\href{https://arxiv.org/html/1706.03762v7}{Vaswani et al. (2017)}~\cite{vaswani2017} and \href{https://aclanthology.org/P19-1656/}{Tsai et al. (2019)}~\cite{tsai2019}. The model compares a representation called a query with representations called keys, then uses the resulting weights to combine their values. Masks keep unavailable observations and future
samples out. The calculation weights information; it does not establish that
attention understands all context.
% Source registry IDs: context-native-01a12f62, context-attention2017-v7, context-mult2019, context-policy4

To see what those weights mean, follow one permitted observation through the
calculation. An encoder turns a stated signal segment into a list of numbers,
called a vector. Learned calculations produce three such descriptions: the
query describes what this part of the model is looking for, a key describes
an observation it can compare with that query, and a value carries information
that observation can contribute. These are roles in a calculation. Their
entries are learned features, not automatically electrical amplitudes, muscle
forces or identified nerve cells. The original channel, time and quantity
still need their own labels.

The proposed technical method writes the operation as

\[\operatorname{Attention}(Q,K,V)=\operatorname{softmax}\!\left(QK^{\mathsf{T}}/\sqrt{d_k}+M\right)V.\]

Read it from the inside out. Q holds the query vectors, K the key vectors and
V the value vectors. A query and key have the same number of entries, their
common width $d_k$. For one query--key comparison, multiply corresponding
entries and add those products. That sum is a dot product. The representations
are learned; multiplying and adding form the numerical matching score. It is not a measured biological
connection strength or proof that the matched observation caused the movement.

The superscript T on K means transpose. It rearranges rows into columns so
the matrix multiplication can calculate a score for each query against each
key. Transposing does not reverse the recording's time order or change which
electrode supplied a value. Dividing the scores by $\sqrt{d_k}$ means dividing
by the square root of the common vector width. With comparably sized components,
a wider dot product can have a larger scale; this division counteracts that
typical growth before scores become weights. It does not calibrate the
prediction's uncertainty or establish a physical unit for a learned feature.

M is the admitted mask. For each query it specifies which observations may
contribute. In the additive form shown, an allowed entry can leave a score
unchanged, while an excluded entry can be represented by negative infinity
so that it receives zero weight. The actual numerical implementation still
needs checking. An unavailable input is not an observed zero, and a query
with no admitted keys has no admitted weighted answer; the interface must
preserve that unavailability rather than manufacture a normal-looking result.

Softmax converts the permitted scores into relative weights. It applies the
exponential function to a score, giving a positive quantity, then divides
by the sum of those quantities for that query. The permitted weights therefore
sum to one when an admitted answer exists. A larger score receives more weight;
softmax does not simply choose the largest score and discard every other
permitted observation. The final multiplication by V multiplies each value
vector by its weight and adds the results. That combined representation still
needs a declared output calculation to predict a named target in its stated
units. A neat set of weights is not yet a correct reach.

The mask must enforce the claim all the way through the calculation. For an
EEG-only prediction, withheld muscle and movement targets cannot re-enter as
queries, keys, values or features computed upstream. Future observations must
also stay out of a prediction issued before they are available. A timestamp
alone cannot establish this: filtering, alignment or scaling could already
have used a later sample. Calling the attention mask causal cannot undo that
earlier leak. Shared cues and movement artefacts still need their own controls,
even when every matrix operation is correct.

No query, key, value weights, EEG encoder or gradient route have been selected
or executed for this proposed loop. The calculation supplies an inspectable
method. Independently measured movement must still supply its test.

A pattern can fit inside the model and still describe the wrong movement.
Grounding therefore means testing the declared prediction against qualified
observations under the same conditions. Independently measured movement would
check the outcome, while muscle activity would add an observation that
constrains the brain model. Even a matching reach cannot identify a unique
private experience. The proposed paired trial would have to show which of
these observations adds useful information.

\subsection{Muscle activity tests the brain model}
Record brain and named-muscle activity during the same reach. In the proposed
paired trial, EEG would sample scalp electrical activity and surface EMG would
sample activity associated with the named muscles. Paired training could use
the muscle observation to constrain their relationship. Then withhold the
muscle target and movement reference from the EEG-only predictor: can the
permitted brain observations predict a specified muscle-activity observable
or movement?

Pairing and calibration remain proposed. Shared cues and movement artefacts
could give that predictor a shortcut, so the comparison needs controls for
both. The muscle observations would cover only the named muscles, and EMG
supplies neither afferent nerve activity nor direct force. Sensory claims
still need known stimuli and measured feedback. Several neural states could
remain compatible with even a well-matched pair; the pair could narrow a
task-specific explanation without revealing the complete hidden brain.
Checking the movement itself requires the next independent observation.
% Source registry IDs: context-native-01a12f62, context-policy4, overhaul-brantley2018
\par\noindent The candidate reference at \url{https://www.nature.com/articles/sdata2018133} has not yet been checked for exact support or reuse rights.

\subsection{Measured motion tests the virtual body}
Measure the same hand and arm from the aligned left and right tripod views.
The proposed motion-capture relay-composite would align the two point maps
in a common clock and coordinate frame, then join them into one composite
while retaining each point's origin. Before that joined map can check a
prediction, its calibration, alignment and capture precision would have to
be independently verified.

The movement map would constrain the external movement model: surface motion
and pose in the task's physical frame. To explain that motion mechanically,
the internal musculoskeletal model would describe specified bones, joints
and soft tissues. Organism, developmental stage and task would determine the
appropriate anatomical sources. X-ray projections could inform bone and
joint relationships; appropriate MRI references could inform specified soft
tissues. Registration, rights and source-specific limits would need independent
checks. No scan set is selected for this proposed comparison. Anatomy still
leaves activation, forces and neural function to test.

With a qualified composite, the musculoskeletal-simulation physics-engine's
joints, muscles, contacts and physical limits would face a measured comparison.
The simulation methods in \href{https://journals.plos.org/ploscompbiol/article?id=10.1371/journal.pcbi.1006223}{Seth et al. (2018)}~\cite{opensim2018} and
\href{https://myosuite.readthedocs.io/en/latest/}{MyoSuite documentation}~\cite{myosuiteDocsUndated} would support this comparison: fit the model on training
data, then compare its predicted reach with held-out measured movement. A force claim would require an additional force
measurement. The body's form limits possible motion; measured function would
test the mechanics and show where they need refinement.
% Source registry IDs: context-native-01a12f62, context-component-names-01a12f81, context-policy4, context-opensim2018, context-myosuite-docs, context-precision-01a12f77


A lifelike hand could expose a missed contact or a late reach. The comparison
still needs independent capture alignment, precision, mechanics and biological
fidelity; none has been demonstrated for this proposal. Captured geometry
cannot establish hidden neural state, internal forces or unique subjective
experience. It would check the declared movement prediction within those limits.

\subsection{Two pairs. One conditional learner}
The learner would have two connected jobs: relate EEG observations to muscle
activity, and relate measured movement to predicted mechanics. Large movement
model names the proposed family; deep neural network describes a possible
implementation. Large refers to intended scope, not a demonstrated parameter
count. Encoders would accept only observations permitted for the named test.
Two branches could use the attention and multimodal computations of
\href{https://arxiv.org/html/1706.03762v7}{Vaswani et al. (2017)}~\cite{vaswani2017} and \href{https://aclanthology.org/P19-1656/}{Tsai et al. (2019)}~\cite{tsai2019} to summarize those relationships
in a shared task representation, with separate output parts for each direction's targets.
% Source registry IDs: context-native-01a12f62, context-policy4, context-attention2017-v7, context-mult2019


Signal features belong to branches, and branches belong to the learner.
That hierarchy describes the computation, not two anatomical thinking systems.
A common clock, task and physical frame would connect the branches; independent
measurements would check their errors. Form constrains the prediction, and
checked function would guide refinement. Several hidden states could still
fit the same reach, so neither a shared representation nor conditional
predictions would establish an exact inverse nervous system. The architecture
has not been selected, implemented, trained or independently accepted.
To choose it, we would have to compare it with simpler baselines on
appropriately reserved observations.

\subsection{Learn the movement, then test what transfers}

A recording of a reach already contains a possible lesson for a computer: use the earlier movement to predict where the hand goes next. It need not contain EEG to pose that question. A separate proposed comparison would ask whether learning from one person's grasp-and-lift movement helps predict another person's movement. The WAY-EEG-GAL description identifies continuous recordings of lifting series as well as separate windowed trials and lift-summary tables.\citep{luciw2014way} The continuous movement could supply both the history and the later checking observation. It would not require matching AllLifts Run or trial rows, and it would create no measured neural pair between the two people. The original study of five complete weight series remains a different design.

To transfer anything useful, we must say what corresponds. Homologous structures share an evolutionary anatomical origin. Mechanically analogous structures can perform a comparable job without that shared origin: a robot gripper's opening and a human thumb--finger opening might both describe how an object is approached. That example would require a separate robotic study; the proposed two-person study would not test it. Each body keeps its own joints and possible directions of movement, positive-axis conventions, proportions, muscle--tendon roles, loads, inertia, contact and timing. Two similar-looking limbs need not generate the same forces or move at the same rate. Calculating a corresponding movement for another body is retargeting; the calculation remains derived, alongside separately identified tracking observations, estimates, predictions and simulations.

The proposed description of movement contains thirteen quantities. Three give the index sensor's position relative to the wrist, three do the same for the thumb and three for the object. One describes the distance between the two digit sensors; the final three describe wrist velocity using only information already available. Thirteen extra mask values say which features are actually available. Each person needs checked position units, sensor placement and a wrist orientation that turns positions into the intended frame without reflecting the body. After the checked conversion to metres, the middle, or median, wrist-to-index-sensor distance from that person's assigned training support supplies a positive length scale for position and velocity. This is a sensor-distance scale, not a measured bone length; development observations do not set it. Velocity is also multiplied by the same 0.1-second interval in both people. The thirteen prediction features then have no physical units. To recover velocity errors in metres per second, multiply by that person's length scale and divide by 0.1 seconds. Missing features contribute no target error; writing a masked zero does not make a missing observation measured. These features do not supply all finger joints, individual muscle forces or the person's inertia.

The model would use fifty movement frames at 100 observations per second to predict 100 milliseconds beyond the latest filtered input actually available. The source rate must first be checked as 500 observations per second. A proposed filter combining 51 present and earlier samples would precede keeping every fifth sample, producing 100 observations per second. Its linear-phase delay is 50 milliseconds. That delay stays in the account of when the information became available; moving a timestamp backward cannot make the predictor respond earlier. Filter coefficients, the frequencies it preserves, timestamp conventions, gaps and handling at recording edges still need checking. Only the named movement columns may enter. The source's kinematics matrix also contains force-plate signals, which cannot quietly become movement features.

Knowing which input was read is different from knowing whether two inputs are independent. Trustworthy published participant attribution, an unchanged exact archive entry and a checked capture can identify the participant record and recorded input without requiring one particular spelling of an ID inside the file. They do not authenticate a biological person or prove that two files came from different original acquisitions. A longer recording and its trimmed copy could have different names and file fingerprints while sharing the same movement. Shared samples, trims, derived versions and overlaps must therefore remain together as one acquisition group before filtering, cutting windows, scaling or assigning roles. The complete public non-ST roster and these relationships still need establishing. All 294 rows of the earlier AllLifts table remain unavailable for training, evaluation and assigning trials or acquisition groups. ST, competition storage and reserved LINK recordings remain closed.

Qualified P1 groups would provide the earlier movement learning. P2 groups would be divided between target training and development checking. At least two independently established P2 acquisition groups are needed. If there are $G\geq2$, a fixed hash calculation ranks their established original-parent identifiers; the first $\lfloor3G/5\rfloor$ become training groups and the rest become development groups. The floor brackets mean rounding down, so the nominal 60/40 division is explicit even when the number of groups is small. The exact identifiers, how they are written as bytes and the complete roster must be settled before ranking or seeing outcomes. The order is not chronology and need not balance lifting conditions. If uncertain ancestry forces all P2 recordings into one group, the comparison stops. Cutting that group into windows cannot create independent training and checking acquisitions. No actual group or movement segment has been assigned by this proposed identity method.

The three training arms keep the target task and model the same. B0 first learns from correctly mapped P2 training movement. B1 and B2 first learn from exactly the same P1 windows in the same order. B2 uses the proposed correct correspondence. B1 deliberately uses the thumb sensor for the index role, the object sensor for the thumb role and the index sensor for the object role, leaving the wrist anchor fixed. Both use the same correct source calibration and scale before this change, then the same correct target adapter and ordered target-training batches. Sampling gives each original acquisition group balanced treatment and records repeated as well as unique windows. The learner might undo the false correspondence; changing roles can also change which features are available. Equal numbers of updates do not give the arms equal diversity of source movements. Those limits must stay in the explanation of any difference, rather than prompting a more convenient control after results are known.

The proposed network is a single gated recurrent unit, or GRU, layer: a calculation that updates a numerical memory as it reads the sequence. It has a memory width of 32, accepts the thirteen movement features and thirteen masks as 26 inputs, and predicts thirteen outputs. That software memory is not measured biological memory. Each arm uses the same initialization for a given fixed random seed, with seeds 17, 31 and 73. Every arm/seed receives 200 earlier-learning updates and 100 target-training updates, using batches of 16 examples. Adam, the proposed weight-update rule, has learning rate 0.001, beta settings 0.9 and 0.999, epsilon $10^{-8}$ and zero weight decay. There is no clipping, changing learning-rate schedule or early stopping; the final update supplies the checkpoint. During training, squared prediction errors are averaged over observed target features; a batch with no observed targets must stop. Across three arms and three seeds this would be 2,700 updates. None has been performed for this comparison. No external checkpoint is selected, and the numerical updates, filter, precision and actual computing limits still need independent checks.

For each development acquisition group and seed, the proposed score adds the squared errors for the observed target features, divides by how many such targets were observed, then takes the square root. It then averages those group scores equally, so a long recording does not automatically outweigh a short one. A comparison with no observed targets has no valid denominator and must stop. The report must give feature counts, missing observations, errors converted back to physical units, uncertainty and a persistence forecast---continuing the latest available movement value---using the same information endpoint. Every worse group or feature remains visible. B2 would need lower normalized error than both B0 and B1 in each fixed seed, with no applicable original requirement failing, to justify another investigation. Two people and development checking do not establish performance across a population, untouched final confirmation or biological teaching.

Larger movement collections could supply more examples after their exact files, permissions, relevant movements and original-recording overlap are checked. \href{https://www.addbiomechanics.org/download_data.html}{AddBiomechanics} distinguishes recorded markers and ground reactions from fitted joint movements and torques estimated by inverse dynamics---calculating forces or torques consistent with motion and a body model. Those estimates are not new direct force measurements. The selected \href{https://amass.is.tue.mpg.de/license.html}{AMASS} asset would need its research-use, non-commercial and redistribution terms checked. The same original recording in two collections stays one source, not independent training and checking data. LaBraM's earlier training inventory flags the WAY source family, but exact checkpoint, participant and trial exposure remains unknown. A later decision to reserve a recording cannot remove its possible influence from already trained weights. No such checkpoint is selected for this comparison.

Scans have a more specific job than making the body look complete. Before adding X-ray/CT or MRI, check which anatomical sources the chosen musculoskeletal model already uses. A new scan is useful when it fills a stated gap or provides a genuinely independent check. Its person or species, scale, registration, permissions and uncertainty still need checking. Registration establishes the correspondence between that image and the body being studied. Another person's scan remains an external reference; a generic body model does not secretly measure this participant's geometry, muscle activation or forces. No scan set is selected here.

The actual eligible recordings, acquisition groups, units, frames, clocks and calibration are still unresolved for this continuous comparison. It reports no new fit or successful transfer. A movement description that proved useful could later be tested alongside genuinely paired neural and body observations, using only the inputs allowed for that prediction. The transferred or calculated target would not become a simultaneous measurement. The useful next answer is where the movement forecast improves, where it worsens and which measurement or assumption explains the difference. Neural equivalence, causal physiology, delivered sensation and a skill that remains after guidance ends would each need a further experiment.

\subsection{Pretrain. Pair. Verify. Refine}
Those comparisons start with the observations we may use. Data rights,
calibration and held-out subjects, sessions and tasks determine which can
enter the study; resolve them before processing data or choosing a model.
Open-source code does not grant blanket permission to use datasets,
checkpoints or clinical interfaces. Qualified data could first train
representations for each recording type. LaBraM, described by \href{https://arxiv.org/html/2405.18765v1}{Jiang, Zhao and Lu (2024)}~\cite{labram2024}, is one EEG encoder
candidate, not a complete brain-body writer. Synchronized paired observations could
then train the declared mapping from permitted inputs.

That first split has to survive the model's past. Before fitting a scaler,
training a tokenizer or selecting an encoder, identify the exact proposed
checkpoint and its pretraining corpus: versions, original recordings,
participants, sessions, tasks and every derivative used. The same recording
can appear in several collections; overlapping windows can appear under
different filenames. Keep those copies and derivatives with their original
recording and participant when defining the applicable splits. A new archive
name does not make the person or trial new.

The current overlap audit flags the WAY-EEG-GAL source family in LaBraM's
pretraining inventory. The exact checkpoint, participant and trial exposure
remains unknown. No LaBraM checkpoint has been selected for the separate
proposed hand/arm transfer comparison, and its target development observations
are development data throughout. A model tested there cannot be described as
independently confirmed on untouched final recordings. The audit needs the
actual source and checkpoint identities before that stronger claim can be
assessed. Unknown overlap is a limit on an independence claim; it is not a
blanket prohibition on useful pretraining or transfer.

For components fitted in this study, means, scales, tokenizer parameters,
codebooks and encoder updates must use only their assigned training support.
Development observations can make the predeclared selections; final held-out
observations must not set the preprocessing, representation, mapping, stopping
point or reward. A separately admitted external pretraining corpus can be
useful, but its exposure remains part of the model's history. Removing a record
from the later fine-tuning set cannot remove its influence from weights already
trained on it. Freezing a previously exposed model does not turn that record
back into independent final confirmation.

LaBraM's candidate sequence separates three learning problems. First, a neural
tokenizer would learn a way to describe specified signal segments. A tokenizer
is the calculation that turns those segments into model tokens. In this route,
a learned set of numerical patterns, called a codebook, supplies discrete
code identities for the segments. A code identifies a model pattern; it is not
an EEG voltage, a neuron identity, a muscle command or a word whose meaning is
already known. The codebook's learned numerical entries and the code labels
that select them are different things. Their signal support and tokenizer
version have to stay with the training record. The segment definition,
tokenizer objective and numerical implementation remain choices to qualify;
none has been selected for our paired learner.

Second, masked neural-code prediction would train an encoder using those code
targets. Hide selected signal segments from the predictor, give it the admitted
remaining context, and ask it to predict the code identities assigned to the
hidden segments by the specified tokenizer. Its output assigns scores or
probabilities to possible codes. The training loss penalizes a poor prediction
of the saved target code at the hidden positions. Computing that target from
the original training segment is legitimate supervision; leaking the hidden
segment, its target code or a feature calculated from it into the predictor
would defeat the exercise. Predicting an easy-to-leak answer would teach us
little about the representation.

Hiding a segment for learning and excluding an observation by permission serve
different purposes. A learning mask makes a stated prediction exercise; an
input-availability or causal mask limits what the named route may know.
Pretraining must declare its permitted context. An offline masked-segment
exercise does not by itself establish that an online prediction uses only
available past and present inputs. That later claim needs the entire input
route checked under its deadline, including the preprocessing.

Third, downstream fine-tuning would adapt the representation to the declared
task. Actual paired observations would supply the named EEG, muscle or movement
targets for the permitted direction. The outputs would then face their own
units, uncertainty, baselines and independent measurements. A useful code
prediction score cannot substitute for a movement error, a force measurement
or retained biological learning. Likewise, the previously reported neural
next-bin training is its own objective; it is not evidence that this staged
tokenizer route has already run.

Movement-only pretraining and unpaired representation learning remain useful
candidate routes. Transfer may use homologous structures with shared anatomical
origin or mechanically analogous structures with comparable jobs, after their
actual mechanics, scale, timing, contact and missing features are qualified.
Such training could supply a movement representation before every source has
simultaneous neural observations. Retargeted or learned movements must keep
their derived identity; they do not become measured neural--muscle--movement
pairs. The proposed transfer and the paired prediction ask complementary
questions, each with its own matched controls and independent target. Neither
is a new fit or a successful transfer reported here.

Paired prediction would give later refinement a baseline to beat. Before
optional reinforcement learning, post-training or task-specific fine-tuning,
the physical predictions and reward need checking against independent
outcomes; a reward the computer can calculate supplies no biological proof
by itself. The offline reinforcement-learning method of \href{https://arxiv.org/abs/2006.04779v3}{Kumar et al. (2020)}~\cite{cql2020} is a
candidate only where qualified actions, transitions and rewards support it: labels alone cannot tell us what an untried action
would do. Backpropagation can update differentiable neural components, but
it cannot automatically pass through every physics engine. The simulation
methods in \href{https://journals.plos.org/ploscompbiol/article?id=10.1371/journal.pcbi.1006223}{Seth et al. (2018)}~\cite{opensim2018} and \href{https://myosuite.readthedocs.io/en/latest/}{MyoSuite documentation}~\cite{myosuiteDocsUndated} do not establish
universal differentiability for this proposed route. Simpler baselines and frozen final
tests remain through every phase.
% Source registry IDs: context-native-01a12f62, context-policy4, context-labram2405, context-cql2020-abstract, context-opensim2018, context-myosuite-docs


This proposed paired training sequence has not been carried out.

\subsection{Respond quickly. Improve carefully}
A frozen accepted model would provide the proposed response: predict from
permitted past and present observations, then prepare encoded outputs for a
qualified interface. Inspecting the errors would be a separate job. That path
would calibrate and propose a candidate update, to be independently evaluated
with a restorable checkpoint before it could replace the running model.
Updating an estimate of the current state therefore differs from changing
learned weights; examining an error would not silently change the predictor.

This fast-and-slow comparison is a Kahneman-like analogy for the proposed division of work. It does not identify two biological systems or assign inherent thinking speeds to EEG and spatial data. Real-time performance would have to be measured from acquisition through filtering, alignment, inference, output, feedback and rendering, including delay tails and jitter. No accepted model, real-time demonstration or running system that adapts under these rules is established here. The loop has no measured end-to-end latency or accuracy guarantee.

That timing comparison must also count missed deadlines. An observation that
arrives too late cannot repair the response already due. Test the complete
chain while it coexists with the declared healthy workloads, and report the
conditions as well as delay tails and jitter. Timing one model calculation
alone cannot establish that acquisition, prediction, output and feedback meet
the task's deadline together. These remain proposed measurements, not a
real-time result from this loop.
% Source registry IDs: context-native-01a12f62, context-policy4

\subsection{Write sensation. Enact action}
What reaches the trainee when the computer has made its prediction? A sensory
transducer would have to deliver a sensory output; a controller and actuator
would have to produce action in a virtual, robotic or later qualified biological
body. Input and output describe the computer boundary. Afferent and efferent
describe physiological direction relative to the central nervous system.

A display or haptic interface supplies sensory feedback, which differs from direct neural writing. Direct cortical writing could have a sensory role without being a peripheral afferent signal, just as a robot command is not an anatomical efferent nerve. EEG and EMG recorders do not write, and predicting an EEG-like pattern does not deliver sensation. Each proposed output needs an independently qualified interface, permission for its use and a measurement of what it delivers. No such physical interface, stimulation, acquisition or actuation is established by this proposal.

Comparing a recorder with a sensory or action device requires evidence about what each measures or delivers. The clinical source could not be read, so it supplies no methods evidence here. Implanted precedents cannot establish that EEG or EMG recorders are equivalent to devices that deliver sensation or produce an action. Those physical additions remain unpriced and outside the acquisition and virtual-modelling scope.
% Source registry IDs: context-write-native-01a12f66, context-native-01a12f62, context-policy4

\subsection{Close the loop with evidence}
Return to the cup with two independent checks: withheld body observations
for the brain-to-body prediction, and the recorded neural target for the
conditional estimate from the body. Agreement between the models or a high
reward could still accompany a wrong answer. Neither establishes a unique
inverse brain or someone's private experience. A later sensory or action
interface would need a further measurement of what it delivered, under which
conditions, with what error and delay.

Approximation gets us towards plausibility and virtualization gets us to
precision. Independent measurements, repeatability, causal checks and stated
error bounds must earn that precision; the proposal supplies neither perfect
biology nor an automatic result. The original held-out and real-reference
controls would test where each prediction works and where it fails. The
retained-learning controls ask a further question: what remains when useful
guidance stops, with any necessary movement bridge kept constant?

Physical sensory writers, transducers and actuators remain unpriced additions
outside the acquisition and virtual-modelling costs. Those costs cannot tell us
the price of delivering sensation or action. No completed loop, training or
physical-writing result is claimed here. \emph{Follow the action, then test the
prediction} develops the tasks that would challenge the mechanism;
\emph{Measure both ends. Test the translation.} gives the detailed measurement
and validation protocol. Comparable observations, predictions and errors would make those tests possible.
% Source registry IDs: context-native-01a12f62, context-write-native-01a12f66, context-precision-01a12f77, context-policy4, context-spec15, context-glossary13, b2b-circular-checklist-original

% Proposed full-depth accessible insertion; not a summary of the technical text.
% Inherit the selected master's Noto Sans/Fira Math setup.
\section{More data must survive a harder test}

Give a model another thousand examples of the same reach and it may get better at that reach. Give it examples from another person, another task or another kind of body and you change what it has to learn. That is the attraction of a Large Movement Model: the relationships learned from a varied collection might help it cope with a movement it has never encountered. It is also a question we can get wrong. A broader collection could help, add nothing or teach a correspondence that makes a particular prediction worse.

The proposal therefore makes two bets that need separate tests. One is about amount: does more qualified training data improve the chosen prediction? The other is about diversity: does a wider range of qualified examples help when their total amount is held comparable? Neither bet has been demonstrated here. No body-wide world model has been trained or validated by this proposal.

\subsection{Four source groups, distinct observations}

The proposed training collection has at least four groups. One supplies neural observations. A second supplies muscle and other myographic observations: records of muscle activity whose actual measurement type must be stated. A third supplies measured movement and displacement. A fourth supplies examples generated by physics-based motion or musculoskeletal simulation. Each would teach the model something different. Listing them does not mean that four complete collections have been acquired, or that nothing else could be useful.

The intended neural range includes suitable brain, spinal, peripheral and relevant visceral recordings. A recording's site, sensor and role in the task determine what it can tell us. The current focus is motor--sensory. The programme uses Myoception for sensing muscle-driven movement and the feedback that comes with it, including body-position and movement feedback, touch and other task-relevant sensations. This proposed programme use of the name establishes neither a new external scientific standard nor a simultaneous record of the whole body. The intended range remains broader than the available paired examples.

Touch, body-position and movement feedback (proprioception), contact, force and context cannot be treated as interchangeable labels. A known tactile stimulus says what the experiment delivered. A contact sensor reports the event it was built to detect. A calibrated force instrument measures a specified push or pull. Movement measurements can help estimate where a limb is without directly recording the sensory nerve signals or how its position feels to the person. EMG records muscle electrical activity; it does not directly record peripheral nerve or sensory activity. A stimulus, a report of sensation, a recorded neural feature and a model's estimate therefore keep their different meanings even when they concern the same action.

The simulation group is a training source in its own right. A qualified virtual body could supply examples of motion, changes between states, calculated muscle--tendon and contact forces, or movement patterns used as a starting point for learning. Each example would retain the simulator and body-model version, parameter settings, numerical solver, controller and its original sources. A controller that requests a joint position differs from one that drives muscles. Calling both muscle control would conceal information about how the movement was produced.

If a measured motion is transferred to another body model, keep the original recording and the transformation with it. If the computer then calculates internal states or forces, those values remain simulated. A calculated contact force does not become an independently measured one because a measured trajectory helped produce it.

The generator first needs calibration against suitable measured support, with uncertainty and data roles kept explicit. Its adjustable settings would use training support, leaving independent tests untouched. A proposed comparison would keep a fixed amount of real data in both arms and add simulation examples in one arm. It would match the permitted information, computing work and original decision rules, then check the contribution against held-out real observations. Record how often each arm sees real and simulated examples and how many updates it makes. Keeping the same real recordings does not ensure that the models see them equally often. Count simulated examples and their variety separately from real acquisitions. Real test outcomes must never be used to tune the generator, calibrate it, select rewards or choose the model that will face those same tests.

Simulation could teach a useful relationship, or pass on its own wrong contact rules, dynamics or body correspondence. Those failures must remain visible. Generated examples cannot independently verify the biological assumptions that generated them. The methods, implementation and use conditions of any proposed generator still need checking; no frontier implementation or dataset is selected here.

These four source groups do not replace the programme's five component roles. Brain signals, muscle signals, measured movement, a virtual body and a learning model describe what each part does in the system. The source groups describe what can supply training examples. The virtual body can supply examples and constraints, while the learning model can use several groups. Other useful information can be added if its meaning and measurements can be checked.

\subsection{Count what the model actually receives}

A recording can be divided into many short windows, and the same window can be shown to a model many times. Those operations increase the number of training presentations. They do not create another person or another independent recording. For every input type, the study would report the original acquisition groups, the people or animals, recording time, trials, available target values, distinct windows and repeated presentations. An acquisition group keeps together material that came from the same underlying recording, including copies, overlapping extracts and differently named versions.

Diversity needs a description too. Different people, recording days, movements, body regions, sensors, apparatus and environments expose different changes. A dataset containing one new sensor is not equivalent to a dataset containing twenty new people. One more dataset is therefore an inadequate measure of breadth. The proposed study would identify which aspect changes and which aspects remain comparable, including the quality and missingness of the measurements.

To test amount, we would increase eligible training material from a declared source distribution. To test diversity, we would hold a stated amount comparable while changing selected aspects of the collection. A design that crosses both changes could ask whether their effects depend on each other. The model design, permitted inputs, prediction target, preprocessing and selection rules would remain fixed. The experiment would report how often examples were reused, the number of updates made to the model and the computing work actually consumed. Equal numbers of updates need not cost the same or expose the model to the same information.

Some comparisons will resist a clean match. Adding a sensor type may also change the task, participant population or recording quality. In that case we would report the combined change and its uncertainty. Calling it diversity alone would assign the result a cause the comparison has not separated.

These are proposals for a later, separately specified study. Before running it, the exact recordings, training budgets, random starting conditions, controls and pass/fail rules would have to be fixed. The proposal does not change the already adopted continuous WAY-EEG-GAL development comparison, its participant roles or its prescribed updates. It does not open reserved LINK recordings, ST or other protected tests. Calibration and numerical scaling would use training support only. A recording already seen by the model cannot be made unseen by giving its windows new filenames.

\subsection{Ask where the improvement holds}

A new recording day from a familiar person tests something different from a new person. A new person tests something different from a new movement, a new apparatus or another species. Generalization means carrying a learned relationship into a stated new setting, so each of those changes needs its own result. The study would keep the relevant original recordings together before making windows or choosing settings, and would audit any earlier pretraining exposure.

Each proposed transfer would face a model trained on target data alone, the applicable deliberately incorrect correspondence and appropriate simpler predictors. The original tests for false associations, shuffled pairings, shared cues, recording artefacts and access to future information would remain. Comparing a source-trained model only with an untrained model would not show whether the source data adds useful information beyond a simpler alternative. Nor could a better average erase harm to a recording or output protected by an original no-harm rule. Unhelpful transfer, plateaus and outright failures belong in the result.

\subsection{What a power law would mean}

One possible pattern is that error falls as training amount rises, with each further increase buying a smaller improvement. A proposed mathematical description is
\[
E(N)=E_{\infty}+A\left(N/N_0\right)^{-\alpha}.
\]
The letters name quantities that the study would define before fitting the curve. $N$ is the chosen training amount and must be greater than zero. $N_0$ is a fixed positive reference amount in the same units, so dividing one by the other leaves a pure number. $E(N)$ is a nonnegative error on the specified held-out population. For this proposed decreasing family, the floor $E_{\infty}$ must be zero or positive, the amount-dependent size $A$ must be zero or positive, and the exponent $\alpha$ must be positive. The floor and $A$ have the error's units; the exponent has none. If $A$ is positive, this curve falls as training amount increases. If $A$ is zero, it stays flat however large the positive exponent. The exponent alone therefore cannot establish the direction. If the error measure has an upper limit, the predicted values must stay inside it over the declared fitting and prediction range as well. A fitted floor would describe this candidate curve; it would not establish an irreducible limit of biology.

There is no estimated exponent or demonstrated law here. Choosing a decreasing candidate does not require the measurements to improve. A flat line, a plateau, rising error, a reversal or a different pattern for each domain may describe them better. Harm from added data must be reported, even when the proposed curve cannot represent it. The study would choose its fitting range, alternative curve shapes, uncertainty method and checks on withheld data-amount levels in advance. It would ask whether the fitted relationship predicts levels left out of fitting, rather than merely drawing a pleasing line through the points it has already seen. Overlapping windows from one recording do not provide the independent repetitions needed to make uncertainty vanish.

The measurement being fitted also stays fixed. An error in rescaled muscle-signal amplitude, an error in camera pixels and an error in physical joint angle describe different things. Combining them into one percentage called biological accuracy would need a defensible, independently checked scoring rule. The proposal instead retains the original units and reports each domain, the number of observations, absent values and uncertainty, alongside any justified common summary. More training data, a larger computational model and averaging over more outputs remain different changes.

\subsection{Two useful sources do not make one shared event}

FALCON M1-A illustrates a possible neural--muscle link. The fixed release describes simultaneous motor-cortex activity and processed electrical activity recorded with electrodes within the muscles of one rhesus macaque. That acquisition is called intramuscular EMG; it differs from a surface recording taken over the skin. Processing the signal is a separate step from where it was recorded. Its three split labels do not mean three animals. The reviewed metadata does not establish a released synchronized motion-capture trajectory. Some processing of the muscle targets uses information on both sides of a time point, including later information. A model scored against those targets would therefore need a separate test before being described as purely causal in real time. Its small minival subset is drawn from calibration data to check submission format. Because those examples are already part of calibration, they are not an independent validation population. \href{https://doi.org/10.48324/dandi.000941/0.241029.1405}{FALCON M1-A release}; \href{https://papers.nips.cc/paper_files/paper/2024/file/8c2e6bb15be1894b8fb4e0f9bcad1739-Paper-Datasets_and_Benchmarks_Track.pdf}{FALCON methods}.

AJILE12 supplies a different candidate link. It connects recordings inside the skull with upper-body positions estimated from video in twelve clinical participants. The positions are two-dimensional image coordinates, measured in pixels, rather than instrumented three-dimensional motion capture. The reviewed release does not establish paired muscle electrical recordings. Its pose estimator is adapted to each participant using that participant's labelled video. Whether a particular day chosen for a later test contributed to training that estimator remains unverified, as does a numerical bound on alignment error. Even when a neural model has never seen a participant or day, the method used to create the movement labels may already have been adapted to that participant. That dependence belongs in the account of what was held out. \href{https://doi.org/10.1038/s41597-022-01280-y}{Peterson et al., 2022}; \href{https://doi.org/10.48324/dandi.000055/0.220127.0436}{AJILE12 release}.

Those sources could contribute different parts of a larger research programme. Their recordings do not share an event, a body or a clock. Combining a macaque's brain--muscle example with a person's brain--pose example would not produce an observed brain--muscle--movement triple. Separate, qualified pretraining could learn useful patterns from each source, followed by a declared transfer whose predictions are checked against actual target measurements. The assumed correspondence would still need testing. The release descriptions identify possible source relationships. They do not establish that the recordings have been acquired or that a model has learned from them. Their exact use conditions, file integrity, units, clocks, earlier exposure and the independence of the chosen partitions still need checking before training.

\section{A world model must account for the body it moves}

Watch a hand close around an object. From outside, the fingers change position, touch begins and the object may move. The body model also needs an inside: muscles and tendons exert forces, joints constrain motion and the load resists it. The proposed world model would relate these events, so that a forecast of movement can be checked against what the declared mechanics and observations support.

Its initial interior consists of specified muscles, tendons, joints and mechanical states. Its exterior consists of supported movement, objects, contacts and environmental conditions. These are deliberately stated parts of a body and a task. Calling the model a world model does not establish complete organs, a body-wide neural recording or access to private sensation. An absent measurement stays absent, however convincing the picture becomes.

The existing component names describe proposed jobs. NeuroTranslate would map declared observations to named targets. MoCap ReComp would join source-labelled movement observations in time and a common physical frame; a collection of estimated video poses does not become that proposed instrument by changing its label. MyoSim Replay would calculate the declared body's mechanical response. Their proposed combination is the Large Movement Model. A neural composite would bring identified recording types and their source-dependent representations into a model. It would not establish that a cell in one person matches a cell in another, or build an observed whole nervous system from unrelated parts.

\subsection{Give every internal value a meaning}

A proposed mechanical state could include joint configuration $q$, generalized velocity $v$, muscle-model activation $a$ and environmental state $c$. Configuration describes how the model is arranged; velocity describes how it is changing. A rotational joint may use radians and radians per second, while a sliding coordinate needs length and length per second. The actual simulator determines which coordinates it uses. Its velocity coordinates need not be identical to ordinary derivatives of every stored position coordinate, so that conversion must be stated.

The model could use permitted earlier observations, task context and a declared control input to predict a range of possible later states. Predicting a range exposes uncertainty that a single animation can hide. During use, the predictor may receive only the inputs allowed for that route. A future measurement used as a training target does not become available before it occurs.

The instruments observe different things. A neural recording supplies its named electrical or imaging signal. EMG supplies electrical activity associated with muscle; a processed EMG amplitude is not automatically a measured activation state or force. Tracking supplies measured or estimated positions, with the camera geometry, calibration and uncertainty attached. A force calculated inside the virtual body remains a calculated force until an independent reference tests it. An availability mask tells the model which values are absent. Filling the numerical array with zero for computation must never be described as having measured zero.

The mechanical accounting can be expressed as
\[
M(q)\dot v+h(q,v)
 =R(q)F_{\mathrm{mt}}+J_c(q)^{\mathsf T}f_c+\tau_{\mathrm{ext}}.
\]
The left side describes the movement changes and the model's declared velocity and gravitational effects. $M$ represents how its mass is distributed and resists acceleration; $\dot v$ describes change in generalized velocity. The term $h$ collects the specified velocity-related and gravitational contributions. On the right, $F_{\mathrm{mt}}$ contains muscle--tendon forces, and $R$ maps their pulls into forces or turning effects on the model's coordinates. Contact forces $f_c$ act through a mapping $J_c^{\mathsf T}$. The final term $\tau_{\mathrm{ext}}$ accounts for other declared external contributions.

Each side must use compatible coordinates and force or turning-effect units for each component. The equation is proposed bookkeeping for a chosen articulated-body simulator. It is not a discovery of all the forces inside a living person. Its geometry, mass distribution, muscle parameters, loads and contact rules need evidence. Established OpenSim methods show how musculoskeletal models can be specified; they do not validate the proposed combination here. \href{https://journals.plos.org/ploscompbiol/article?id=10.1371/journal.pcbi.1006223}{Seth et al., 2018}.

\subsection{Make the visible movement answer to an independent measure}

A predicted path and its reference must describe the same interval, body frame and units. Internal force claims need suitable force measurements as well: different combinations of muscle forces can be compatible with a similar visible path. The tests would examine contacts, allowed joint motion, energy accounting, gravity, friction and sensitivity to numerical solver settings within the claimed domain. A score can improve because the model exploits a numerical correction or has been given future targets. That result would identify a failure in the computation or test, rather than improved forecasting of the intended movement.

The proposed comparison would also remove or replace inputs and constraints one at a time, using matched conditions, to ask what the neural representation, muscle observations, movement prior and physics each contribute. Independent measured outcomes would decide whether those changes help. Later reward-based policy training would need suitable observations of actions, their consequences and rewards, plus a reward calculation checked against those independent outcomes. A dataset of positions alone does not supply all those consequences.

A useful first result could be a forecast for one body region with missing states plainly identified. Predicting it, replaying a recording, controlling a virtual body, intervening in a real body and leaving a person with a lasting skill would still be different achievements. None is established by a more lifelike render.

% R4 full-depth accessible insertion candidate; existing manuscript untouched.
% Inherit the selected master's sans-serif text, code and mathematics fonts.
\subsection{One muscle pulse, two simulated trajectories}

Give the computer hand one muscle command. Does its calculated movement change? That is the question tested here. The installed \texttt{myohand\_r} mechanism represents a human right hand attached to a fixed passive torso. Two copies of the same starting state were compared, with a short excitation added to one of them. No recorded brain activity, muscle electrical signal or measured human movement drove or verified this test. It checks how the simulation responds to a command; the larger biological and training questions remain open. The software was \texttt{myo\_sim} 0.2.3 at commit \texttt{93b0ca8f4ec90c9899ee7f05fee561e9911da91b}, MuJoCo 3.14.0 and NumPy 2.5.3.\cite{myohandMyoSimPinned2026,myohandMuJoCo3140,mujoco2012,myohandNumericalRecord20261105}

\subsubsection{Change the input; keep the starting state}

The compiled hand has 23 hinge coordinates: 23 scalar angles that describe its movable joints. Angles use radians. It also has 39 muscle--tendon actuators, the parts that calculate a muscle-driven pull, and 39 activation states, the internal numbers describing the muscle models' activation. There are zero sensors. A calculated activation is therefore not a measurement of muscle electrical activity, and a calculated force is not a force-instrument reading.

The two branches began with matching complete recorded initial states, including joint angles, velocities and activation states. One branch received no additional excitation. The other received excitation 0.2 at the actuator named \texttt{FDS2\_r}; all other commands stayed zero. Excitation is dimensionless here: 0.2 is a value on the model's control scale, not 0.2 volts, 0.2 newtons or an EMG amplitude. The command drives the muscle model rather than directly requesting a joint angle.

Each branch advanced through 100 steps of $0.002\,\mathrm{s}$, for a nominal $0.200\,\mathrm{s}$ of simulated time. Both used the simulator's native Euler setting, its specified numerical rule for advancing the state. Counting steps from zero, the pulse was present at steps 10 through 59. Those commands apply over the pre-step interval $0.020\leq t<0.120\,\mathrm{s}$. The unpulsed hand could still move under gravity and passive mechanics. Comparing the two trajectories accounts for that shared motion; comparing only the pulsed hand with its starting pose would answer a different question.

A forward evaluation recalculates quantities that depend on the current state, such as the endpoint force. One was made at the beginning and one at the end of each branch. There were therefore 200 explicit integration steps and four explicit forward evaluations. The executed construction path also implies two compilation calls and one internal forward evaluation while the model was being built. These counts do not expose or count every operation inside the native solver.

The built mechanism must be described as it actually was. Leaving its optional \texttt{inertia\_floor} setting at \texttt{None} did not leave all inertias untouched: the native builder changed each diagonal rotational-inertia component of four fixed scaffold bodies from $1$ to $10^{-4}\,\mathrm{kg\,m^2}$. Rotational inertia describes resistance to a change in turning motion. The comparison includes that construction change. No extra physical settings were retuned and no fallback to the alternative \texttt{implicitfast} setting was used. Position, orientation, body scale, contact behaviour and fidelity to a real hand remain uncalibrated by this test.

\subsubsection{Read the difference at the same final state}

Write $P$ for the pulsed branch, $C$ for the control and $T=0.200\,\mathrm{s}$ for the nominal endpoint. The reported comparisons are
\[
\Delta q_{\infty}=\max_{1\leq j\leq23}|q_{P,j}(T)-q_{C,j}(T)|,
\qquad
\Delta a=a_{P,\mathrm{FDS2}}(T)-a_{C,\mathrm{FDS2}}(T),
\]
\[
\Delta F=F_{P,\mathrm{FDS2}}(T)-F_{C,\mathrm{FDS2}}(T).
\]
The first expression subtracts the two final angles for each of the 23 hinges, ignores the sign and reports the largest separation. It does not report the largest separation at any time during the run. The next two expressions keep the sign: pulsed activation minus control activation, and pulsed tendon-actuator force minus control force. Activation has no physical unit; the force uses newtons and the model's sign convention. Both endpoint forces were recalculated at their final states, avoiding a comparison with force values left over from an earlier step.

% Shared by all four proposed R3 inserts; no added packages or font defaults.
\begin{table}[htbp]
\centering
\caption{Final contrasts for one simulated right-hand mechanism. All values compare the pulse branch with its matched control at nominal $T=0.200\,\mathrm{s}$. The angle value is the largest absolute separation across 23 hinges at that endpoint; activation and force retain the pulse-minus-control sign. Activation and force concern the selected \texttt{FDS2\_r} actuator. Force uses the synchronized final forward evaluations and the model's sign convention. These are simulated contrasts, not errors against biological measurements. The displayed digits preserve the recorded numerical values, without establishing biological accuracy.\cite{myohandNumericalRecord20261105}}
\label{tab:myohand-numerical-response}
\begin{tabular}{p{0.53\linewidth}p{0.26\linewidth}p{0.10\linewidth}}
\hline
Quantity & Contrast & Unit \\
\hline
Largest absolute final joint-coordinate separation, $\Delta q_{\infty}$ & $0.5027347840656727$ & rad \\
Selected final activation difference, $\Delta a$ & $0.05205729904534248$ & none \\
Selected final tendon-actuator force difference, $\Delta F$ & $-6.407722136774943$ & N \\
\hline
\end{tabular}
\end{table}


A separate audit checked the saved evidence using standard-library software. It completed 132 checks without a failed check and did not rerun the hand. The required saved numerical channels contained no non-finite values, and the recorded native warning counts were zero. Four saved lists describing the endpoint mass matrices contained positive values. Those lists are numerical diagnostics of how the model resists acceleration. The raw matrices and the operands used for residual calculations were not independently exported, so the audit could check the stored results without independently solving those matrices again.

No contact was observed in this case. At the final control endpoint, the recorded soft-limit excess was $0.0004121829780273552\,\mathrm{rad}$. A soft limit resists motion near a boundary; this observation prevents it being described as an absolute hard stop. The check supplies neither proof of energy conservation nor evidence that changing the numerical settings would give a converged answer. No physiological threshold was used to decide whether a contrast was large enough.

The pulse changed the simulated state under the stated conditions. The table compares two calculated trajectories; it does not measure prediction error against a person. Nor do 200 connected time steps become 200 people or 200 independent experiments. No statistical $p$-value follows from that count. The retained decimal places identify the numerical result without establishing accuracy against biology or repeatability across other cases.

\subsubsection{Keep the unanswered questions visible}

The recorded run used CPU dynamics. Its BLAS numerical-library pool reported one thread, and its dynamics pool size was zero. Temporary parallel work during compilation was not measured. The 263 observed loaded-file identities describe what that bounded inspection saw, rather than every system dependency. Resource use and speed still need analysis. The reported process high-water memory value was 5,187,464 KiB; it can carry history from before an exec transition and cannot be assigned to the hand model alone. It establishes no model-only 3 GiB memory bound. Tool/native timing near 448 ms and timing inside the probe near 491.927 ms use unreconciled boundaries. These values establish neither complete useful end-to-end latency nor a speed improvement, and they do not partition every CPU/GPU cost.

The hand supplies calculated activation, force and motion, with no sensory measurements. Myoception names the programme's proposed sensory side of muscle-driven movement, including touch and body-position feedback. This test does not record those sensory nerve signals or demonstrate a biological feedback loop. It also trains no controller or Large Movement Model and tests no generalization. The wider proposal still has four training-source groups: neural observations, muscle measurements, measured movement and physics-generated examples. Using this generator for training would first require qualified mechanics, controls, calibration, uncertainty, source rights and data lineage, then independent tests against held-out real measurements. More examples and more varied examples remain separate hypotheses. The diagnostic reopens no protected tests or quarantined sources.

The present record has no successful scaling curve or trained body-wide world model. In LINK, eighteen fits and 720 saved epochs still fail both model-admission gates and all five benefit claims. Adding a broader proposal does not undo that outcome. New data must face a new qualified comparison, with the original failures intact and a test strong enough to expose another one.

% Prospective task/data protocol; no new experiment or admitted model is reported.
\section{Follow the action, then test the prediction}

Watch a ball pass from one hand to the other. Where were the hands? When did one
let go? When did the other catch it? A model might follow the movement and still
get the handover wrong. That is why the next proposed study needs more than a
picture that looks convincing. It needs observations that can disagree with the
prediction.

We aim to connect a stated intention with an action as closely as possible in a
defined task. Understanding the process between them should make that connection
more specific. The test must still say whose movement it concerns, under which
conditions, how much error is tolerated and which independent observation checks
the result. It cannot turn success in one task into a complete map of the mind.

\subsection{Three recording views of the same event}

The scalp recording, arm sleeves and left/right tripod views introduced in
\emph{From a reach to a testable prediction} would have to agree on which event
they observed. The two surface-EMG sleeves would record one arm each; scalp EEG
would record scalp electrical activity, EMG muscle electrical activity, and
calibrated cameras the hands and arms. Their sites, units, capacity, clocks and
delays need checking before a handover can be called a shared observation.
OpenBCI or an equivalent configurable system remains an unacquired, unqualified
candidate.

For this hand task, the motion-capture relay-composite would join the left and
right observations. Earlier top and bottom dot maps describe another proposed
configuration. Each needs its own calibration and visibility checks. An unseen
finger remains unseen even if a model can draw a plausible replacement; the
display must keep that prediction separate from measured points.

The virtual body would then predict a response for comparison with the measured
movement. Its joints, muscles and contacts need appropriate reference evidence.
EMG does not directly measure force or incoming sensory-nerve activity. If the
claim concerns contact or force, the experiment needs an observation that can
check that claim.

\subsection{Let the task ask the next question}

Start with rapid hand rubbing. Contact and movement happen together, so matched
conditions are needed to find what each contributes. Record the requested
knuckle movement and the stretch with interlaced fingers, palms turned outward
and arms extended as distinct actions. The cold-hands example explains the
rubbing; it does not establish a cold-exposure experiment or a measured sensation.

Then pass a juggling ball between the hands---the proposed hot-potato task---and
progress to the preferred two-ball juggling task. One hand's release changes what
the other must do. The model now has to account for both limbs, the object and
the timing of a catch. Check it against simpler predictions from the cue or
previous movement; those may explain an apparent success without useful brain
information.

The boxing stage would distinguish stated punch types from independently
observed movement. It would become NeuroBoxer, a proposed virtual game. A
prime-era Mike Tyson opponent gives the game a recognizable aim; it supplies no
physical opponent or guarantee about injury. Boxing also uses the hands and
wrists for stabilization. Footage of an expert can show movement and tactics,
subject to its rights and camera limits. It cannot supply that expert's brain
signals, muscle signals or force measurements.

Toy archery adds another relationship. One hand holds while the other draws and
releases; the fingers change as the arrow leaves. The proposed apparatus has
arrows with small, cup-shaped suction tips and an over-arm pouch with several
arrows. NeuroArcher would turn the measured
coordination into action on a virtual target. A hit in the game still needs the
appropriate test before it can mean anything about biological control.

The complete snap described in \emph{From a reach to a testable prediction} tests another
relationship: preparation, contact and return must all agree with the specified
movement. The cup reach remains the simpler precision comparison. Both need
organism-specific anatomy, contact, timing and real-reference evidence; a
human-looking animation cannot establish that a monkey can perform the snap.

\subsection{Hands first; a brain contribution needs its own test}

The first game milestone would use measured body motion. A later milestone would
ask whether permitted EEG observations can control a specified virtual action
without using withheld muscle or movement targets. Eye activity, small covert
muscle contractions, residual movement, cues and information from the future
could all make a prediction look better. Those alternatives need declared
controls. Looking hands-free is not enough.

Keep the original brain-only, body-only and paired comparisons, including the
simple prediction that the target stays as it was. Set the allowed history and
the original success rules before fitting anything. Keep required people,
sessions and tasks separate before making time windows, and learn preprocessing
from training data only. A pretrained model's earlier exposure must be checked
before a recording can serve as independent final confirmation.

Recorded body observations could also predict a stated recorded EEG feature.
That reverse prediction is conditional: several brain states might fit the same
movement. It does not deliver a sensation. Qualified data and simpler baselines
come before paired fitting or fine-tuning. Reinforcement learning would need
specified actions, resulting transitions and rewards, rather than merely a
prediction error. A reward inside a simulation is not a biological result.
Final held-out tests stay outside training and selection; a running model stays
frozen until a candidate update is independently evaluated.

\subsection{A recording device does not become a writer}

The output boundary in \emph{Write sensation. Enact action} applies to these
tasks too. EEG and surface EMG record; neuromuscular electrical stimulation
(NMES) and functional electrical stimulation (FES) can cause contractions or
functional movement. A body-output experiment or sensory writer would need its
own qualified interface and permission. Physical spinal work remains read-only,
and these physical additions remain outside the acquisition and virtual-work
budget.

A proposed blindfolded condition removes visual feedback; it does not prove
neural control or learning. To ask what was retained, remove trainer information
while keeping the necessary movement bridge and comparable task conditions.
Check what the trainee can still do against the original controls. These
recordings, writer experiments and retention tests have not been performed
here. The current LINK failures remain failures; the sequence gives the next
paired study questions it would have to answer.

% New prospective section. Not assembled, rendered, published or scientifically accepted.
% Preserve earlier studies, result tables, citations and frozen failures as historical evidence.
\section{Measure both ends. Test the translation.}

The following prospective protocol tests the paired learner described in ``Two pairs. One conditional learner'': can permitted observations on one side of a sensorimotor task predict withheld measurements on the other, on recordings the model has not seen? The apparatus remains non-invasive EEG, non-invasive surface EMG, a calibrated motion-capture relay-composite, an open musculoskeletal-simulation physics-engine with a lifelike explanatory view, and an open multimodal model. Its acquisition and validation are unfinished.

The earlier brain--muscle and measured-movement--virtual-body comparisons define the two pairs. Here the protocol specifies their records, targets, calibration and independent tests. Anatomical form constrains possible mechanisms; measured function must test the proposed connections. A conditional learner's combined prediction must answer to measurements outside itself.

\subsection{Keep the observations distinct}

Keep raw EEG, raw EMG, derived features and movement estimates in separate
records. Identify the participant and session in each record, then declare the
sensor sites, references, coordinate conventions, units, sampling and timestamp
semantics, and clocks. Trace every processing step and preserve missing intervals and uncertainty. Retain task
instructions, stimulus events, observed behaviour and any separately obtained
participant report. Independent participants or recordings cannot supply a
simultaneous brain--muscle--movement observation.

Surface EMG records muscle electrical activity. A sensory-loop claim also needs
known stimuli and appropriate measured feedback, contact or force channels;
muscle recordings do not provide incoming sensory-nerve traffic. Preserve the
actual spatial resolution of EEG measurements and source estimates. Templates
and explanatory cutaways do not register participants or localize individual
neurons.

\subsection{The spine shapes the loop}

The spinal cord changes how sensory input becomes muscle output. Account for
local coordination, sensory gating, reflex gain and some lasting adaptations,
rather than reducing it to a transmission delay. Separate a response to the
current posture or task from a change retained after training. Keep both apart
from adaptation in sensory receptors, peripheral pathways, neuromuscular
junctions and mechanical tissues. These processes interact; changed output alone
does not locate a memory.

Human H-reflex conditioning and the constrained mouse adaptation studied by
Lavaud and colleagues ask different questions.\cite{thompson2009retained,thompson2013retained,lavaud2024}
The human experiments do not establish brain-independent skill learning. The
mouse study does not establish unrestricted walking or reaching. Circuits
required for acquisition or recall are not thereby exclusive memory-storage
sites. Our human comparison uses the abstracts, selected results and limited
methods passages. The mouse comparison also draws on relevant author-manuscript
passages; equivalence to the final published edition and the complete methods
and statistics have not been established here. Passive proprioceptive training
motivates a peripheral observation model; its psychophysical and spindle-afferent
experiments were separate.\cite{ackerley2022retained} Those observations do not
establish complex memory stored in peripheral axons.

Define the function before choosing its computational substitute. Candidate
functions include gain and inhibition, phase- and load-dependent sensory
responses, coordination across supported muscles or joints, relevant rhythms,
and acquisition, retention and expression of a declared adaptation. Give each
function its own task, species, anatomical scope, measurement and reference.
Naming a mechanism does not show that the model implements it. A hand task
cannot inherit an unmeasured locomotor circuit.

Record task context and history. Distinguish the state within a trial from a
learned parameter change; retain elapsed time, intervening activity, resets,
drift and uncertainty. Recorded spinal or peripheral activity needs supported
channels of its own. Where those observations are absent, label internal state
as estimated or as a simulated controller. EEG, EMG and movement cannot be
relabeled as recorded spinal activity or identified interneurons.

Declare controls and endpoints before fitting or experimentation. Test
immediate context effects separately from retention after a declared interval
and from peripheral or tissue changes. Preserve the original held-out subjects
and tasks, eligible independent measurements, matched baselines and ablations.
Check alternative explanations, time shifts, missing inputs and artifact
leakage. Lower movement error alone cannot locate a biological adaptation,
establish its cause or certify useful plasticity.

To externalize a specified function, propose a substitute that can be tested
through declared peripheral and cerebral interfaces. Compare latency, state
dependence, stability, uncertainty and measured outcomes with the specified
biological function. Computational feedback, peripheral stimulation and restored
biological function require different evidence. A matched virtual trajectory
does not demonstrate spinal replacement, relocated learning or clinical repair.
Keep the spinal-cord site, measured function and proposed computational
substitute distinct. Physical spinal scope remains read-only.

\subsection{Give the two directions different targets}

For brain-to-limb prediction, give the model EEG and only the declared context
available at inference. Test its predicted muscle and movement observables
against withheld measurements. A brain-only condition receives no measured EMG
or movement input. If those inputs help, evaluate them in a separately labeled
hybrid condition.

For limb-to-brain prediction, give the model muscle activity, movement and
declared context. Its answer must allow for different patterns in the supported
recorded EEG features. A conditional distribution assigns probabilities to
specified feature values or ranges, given those inputs. Compare that prediction
with independently retained EEG, then report uncertainty and alternatives.
Movement can narrow the possibilities without identifying a unique hidden
brain state or subjective experience. Agreement between learned models is a
consistency check; it cannot supply the measured target.

\subsection{Reading a signal and writing an output}

The protocol records EEG and EMG; its two prediction directions do not themselves deliver sensation or enact a movement. ``Write sensation. Enact action'' defines the computer and physiological boundaries. Testing a writer would additionally require a separately qualified sensory writer or transducer, a specified actual response and distinct device, measurement, safety and clinical evidence. Indirect sensory feedback and direct neural writing remain different interfaces; a predicted EEG feature is neither. Enaction needs an appropriate virtual, robotic or biological output interface. A robot command has an action role, not the anatomical role of an efferent nerve.

Physical sensory writers, transducers and body actuators remain excluded from the acquisition and virtual-work plan, unqualified and unpriced. This proposal authorizes no stimulation, participant acquisition, biological actuation, new equipment purchase or training.

\subsection{Align before learning}

Calibrate the acquisition views. Verify their clocks against recorded events,
offsets and drift, while preserving occlusion and dropout. Keep a missing
observation missing: a predicted finger position is not a measurement. Declare
frames, scale, joint conventions and acceptable spatial, temporal and end-to-end
errors before selecting equipment or models. Position, joint, contact and force
claims each need an appropriate independent reference.

Specify how the virtual body represents subject scale, muscles, tendons, joint
constraints and contact. Declare its integration settings and observation
models. Biomechanical simulation predicts mechanics; VFX rendering makes those
predictions visible. Judge individual parts and their assembly against the
original anatomical, physical, functional and real-reference-series criteria.

\subsection{Build an open corpus without inventing pairs}

Inventory neural structure, neural activity, muscle structure and muscle
activity as separate domains. Use measured movement and task context as
constraints. For each source, record rights, version, access conditions,
species, participants, tasks, montages, units, clocks, preprocessing and missing
paired modalities. Broad open-data coverage remains a continuing objective,
not an assertion that every dataset is accessible, licensed, compatible or
already used. Check code, pretrained weights and training-data rights separately.

The candidate uses an encoder for each modality, explicit input-availability
masks, a shared conditional representation and separate decoders for the two
directions. Keep simple EEG-only, EMG-only, movement-only and hybrid baselines.
Define pretraining, post-training and fine-tuning tasks before fitting. An
existing encoder supplies a candidate component, not an already established
brain--body model.

\subsection{Make the reward answer to measurements}

First test paired supervised prediction in each direction. Compare any later\linebreak[4]
reinforcement-learning or other closed-loop refinement with that baseline.
Evaluate independent movement measurements, task outcomes, relevant contact or
force, latency and physical validity. Freeze reward units, terms, weights and
decision thresholds before evaluation. The candidate cannot invent the target
against which it earns a reward.

Backpropagation updates model parameters. Reverse prediction estimates another
observable. Measured biological feedback comes from the body. Identify which
process the study actually implements. A supervised prediction error alone does
not establish reinforcement learning: a proposed control study must also declare
its state, action, transition and reward.

Partition eligible participants, sessions and novel tasks before preprocessing
or model selection. Preserve TRAIN-only normalization, frozen seeds and
checkpoint selection, untouched held-out evaluation, uncertainty and original
negative controls. Include informative missing-modality, time-shift, artifact,
drift and alternative-inverse comparisons. Keep reserved recordings closed and
retain completed fits and failed benefit conjunctions. Evaluate frozen inference
separately from any proposed adaptive update. This prospective protocol admits
no new fitting, optimizer update or model/backend promotion.

\subsection{Show what was measured and what was inferred}

Show the selected source clock and the inputs available to each prediction.
Place recorded EEG, muscle activity and movement beside their predictions, with
uncertainty, errors, gaps and stimulus or contact context. Identify recordings,
inferred features and simulation as separate layers. Preserve site/channel
inspection and source-time play, pause, seek and step. Unsupported peripheral
or spinal pathways remain visible gaps. An animated cord cannot substitute for
a spinal recording or supported anatomical localization.

\section{Demonstrated Results and Open Questions}

The current LINK comparison completed and failed its original improvement
requirements. Earlier results support certain recording-specific predictions
and reject fixed-core transfer in favour of a locally calibrated specialist
on a separate old-day recording. The muscle-driven tests distinguish successful
reduced human-arm correction from rejected native macaque controllers. Source
and geometry defects keep their withdrawal and repair records. These are
bounded results that a reader can inspect, each with its own tested scope.

The three later hand-related results must stay separate. The neural classifier
correctly identified 116 of 150 later attempted-finger labels in one locally
calibrated participant. In the linked simulation, 41 of 60 supported source
intentions became correct contacts while all 54 issued supported commands were
executed and 13 unintended presses remained. The separate goal-driven hand
policy completed 32 of 32 constrained simulated phrases. Those fractions have
different denominators because they answer different questions.

A simulated hand can execute its command accurately even when the neural
decision chose the wrong finger. A hand policy can learn a requested key
sequence without recorded brain activity driving it. A simulation can be
mechanically checked without proving that a person learned the same action.
None of these results establishes biological teaching, a neural policy for the
phrases, simultaneous voluntary brain--spine--nerve--muscle control, universal
mastery or a unique translation of private experience.

The next comparison depends on measurements and checks still missing. The
proposed neural-context branch has not run, and current fitting is not admitted;
reserved recordings stay closed. Complete independent checkpoint/update,
random-state, reload and trajectory evidence is absent. Real-world joint and
force measurements would test the mechanics, while subject registration and
resolved units and clocks would establish which body and interval the values
describe. Source-aware functional maps, lifelike 3D reference series, whole-GUI
acceptance and the limits of disconnected service each need their own review.
Another report check cannot answer those questions.

Can measured neural and bodily relationships make physical learning more
effective and leave a trainee more capable after assistance stops? That remains
the question to test. The present software helps formulate it and exposes where
its assumptions fail; it supplies no skill file to accept on faith. Progress
requires experiments specific enough to show a benefit, or show where the
proposed route fails.


% R6 current source scope; no historical glossary presented as current.
\section*{Scope of this source edition}
These unrendered manuscripts retain the computational studies and add the
expanded training proposal and one recorded numerical hand diagnostic. The
diagnostic compares two simulated trajectories. It supplies no new neural fit,
reserved-data evaluation, biological teaching or validated whole-body model.
The definition supplement and source projection now contain the current
registry's 200 entries, revision 34, including their distinct proposed scopes.
No PDF or native export has been generated from this source edition.

Whole-manuscript scientific, editorial, citation, rights, font, layout, native
and export assessment remains unfinished. Current vocabulary does not change
the original numerical results or establish new model capabilities.

\section*{Evidence, scope and edition status}
\addcontentsline{toc}{section}{Evidence, scope and edition status}
The computational studies in these unpublished working manuscripts lead to a
separate proposed paired experiment. No new fitting, reserved-data evaluation,
model promotion or biological teaching is reported here. The negative results
constrain the next test; they do not finish it. The common optimizer still needs
independent implementation checks, including final-checkpoint and update
verification. Reserved confirmation, complete GUI and offline checks, and
applicable anatomical and functional reference comparisons also remain open.

The original tolerances and reserved recordings remain unchanged. A failed
comparison cannot become a successful one through revised wording. Scientific,
clinical, physical, real-reference, rights and runtime claims still require
their original validation. Manuscript sources, citations, figures, rendered
outputs and publication have not completed their checks. Source support for a
negative report establishes neither full scientific validity nor a qualified
interface or release.

This source edition uses the current shared glossary, revision 34, with 200 entries.
Its current JSON projection retains every canonical term field, alias, proposed-name record and source history.
The definition supplement keeps compatible earlier explanations and updates current role meanings; the six additional entries remain proposed programme terms.
The full source fingerprint and current definitions appear at the start of the supplement.
Current terminology does not revise earlier measurements, reopen reserved tests or establish a proposed capability.

A generated illustration can explain a mechanism; it cannot replace
the working model or measured evidence.

\ifdefined\RSevenSixFourTechnicalPackagingRecord
\paragraph{Reproducibility record.}
The retained packaging record describes where files are stored. Moving opaque
resources outside the candidate HTML, the browser page's textual entry file,
reduced that entry from 113,298,857 to 34,295,597 bytes. The complete candidate
directory was 93,831,812 bytes and online delivery required its assets and whole
directory. The separate static offline export compiled from the original source
was 113,581,899 bytes. Views dependent on the application's service interface
(API) were unavailable offline. The recorded review could not use the
candidate's browser route and did not promote it; these receipts establish no
newer browser qualification. The counts measure neither interaction nor loading
latency, memory, crashes or GPU performance. Whole-directory online delivery and
original-source offline access remain separate claims.
\else
The packaging comparison in \textit{Application Performance} distinguishes
whole-directory online delivery from the separate original-source static export,
whose API-dependent views were unavailable. The candidate's browser route remains
unverified and unpromoted. File sizes establish no interaction, loading, memory,
crash or GPU performance result, nor a newer browser qualification.
\fi

The requested public byline is \texttt{@eliminatenoise}. Author identity,
applicable affiliation and repository submission requirements must be resolved
before publication. No affiliation or publication acceptance is implied.

The retained metadata and bounded passages for Logothetis and colleagues,
Tang and colleagues, the quantum-measurement scientific background, and the
NIST measurement-uncertainty guidance support only their stated uses here.
Complete manuscript citation and claim verification remains pending.

\clearpage
% Proposed accessible fragment; use the host edition's sans-serif text/math setup.
% Registry revision 20; SHA-256 cc9676276550560eddc941ee08874019d7566486be709a439c732b0996a1c1e8.
\section{Exactly What the LINK Controls Change}

A model can improve its average by doing better on one recording and worse
on another. That is why the LINK comparison keeps four results visible:
two finger groups on each of two validation dates. A better average cannot
cancel a harmed group.

The study pairs recorded neural features with measured positions of the index
finger and the combined middle, ring and small-finger group in one macaque.
It trains on January 27 and 30 and February 4 and 5, 2020, then selects models
on February 6 and 11. Those positions retain the source's normalized units.
They are not forces or a reconstruction of every finger joint. There is no
simultaneous muscle, peripheral-nerve or spinal recording in this comparison.

Every example contains eleven successive neural bins, ending at the position
being predicted. A twelfth bin supplies the next-neural target used by the
extra-task model. All twelve must be in the same trial, with adjacent recorded
times between 18.5 and 21.5 milliseconds apart. Even a baseline that never
uses the twelfth bin needs it to be present. Everyone is therefore compared
on the same examples. The ten earlier gaps span a history before the current
bin; the nominal 20-millisecond interval does not turn that history into a
measured nerve-conduction delay.

The 192 neural inputs contain 96 transformed threshold-crossing counts and
96 archived power features. Both sets follow the same physical electrode
contact order. A contact is not a permanently identified individual neuron.
Scaling uses only the four training dates, counting each trial-owned frame
once. Repeated overlapping windows do not give that frame extra influence
over the scale, and validation answers do not supply scaling statistics.

The first training date fixes the contact order. The count and power columns
move together when reordered, so both descriptions still refer to the same
declared physical contacts. The count transformation takes the natural
logarithm of one plus each nonnegative count. It keeps a zero count at zero
while compressing larger counts. The archived power features keep their
supplied numerical scale; the two feature types are not converted into one
physical quantity merely because they enter the same model.

The scaler calculates a separate mean and population standard deviation for
each transformed feature. Population standard deviation means subtracting the
mean, squaring those differences, averaging over all the fitting frames and
taking the square root. Those calculations use float64, the declared numerical
format with more precision than float32. Each unique frame belonging to an
eligible training trial contributes once across the training dates, even when
many overlapping windows later reuse it. The scaler therefore describes those
frames, rather than a new weighting produced by how often they appear in
windows. Validation supplies neither a mean nor a spread.

Normalization subtracts the training mean and divides by the saved training
spread. Exactly zero spread is replaced by a divisor of one, which avoids
division by zero for a constant training feature. There is no other minimum
spread: a small nonzero value is not raised to a chosen floor. There is no
clipping to trim large normalized values. This is the exact retained numerical
rule, not a general recommendation to change it after inspecting a result.
Replacing a zero divisor does not fill a missing observation or prove that
the underlying biological signal is constant.

The transform is calculated in float64, then its output is cast to float32
for the model. Casting means representing the computed values in the other
numerical format; rounding can occur. It does not improve an electrode's
physical accuracy or settle an uncertain unit. The position targets keep
their archived float32 values. The constant reference uses the mean of unique
trial-owned training positions. Nonfinite measurements, unknown contacts or
incompatible data descriptions cause a failure, rather than silently dropping
rows or filling them with invented values.

The source descriptions leave several conflicts visible. A threshold-crossing
attribute says volts, while the quantity is interpreted as event counts; the
comparison makes no voltage conversion. It also retains the archived power
values without an unexplained multiplication by four in the author's loader.
Neither choice establishes absolute power calibration. The documented
32-millisecond interval differs from the retained nominal 20-millisecond
clock. The observed timestamps and their allowed jitter define this recipe;
the conflict remains unresolved, rather than being repaired by silently
changing the interval. Full author trial counts must agree with trial-stop
timestamps, and time-series references that are uniformly one entry short
are identified separately. Changing units or discarding an inconvenient
boundary would change which experiment was being compared.


The recurrent model reads each window from its declared starting state.
Its numerical memory is reset between windows; it does not carry the previous
window's final state into the next one. Adjacent windows may reuse some source
samples, but that overlap is not an extra memory channel. Each prediction
must be accounted for by the input window and the fixed model it declares.

Within the window, each new bin updates the hidden values. The position output
combines the final hidden values with another description: for every neural
feature, its mean across all eleven input bins. That mean retains the feature's
average level while discarding its within-window order. The final recurrent
state can depend on the sequence. The output therefore has access to both
descriptions; a benefit from history alone cannot be attributed uniquely to
recurrent ordering. That is one reason for the separate order control.

A separate linear output predicts the next neural feature vector. It remains
present in both candidate architectures, but has zero objective weight for
position-only fitting. Its presence does not give that model an additional
movement target. Likewise, an old software name for the position output does
not turn the normalized position estimate into velocity. These statements
describe the retained recipe; they do not independently verify every state
reset, prediction or update in the completed execution.

The first control asks what happens when the training pairing is broken.
Imagine each example's neural history and its next-neural target held together
as one packet. The control moves that packet to another example within the
same training date, while the position answer stays where it was. It moves
whole eligible windows, not individual channels or bins, and not whole trials
as blocks. Overlapping windows can move separately. The same fixed shuffle
serves every epoch of that run. It is shared by the two model types starting
from the same seed, but differs across seeds and recording dates.

The two validation recordings remain unchanged for this control. We therefore
fit on mismatched training pairs and ask how the resulting model predicts the
ordinary measured validation positions. Some windows can accidentally stay
in place; common task structure also remains. The control challenges the
declared association without proving that every other association has gone.

The second control asks whether the declared order helps. Number an input's
bins from zero to ten, with ten being the current bin. The control presents
them as nine, eight, seven, six, five, four, three, two, one, zero, ten.
The ten older bins reverse; the current bin stays last. Position answers and
next-neural targets remain unchanged. This happens in both training and
validation. The operation is fixed, uses no random draw and preserves all
the values in each window. A model could learn the reversed arrangement,
so this is a particular order comparison rather than a guarantee that timing
information has vanished.

Both model types start from each of three fixed seeds: 20261007, 20261008 and
20261009. The starting tensors are checked against the saved initialization
for that seed. A separate generator restarts from the same seed in each
paired run and draws one training-row order per epoch. For the association
shuffle, another generator starts at the seed plus one million and advances
across the four dates. Matching seeds specify matching procedures; they do
not themselves prove that the completed execution followed every procedure.

There are 40 epochs per run, with batches of 256 and a retained final partial
batch. Every epoch remains in the record. Each seed selects its lowest
average validation error across the four date/group cells, taking the earlier
epoch on an exact tie. Final predictions average all three selected seed
models equally. Neither control competes to become the chosen measured model.
Two model types, three seeds and three conditions produce the reported
18 runs and 720 saved epochs.

The simple linear references try five fixed penalties: 0.01, 0.1, 1, 10 and
100. A tie selects the smaller penalty, then the current-bin reference ahead
of the history reference if their selected errors tie. A tie between the
two measured recurrent ensembles selects position-only. These rules prevent
a favourable seed, epoch or tie decision being quietly substituted later.

The protocol asks five further questions: does the position-only model beat
current-bin ridge; does it beat its reversed-history control; does the
extra-task model beat current-bin ridge; does it beat its own reversed-history
control; and does the extra-task model beat position-only? Each
comparison requires all three conditions together: a strictly lower average
error, no higher error in any date/group, and a negative upper end of its
paired 95-percent error-difference interval. Lower error is better, so a
candidate-minus-reference interval wholly below zero favours the candidate.

That interval repeatedly samples whole supported trials within each validation
date, keeping every sampled trial's supported frames and using the same draws
for the models and groups. It uses 1,000 draws with fixed seed 20261007 and
linear percentile interpolation. It describes uncertainty within these
validation-selected comparisons; it does not confirm a result in a new animal
or on a reserved recording.

For model admission, the three-part comparison is made against the selected
best linear reference. Every date/group must also have a defined, positive
ordinary prediction score, called $R^2$, and exceed its matched association
control. Ordinary $R^2$ compares error with predicting
that validation group's mean; the separately named training-mean score uses
a different reference. An undefined value cannot count as a pass. This
model-admission rule does not replace the five separate benefit comparisons.

The retained independent postfit checker records each condition separately
for all five comparisons. Every comparison has a lower average error, yet
every one makes at least one date/group worse. Three intervals have an upper
end below zero. None of that cancels the failed no-harm conditions: all five
benefit claims fail, and both models fail their separate admission rules.
These are the existing decisions, not a newly calculated result.

The remaining checks concern how the learning calculation was performed. The
order-control ledgers have empty provenance fields; the source and earlier
control evidence do not provide a new dynamic reversal check. The correctly
negative postfit report also cannot independently verify every optimizer
update, checkpoint or restored random state. A completed run and a saved file
are different evidence from that verification.

The later recordings remain closed: February 22 and 24, 2020, and June 21 and
22, 2023. Here, ``earlier input'' means earlier in the supplied file's source
indices. It does not establish that archive filtering used no later signal,
or that an online biological loop has the same delay. Predicting recorded
positions is a bounded decoding test. It demonstrates neither sensory
delivery nor a nervous system learning and retaining a skill.

\par
% Dedicated current SOURCE export; historical G26/G29 exports remain outside this cohort.
\section*{NeuroTranslation Definitions and Lookup Index}
\phantomsection\label{programme-definitions}
\addcontentsline{toc}{section}{NeuroTranslation Definitions and Lookup Index}
\noindent Canonical registry revision 34; 200 entries; SHA-256 \texttt{3fffa79b\allowbreak{}c9f404ca\allowbreak{}a7d4910e\allowbreak{}9e37de5a\allowbreak{}97371fdc\allowbreak{}1ad11bc5\allowbreak{}4155d309\allowbreak{}ec5f5394}.\par
\noindent A measurement, a computer state and a claim about biology need different words. These entries explain the current programme meanings. A proposed name does not establish a working interface, a trained model or a biological result.\par
\noindent Trainer and trainee name the current guidance roles. A machine-learning model can still be a learner; its fitted parameters do not measure a person's retained learning. Historical names and original source quotations remain historical.\par
\noindent Lookup labels help locate an entry. Some name related or contrasting concepts, so they are not interchangeable. The current JSON projection preserves every canonical term field, alias, source and definition history.\par
\noindent Physical scale and units, computational detail and statistical aggregation remain separate. Missing ranges or measurements stay unspecified. Both paper editions retain the same definitions and scientific scope; explanatory wording may differ.\par

% Stable term ID: a-checked-message-between-software-parts
\begin{samepage}
\subsection*{A checked message between software parts}
\noindent A structured message and interface specify what the values mean and how software parts exchange them.\par
\end{samepage}
\begin{samepage}
\noindent The checks include recording, ordered channels, quantity, timing where supplied, units and body axes. Passing a check does not prove that the caller’s claimed physical source is true.\par\nopagebreak[4]

\end{samepage}


\noindent\textit{Current registry lookup labels:} Translation packet, typed port and API; packet; translation port; typed interface; API; application programming interface\par

% Stable term ID: a-group-of-recorded-neural-features
\begin{samepage}
\subsection*{A group of recorded neural features}
\noindent Several recorded neural features considered together when predicting an outcome.\par
\end{samepage}
\begin{samepage}
\noindent Their joint pattern can help predict movement. That does not establish that the same pattern has the same meaning in another animal, task or day.\par\nopagebreak[4]

\end{samepage}


\noindent\textit{Current registry lookup labels:} Neural population; population decoding; population model\par

% Stable term ID: a-map-of-physical-connections
\begin{samepage}
\subsection*{A map of physical connections}
\noindent A structural map records physical connections. A functional recording measures activity over time, such as which recorded features produced events during a task.\par
\end{samepage}
\begin{samepage}
\noindent The structure can guide a model, but it cannot supply missing time-recorded activity. A drawing, connection map and synchronized functional recording have different evidential roles.\par\nopagebreak[4]

\end{samepage}


\noindent\textit{Current registry lookup labels:} Structural connectome and anatomical prior; connectome; connectomic connectivity; anatomical prior; structural connectivity; structural reference; recorded function; structural versus functional measurement\par

% Stable term ID: a-model-spanning-the-body
\begin{samepage}
\subsection*{A model spanning the body}
\noindent A future model connecting relevant observations and actions across the brain, spinal cord, nerves, muscles and movement.\par
\end{samepage}
\begin{samepage}
\noindent The project does not yet have a matched voluntary recording across that whole chain. Separate datasets cannot fill the missing measurements.\par\nopagebreak[4]

\end{samepage}


\noindent\textit{Current registry lookup labels:} Body-wide functional model; body-wide neural fusion; whole-body neural translator\par

% Stable term ID: a-specified-human-related-feature
\begin{samepage}
\subsection*{A specified human-related feature}
\noindent Studying a defined human-related biological component or feature in a model organism.\par
\end{samepage}
\begin{samepage}
\noindent This can overlap with adding donor cells. One selected immune or neural feature does not make an animal a miniature human or establish human-like cognition.\par\nopagebreak[4]

\end{samepage}


\noindent\textit{Current registry lookup labels:} Humanization; humanized; human-related feature\par

% Stable term ID: ability-that-remains-after-guidance-is-removed
\subsection*{Ability that remains after guidance is removed}
\noindent A learned ability that remains without the training guidance and can be tested again later.\par
\noindent Remove trainer commands while keeping any bridge needed to express the trainee’s own control. Switching off that required route can prevent movement regardless of what the trainee learned.\par

\noindent\textit{Current registry lookup labels:} Retained learning; retention; independent pupil control; retained skill; pupil-alone; necessary bridge held constant; pupil-alone with necessary bridge constant; independent learner control; learner-alone; learner-alone with necessary bridge constant; independent trainee control; trainee-alone; trainee-alone with necessary bridge constant\par

% Stable term ID: action2world
\begin{samepage}
\subsection*{Action2World}
\noindent The hand reaches; the cup moves or stays still. Contact, load and the measured outcome tell us whether the action did what we intended, and that feedback shapes the next attempt.\par
\end{samepage}
\noindent The relationship between an action and its consequences for the body and surrounding environment, including the sensory information those consequences produce.\par
\noindent A predicted consequence is not an observed outcome. Specify the object, reference frame, task, units and uncertainty. Footage supplies movement examples without supplying the performer's neural, muscle or force measurements.\par
\begin{samepage}
\noindent\textit{Organizational level:} whole system; network; signal/feature\par\nopagebreak[4]

\end{samepage}


\noindent\textit{Current registry lookup labels:} action2world; action to world; action-to-world\par

% Stable term ID: activity-state-inside-the-muscle-model
\begin{samepage}
\subsection*{Activity state inside the muscle model}
\noindent The model’s internal state that responds to a command and affects its force.\par
\end{samepage}
\begin{samepage}
\noindent This state can change more slowly than the command. The displayed activation comes from the computer model, not recovered physiological activation from the macaque recording.\par\nopagebreak[4]

\end{samepage}


\noindent\textit{Current registry lookup labels:} Muscle activation; activation; activation dynamics\par

% Stable term ID: adjustment-using-evidence-from-the-new-setting
\begin{samepage}
\subsection*{Adjustment using evidence from the new setting}
\noindent Changing or checking a model using examples from where it will be used.\par
\end{samepage}
\begin{samepage}
\noindent Adjusting a decoder, checking probability scores and measuring a body’s physical properties are different kinds of calibration. The new-day model uses labeled examples.\par\nopagebreak[4]

\end{samepage}


\noindent\textit{Current registry lookup labels:} Calibration; recalibration; calibrated transfer; uncalibrated; physiological calibration; physiological fidelity; subject calibration\par

% Stable term ID: agent2agent
\begin{samepage}
\subsection*{Agent2Agent}
\noindent One software agent passes declared information to another.\par
\end{samepage}
\noindent A proposed link carrying specified information between software agents.\par
\begin{samepage}
\noindent Name the agents, message or representation, direction and tested outcome. Software communication does not by itself supply a biological recording, sensory writer, shared experience or demonstrated skill transfer.\par\nopagebreak[4]

\end{samepage}


\noindent\textit{Current registry lookup labels:} agent2agent; agent-to-agent; agent to agent\par

% Stable term ID: agent2brain
\begin{samepage}
\subsection*{Agent2Brain}
\noindent An agent would give a person feedback during a stated task. Tests would have to distinguish help during that task from learning that remains afterwards.\par
\end{samepage}
\noindent A proposed path from a software agent to a person through specified feedback or a separately qualified neural output interface.\par
\begin{samepage}
\noindent Feedback would reach a person through an ordinary sensory route or a separately qualified direct neural interface. Those routes need different devices, evidence and admission: feedback on a screen is not direct neural delivery. EEG and EMG record; they do not write. The proposed agent-to-person path leaves the reverse observation and feedback direction intact. Help during a task must be tested separately from learning retained after it ends. Agent2Brain establishes no stimulation or physical-action authority. Physical spinal scope remains read-only.\par\nopagebreak[4]

\end{samepage}


\noindent\textit{Current registry lookup labels:} Agent2Brain?; agent2brain; agent2brain?; agent-to-brain; agent to brain\par

% Stable term ID: all-outputs-inside-their-ranges-together
\begin{samepage}
\subsection*{All outputs inside their ranges together}
\noindent How often every output of interest falls within its stated range in the same case.\par
\end{samepage}
\begin{samepage}
\noindent A 90\% range for each separate output does not guarantee 90\% coverage of all outputs at once or of a whole body.\par\nopagebreak[4]

\end{samepage}


\noindent\textit{Current registry lookup labels:} Simultaneous coverage; all-output coverage; whole-body coverage\par

% Stable term ID: an-independent-way-the-model-can-move
\begin{samepage}
\subsection*{An independent way the model can move}
\noindent One independent model movement, such as rotation at a joint. The formal shorthand is DOF.\par
\end{samepage}
\begin{samepage}
\noindent The native arm has seven coordinates and no separate finger movements. The human replay controls two goals while bracing other independent joints.\par\nopagebreak[4]

\end{samepage}


\noindent\textit{Current registry lookup labels:} Degree of freedom; DOF; DOFs; independent coordinate; braced joint\par

% Stable term ID: analysis-timing-and-biological-travel-time
\begin{samepage}
\subsection*{Analysis timing and biological travel time}
\noindent An analysis offset aligns two recorded quantities. A conduction delay concerns how long a signal travels through biology.\par
\end{samepage}
\begin{samepage}
\noindent The 40 ms sensory offset estimates an earlier limb state. The motor target’s 4 ms offset comes from sampling. Neither measures neural conduction time.\par\nopagebreak[4]

\end{samepage}


\noindent\textit{Current registry lookup labels:} S1 alignment offset and conduction delay; 40 ms sensory lag; S1 lag; 4 ms sampling offset; conduction delay; physiological latency\par

% Stable term ID: animal-drawing-and-measured-anatomy
\begin{samepage}
\subsection*{Animal drawing and measured anatomy}
\noindent An illustration explains a proposed animal or body. An empirical anatomical atlas is built from verified observations of anatomy; it cannot be established by a convincing drawing.\par
\end{samepage}
\begin{samepage}
\noindent The antechinus illustration is a proposal, separate from acquired macaque and human sources. The implementation has no matched validated antechinus neural/body atlas and makes no living-animal learning claim from the image.\par\nopagebreak[4]

\end{samepage}


\noindent\textit{Current registry lookup labels:} Antechinus proposal; antechinus; marsupial proposal; hypothetical organism; Proposed marsupial setting; anatomy illustration; illustrative model; empirical atlas; anatomical atlas\par

% Stable term ID: another-brain-area-used-in-the-finger-study
\begin{samepage}
\subsection*{Another brain area used in the finger study}
\noindent A brain area recorded alongside the movement area in the human finger study.\par
\end{samepage}
\begin{samepage}
\noindent The formal name is posterior parietal cortex, or PPC. The study provides attempted-finger labels, not measured trajectories of executed finger movement.\par\nopagebreak[4]

\end{samepage}


\noindent\textit{Current registry lookup labels:} Posterior parietal cortex; PPC; Parietal brain area\par

% Stable term ID: assumed-repair-to-the-native-muscle-model
\begin{samepage}
\subsection*{Assumed repair to the native muscle model}
\noindent A separately named model version replaces unusable zero-rate properties with the software’s default values.\par
\end{samepage}
\begin{samepage}
\noindent The source uses a named Schutte muscle model. Restoring 117 default properties permits numerical tests but does not measure the subject’s muscle response rates.\par\nopagebreak[4]

\end{samepage}


\noindent\textit{Current registry lookup labels:} Schutte muscle properties and default-rate fork; Schutte; zero-rate; assumed-rate fork; class defaults\par

% Stable term ID: attachment-order-and-the-chosen-route-segment
\begin{samepage}
\subsection*{Attachment order and the chosen route segment}
\noindent The ordered connection points and the numbered segment to which a wrapping rule applies.\par
\end{samepage}
\begin{samepage}
\noindent The rule uses the source model’s attachment order. The SDK counts from one and includes both range endpoints; the correction does not fit arbitrary points to an error frame.\par\nopagebreak[4]

\end{samepage}


\noindent\textit{Current registry lookup labels:} Attachment topology and SDK range; attachment topology; one-based inclusive range; humerus-to-ulna segment\par

% Stable term ID: axes-pushes-and-turns
\begin{samepage}
\subsection*{Axes, pushes and turns}
\noindent A coordinate frame states the axes used for movement. A handle wrench combines forces and turning moments measured in that frame.\par
\end{samepage}
\begin{samepage}
\noindent The handle record has three force and three moment components. An unknown mapping into the computer-body axes prevents treating those forces as verified simulator inputs.\par\nopagebreak[4]

\end{samepage}


\noindent\textit{Current registry lookup labels:} Coordinate frame and handle wrench; body frame; XY; physical frame; handle wrench; forces and moments\par

% Stable term ID: biological-tissue-repair
\begin{samepage}
\subsection*{Biological tissue repair}
\noindent Restoration or replacement of biological tissue or its function, supported by direct biological evidence.\par
\end{samepage}
\begin{samepage}
\noindent Better simulated tracking, assisted movement and retained behavior are different outcomes. The prototype does not demonstrate neural tissue repair.\par\nopagebreak[4]

\end{samepage}


\noindent\textit{Current registry lookup labels:} Regeneration; tissue regeneration; repair\par

% Stable term ID: biological-virtualization
\begin{samepage}
\subsection*{Biological virtualization}
\noindent For a movement model, matching the path is one question; matching the living system’s response to an intervention is another.\par
\end{samepage}
\noindent A claim that a digital model reproduces specified biological behaviour and causal dynamics within an independently tested domain and stated error bounds.\par
\noindent Independent ground truth must establish measured error and repeatability within declared conditions and tolerances. For causal dynamics, the comparison also needs independently measured responses to appropriate causal interventions; a similar-looking trajectory alone cannot establish them. Approximation can make a model plausible and may contribute to a validated model when its assumptions, errors and domain are tested and disclosed. A musculoskeletal simulation is itself a model to test, not independent empirical ground truth. Numerical detail and repeated verifier results guarantee neither biological accuracy nor complete physiology, subjective continuity, consciousness transfer or implemented Brain2Body capability. This biological claim is distinct from operating-system virtualization.\par
\noindent\textit{Scope:} validation: Specified observed behaviours and causal dynamical responses under declared conditions.\par
\noindent\textit{Organizational level:} whole system; network; organ; local circuit; cell; signal/feature\par
\noindent\textit{Computational granularity:} Declared model and brain-body interfaces\par
\begin{samepage}
\noindent\textit{Statistical aggregation:} Declared validation aggregation; separate from physical scale\par\nopagebreak[4]

\end{samepage}


\noindent\textit{Current registry lookup labels:} virtualization; virtualisation; biological virtualisation\par

% Stable term ID: body-dimensions-and-resistance-to-motion
\begin{samepage}
\subsection*{Body dimensions and resistance to motion}
\noindent Mass is the amount of matter modeled. Inertia describes resistance to changes in motion; anthropometry means body dimensions.\par
\end{samepage}
\begin{samepage}
\noindent The scaled model keeps assumed source masses and applies geometry-related inertia updates. Its length fit does not identify the animal’s actual masses, strengths or individual bone dimensions.\par\nopagebreak[4]

\end{samepage}


\noindent\textit{Current registry lookup labels:} Mass, inertia and anthropometry; inertia; mass; anthropometry; subject scaling; specimen scale\par

% Stable term ID: body-position-stored-by-the-simulator
\begin{samepage}
\subsection*{Body position stored by the simulator}
\noindent The model’s stored configuration, such as its joint positions.\par
\end{samepage}
\begin{samepage}
\noindent The formal variable name is qpos. Its entries depend on the model and are not automatically measured joint angles. The matching velocity variable, qvel, is explained separately.\par\nopagebreak[4]

\end{samepage}


\noindent\textit{Current registry lookup labels:} Generalized position and velocity; qpos; joint state\par

% Stable term ID: body-model-software-and-its-programming-tools
\begin{samepage}
\subsection*{Body-model software and its programming tools}
\noindent OpenSim provides tools for body models, geometry, scaling and simulation. Its software development kit, or SDK, lets code work with those models.\par
\end{samepage}
\begin{samepage}
\noindent Checks on these tools establish the stated model’s numerical behavior. They do not certify the exact upstream model version or the recorded animal’s anatomy.\par\nopagebreak[4]

\end{samepage}


\noindent\textit{Current registry lookup labels:} OpenSim and its SDK; OpenSim; SDK; software development kit; native SDK; OpenSim 4.5.2\par

% Stable term ID: body-position-feedback
\begin{samepage}
\subsection*{Body-position feedback}
\noindent Proprioception means sensing body position and movement. For example, it helps a body monitor where a limb is. Biofeedback uses bodily information to guide action or learning.\par
\end{samepage}
\begin{samepage}
\noindent The implemented correction uses simulated body position. It does not add measured peripheral nerve signals or show that a biological feedback process taught a retained skill.\par\nopagebreak[4]

\end{samepage}


\noindent\textit{Current registry lookup labels:} Proprioception and biofeedback; proprioception; simulated proprioception; biofeedback; Body-state sensing and feedback; proprioceptive feedback\par

% Stable term ID: body2action
\begin{samepage}
\subsection*{Body2Action}
\noindent The proposed link from body mechanics and sensory interaction to a specified, independently measured task action or outcome.\par
\end{samepage}
\noindent A task outcome makes this proposed link testable: the body’s forces, contacts and timing have to be measured in the relevant context. Body2Action connects body mechanics and sensory interaction to that specified action or outcome; it establishes no working controller. World labels are historical conceptual lookup references, not a strict synonym for every environmental interaction.\par
\noindent\textit{Organizational level:} Declared system or interface\par
\noindent\textit{Computational granularity:} Declared process/interface/model\par
\begin{samepage}
\noindent\textit{Statistical aggregation:} Validation domain stated independently of physical scale\par\nopagebreak[4]

\end{samepage}


\noindent\textit{Current registry lookup labels:} Body to action; Body2World; Body to world\par

% Stable term ID: body2body
% Canonical status: "user_requested_proposed_programme_role"
\subsection*{Body2Body}
\noindent A measured movement can constrain a digital body; a digital body can suggest a movement to test. Body2Body asks which quantities can travel between those settings and whether the resulting movement matches an independent reference.\par
\noindent A proposed research direction relating specified physical-body measurements to digital body models and testing mappings or feedback from those models to declared target bodies.\par
\begin{samepage}
\noindent The user’s physical-to-digital and digital-to-physical wording names a proposed direction, not a proven symmetric inverse. Observation-to-model fitting, retargeting and a separately qualified physical output have different requirements. Homologous and mechanically analogous features can support unpaired pretraining or transfer when their mechanics, scale, contacts, timing and missing features are declared. That learned relationship is not an actually measured neural pairing. Keep inferred and simulated quantities distinct from measurements. Biological virtualization still requires independent biological outcomes in a stated domain; no real body actuation, retained skill or whole-organism equivalence follows here.\par\nopagebreak[4]

\end{samepage}


\noindent\textit{Current registry lookup labels:} Body to body\par

% Stable term ID: brain-imaging-related-to-blood-oxygen
\begin{samepage}
\subsection*{Brain imaging related to blood oxygen}
\noindent An imaging method, called fMRI, that measures signals associated with changes in blood oxygenation.\par
\end{samepage}
\begin{samepage}
\noindent It does not record individual electrical spikes. The paper discusses a published model for this kind of image, not a project result showing that brain scans and spikes are interchangeable.\par\nopagebreak[4]

\end{samepage}


\noindent\textit{Current registry lookup labels:} Functional magnetic resonance imaging; fMRI\par

% Stable term ID: brain-pad
\subsection*{Brain pad}
\noindent A pad shows the proposed interface at the brain.\par
\noindent The proposed brain-side (cortical) interface, drawn as a pad.\par
\noindent A contextual display label in the user's Interoceptomics visual instructions, not a general clinical synonym. The name establishes no acquired neural data, registration, tissue access, function, safety, clinical use, species or physical dimensions.\par

\noindent\textit{Current registry lookup labels:} brain pad; cortical prosthesis; cortical interface; cortical interface prosthesis\par

% Stable term ID: brain-stimulation
\begin{samepage}
\subsection*{Brain stimulation}
\noindent An intervention that applies a specified physical stimulus to alter brain activity.\par
\end{samepage}
\noindent The method must be named, such as transcranial magnetic/electrical stimulation or intracortical microstimulation. These differ in access, risks and capabilities. EEG records; a scalp EEG headset does not provide general sensory writing or safe skill upload. Broad neuromodulation is not identical to precise task-content encoding.\par
\noindent\textit{Organizational level:} Declared game/model/interface role\par
\noindent\textit{Computational granularity:} Declared task-specific model or interface\par
\begin{samepage}
\noindent\textit{Statistical aggregation:} Validation scope declared per experiment; no physical-scale inference\par\nopagebreak[4]

\end{samepage}
\noindent\textit{Source:} A brain-computer interface that evokes tactile sensations improves robotic arm control\par
\noindent\textit{Source identifier:} flesher2021\par
\noindent\textit{DOI:} 10.1126/science.abd0380\par
\noindent\url{https://pubmed.ncbi.nlm.nih.gov/34016775/}\par
\noindent Research precedent for motor decoding plus sensory feedback through intracortical microstimulation; invasive and not a consumer EEG capability\par


\noindent\textit{Current registry lookup labels:} Neural stimulation; Neurostimulation\par

% Stable term ID: brain2agent
\begin{samepage}
\subsection*{Brain2Agent?}
\noindent A recording supplies observations that an agent may use for a stated task.\par
\end{samepage}
\noindent A proposed path from specified brain observations to a software agent through declared recording and interpretation interfaces.\par
\begin{samepage}
\noindent Recorded brain observations would reach an agent through a declared recording and interpretation interface. Where applicable, feedback toward the agent remains distinct from Agent2Brain’s proposed agent-to-person path. Neither direction identifies unique private intent or experience. Recording supplies observations; inference and simulation interpret or model them, while actual sensory delivery requires its own interface. Brain2Agent? remains an open proposal.\par\nopagebreak[4]

\end{samepage}


\noindent\textit{Current registry lookup labels:} Brain2Agent; brain2agent; brain2agent?; brain-to-agent; brain to agent\par

% Stable term ID: brain2body
\begin{samepage}
\subsection*{Brain2Body}
\noindent Connecting brain and body function; the 2 means to.\par
\end{samepage}
\noindent The numeral 2 is the preferred spelling. The programme includes body-to-brain feedback as well as outward pathways; the name does not establish complete biological translation or novelty of body–brain research.\par
\begin{samepage}
\noindent\textit{Organizational level:} whole system; network; local circuit; signal/feature\par\nopagebreak[4]

\end{samepage}


\noindent\textit{Current registry lookup labels:} brain2body; Brain2body; brain 2 body; brain-to-body; brain to body\par

% Stable term ID: brain2brain
\subsection*{Brain2Brain}
\noindent Information links between two brain or neural systems.\par
\noindent The project’s label for studying specified information links between brain or neural systems, including trainer–trainee comparisons.\par
\noindent Brain2Brain is the numeral-style public label for the project’s existing brain-to-brain research direction. Each experiment must specify its signal, task and controls; the name itself establishes no shared thoughts, general skill transfer or successful biological teaching.\par
\noindent\textit{Organizational level:} whole system; network; local circuit; signal/feature\par

\noindent\textit{Current registry lookup labels:} brain2brain; brain 2 brain; brain-to-brain; brain to brain\par

% Stable term ID: brain2motion
\begin{samepage}
\subsection*{Brain2Motion}
\noindent The proposed movement pathway connecting Brain2Body with Body2Action and the resulting measured motion, including the returning sensory feedback.\par
\end{samepage}
\noindent The label concerns the movement part of the wider programme, not every sensory/cognitive function. Brain/spinal/peripheral observations, muscle activity, body mechanics and external outcomes stay distinct. A decoded motion signal, simulation or lifelike animation is not a complete physiological explanation or validated skill transfer.\par
\noindent\textit{Organizational level:} Declared system or interface\par
\noindent\textit{Computational granularity:} Declared process/interface/model\par
\begin{samepage}
\noindent\textit{Statistical aggregation:} Validation domain stated independently of physical scale\par\nopagebreak[4]

\end{samepage}


\noindent\textit{Current registry lookup labels:} Brain to motion\par

% Stable term ID: brain2spine
\begin{samepage}
\subsection*{Brain2Spine}
\noindent A movement command meets a spinal cord that already regulates reflexes, coordinates muscles and adapts through experience. Understanding that middle step helps explain how the same intention can produce different actions.\par
\end{samepage}
\noindent The pathway from brain activity into spinal processing: descending signals influence circuits that combine them with sensory input and the body's current state.\par
\noindent A scalp EEG recording does not identify every descending signal or measure the spinal circuit's hidden state. The physical spine remains read only. This label does not name an admitted stimulation interface.\par
\begin{samepage}
\noindent\textit{Organizational level:} whole system; network; signal/feature\par\nopagebreak[4]

\end{samepage}


\noindent\textit{Current registry lookup labels:} brain2spine; brain to spine; brain-to-spine\par

% Stable term ID: brainnet
\begin{samepage}
\subsection*{BrainNet}
\noindent A network of specified links between brain or neural systems.\par
\end{samepage}
\noindent A working network metaphor for connected brain or neural systems within the programme.\par
\noindent BrainNet pictures specified links between brain or neural systems as a network. It is a programme metaphor; it establishes no validated universal neural network, uniquely coined field or available trademark.\par
\begin{samepage}
\noindent\textit{Organizational level:} whole system; network; local circuit; signal/feature\par\nopagebreak[4]

\end{samepage}


\noindent\textit{Current registry lookup labels:} brainnet; Brain Net\par

% Stable term ID: brain-body-interface
% Canonical status: "user_requested_proposed_programme_role"
\subsection*{Brain–Body Interface}
\noindent A brain recording and a measured hand movement answer different questions about the same reach. A Brain–Body Interface would connect the stated observations and any separately qualified outputs so that the proposed relationship can be tested.\par
\noindent A proposed interface linking specified brain observations with body observations or control outputs in a declared task, with each direction and physical route identified.\par
\begin{samepage}
\noindent This names an interface role, rather than the Brain2Body programme or a model spanning the body. A Digital Bridge is the existing selected-information route around an interruption or between specified recording and control channels; it is related, not an automatic synonym for every Brain–Body Interface. Recording, inference, simulated control and physical output require distinct evidence. EEG and EMG record; a feedback or stimulation output needs its own qualified interface. The label establishes neither complete coverage nor a bidirectional biological loop. No new acronym is introduced.\par\nopagebreak[4]

\end{samepage}


\noindent\textit{Current registry lookup labels:} Brain-Body Interface\par

% Stable term ID: choosing-a-model-and-deciding-if-it-passes
\begin{samepage}
\subsection*{Choosing a model and deciding if it passes}
\noindent Choosing settings and applying a stated pass rule using the validation data.\par
\end{samepage}
\begin{samepage}
\noindent The frozen-transfer rule requires a prediction score at least 0.01 above every specified control. A failing candidate stays rejected rather than being given an easier success rule.\par\nopagebreak[4]

\end{samepage}


\noindent\textit{Current registry lookup labels:} Validation selection and admission; validation-selected; acceptance criterion; admission; failure criterion\par

% Stable term ID: command-sent-to-a-modeled-muscle
\begin{samepage}
\subsection*{Command sent to a modeled muscle}
\noindent The input command that starts the modeled muscle’s activation process.\par
\end{samepage}
\begin{samepage}
\noindent The command is bounded and nonnegative here. It differs from the resulting activity state, recorded muscle electricity and instructions for a biological intervention.\par\nopagebreak[4]

\end{samepage}


\noindent\textit{Current registry lookup labels:} Muscle excitation; excitation; actuator command; muscle control\par

% Stable term ID: communication-route-around-an-interruption
\begin{samepage}
\subsection*{Communication route around an interruption}
\noindent An interface carrying selected information around a modeled break or between stated recording and control channels.\par
\end{samepage}
\begin{samepage}
\noindent A bridge can compensate without repairing tissue or supplying muscle energy. The virtual interruption does not establish a biological injury model or treatment.\par\nopagebreak[4]

\end{samepage}


\noindent\textit{Current registry lookup labels:} Digital bridge; compensatory bridge; brain–spinal bridge; BCI; brain-computer interface; virtual pathway interruption\par

% Stable term ID: comparable-roles-in-different-organisms
\begin{samepage}
\subsection*{Comparable roles in different organisms}
\noindent Comparing roles such as reaching or support while keeping each organism’s anatomy explicit.\par
\end{samepage}
\begin{samepage}
\noindent Related roles do not imply identical neurons or muscles. A combined fictional body cannot replace the measured sources and checks needed for each real organism.\par\nopagebreak[4]

\end{samepage}


\noindent\textit{Current registry lookup labels:} Organism role and homology; organism adapter; limb role; functional coordinate; homology; biped/quadruped\par

% Stable term ID: comparisons-that-test-whether-guidance-matters
\subsection*{Comparisons that test whether guidance matters}
\noindent Matched replay uses the relevant recorded activity under matched task conditions. Stale or unmatched replay uses activity from a different time or condition. No-link and inactive or sham comparisons test whether the proposed information route adds benefit.\par
\noindent These are proposed comparison types, not new animal results. Shuffling disrupts timing; a yoked comparison gives matched help from another case. Guidance amount, timing and later trainee control need their own stated comparisons.\par

\noindent\textit{Current registry lookup labels:} Shuffled and yoked controls; shuffled-input control; time-shuffled; yoked signal comparison; Matched comparisons with altered information; matched replay; matched own-session activity; stale replay; unmatched replay; no-link; sham; inactive comparison\par

% Stable term ID: completed-time-interval-for-counting-events
\begin{samepage}
\subsection*{Completed time interval for counting events}
\noindent A stated time interval groups events into one count. A half-open interval includes its start and excludes its end.\par
\end{samepage}
\begin{samepage}
\noindent Inputs use completed intervals. The exact boundary rule, interval duration and timing must remain explicit rather than being guessed from the display.\par\nopagebreak[4]

\end{samepage}


\noindent\textit{Current registry lookup labels:} Completed time bin and half-open binning; bin; 20 ms bin; 50 ms bin; half-open; bin end; cadence\par

% Stable term ID: computer-processors-memory-and-file-size
\begin{samepage}
\subsection*{Computer processors, memory and file size}
\noindent A CPU is a general processor, a GPU is another kind of parallel processor, and RAM is working memory. Bytes and GiB measure digital data size.\par
\end{samepage}
\begin{samepage}
\noindent The training described in this source edition used four CPU threads; GPU acceleration was neither verified nor used. AI-assisted coding does not establish independent biological validation or an authorship credential.\par\nopagebreak[4]

\end{samepage}


\noindent\textit{Current registry lookup labels:} Computing hardware and data size; CPU; GPU; RAM; AMD; AI; Ryzen AI Max+395; GiB; byte; bytes\par

% Stable term ID: correct-label-rate-and-chance-comparison
\begin{samepage}
\subsection*{Correct-label rate and chance comparison}
\noindent Accuracy is the proportion of labels chosen correctly. A confusion table shows which labels were mistaken for others.\par
\end{samepage}
\begin{samepage}
\noindent A uniform guess among ten balanced classes has a 10\% reference rate. Accuracy and label scores do not prove executed movement or confidence in biological truth.\par\nopagebreak[4]

\end{samepage}


\noindent\textit{Current registry lookup labels:} Classification accuracy and chance; accuracy; uniform chance; ten-class confusion; confusion matrix\par

% Stable term ID: count-prediction-improvement-per-recorded-event
\begin{samepage}
\subsection*{Count-prediction improvement per recorded event}
\noindent Bits per spike describes how much a count-prediction score improves over a defined baseline, divided by the number of recorded spikes.\par
\end{samepage}
\begin{samepage}
\noindent The baseline uses the training mean rate. A positive result concerns estimated average counts, not newly recorded activity or verified stimulation instructions.\par\nopagebreak[4]

\end{samepage}


\noindent\textit{Current registry lookup labels:} Bits per spike; BPS; ΔBPS; bits/spike; likelihood gain\par

% Stable term ID: data-for-learning-choosing-and-checking
\begin{samepage}
\subsection*{Data for learning, choosing and checking}
\noindent Separate portions used to fit the model, choose its settings, and check the selected result.\par
\end{samepage}
\begin{samepage}
\noindent The paper uses local time-ordered splits. If test results have already been inspected, later reuse cannot be described as another independent test.\par\nopagebreak[4]

\end{samepage}


{\rightskip=0pt plus1fil\relax
\noindent\textit{Current registry lookup labels:} Training, validation and test partitions; train/validation/test; train/val/test; split; chronological 70/15/15\par
}

% Stable term ID: data-kept-out-of-a-stated-step
\begin{samepage}
\subsection*{Data kept out of a stated step}
\noindent Data reserved from training or from another specified development step.\par
\end{samepage}
\begin{samepage}
\noindent Independence depends on what was actually excluded. A previously inspected case does not become a fresh test just because it is called held-out.\par\nopagebreak[4]

\end{samepage}


\noindent\textit{Current registry lookup labels:} Holdout; held-out; reserved window; independent evaluation\par

% Stable term ID: describing-motion-and-calculating-what-causes-it
\begin{samepage}
\subsection*{Describing motion and calculating what causes it}
\noindent Describing a body’s position and movement is one task; calculating how modeled forces change that movement is another.\par
\end{samepage}
\begin{samepage}
\noindent The paper calls these kinematics and dynamics. A replayed pose and a muscle-driven trajectory are different results; neither alone proves physiological accuracy.\par\nopagebreak[4]

\end{samepage}


\noindent\textit{Current registry lookup labels:} Kinematics and dynamics; kinematics; dynamics; forward dynamics; numerical plant\par

% Stable term ID: ecotype
% Canonical status: "established_scientific_reference"
\subsection*{Ecotype}
\noindent A population adapted to a coastal habitat is an example of the ecological distinction at issue. Ecotype names the adapted population or form; the surroundings themselves are its environment.\par
\noindent A locally adapted population or form associated with a particular ecological setting.\par
\noindent Ecotype is an established population/ecological concept, not an individual’s surroundings or an obligatory third stage after genotype and phenotype. State the organism and evidence for the claimed local differentiation or adaptation. The retained Turesson source is a historical plant example~\citep{turesson1922ecotype}, not a complete modern evolutionary theory or evidence about the proposed nervous-system experiments.\par


\noindent\textit{Current registry lookup labels:} No additional label declared.\par

% Stable term ID: electrical-recording-in-outer-brain-tissue
\begin{samepage}
\subsection*{Electrical recording in outer brain tissue}
\noindent Electrodes in the brain’s outer layer, called the cortex, record electrical signals outside the cells.\par
\end{samepage}
\begin{samepage}
\noindent This differs from a scalp recording. It also does not observe every cell in a circuit; the source file defines which events and channels were recorded.\par\nopagebreak[4]

\end{samepage}


\noindent\textit{Current registry lookup labels:} Intracortical extracellular recording; intracortical; extracellular cortical events; Electrical recording inside the cortex; cortex\par

% Stable term ID: enactome
% Canonical status: "user_requested_proposed_data_organization"
\subsection*{Enactome}
\noindent A reach has a path, a time and a result. An Enactome would organize observations like these so that the action and what followed can be compared across stated conditions.\par
\noindent A proposed organization of specified action observations and the recorded relationships between actions and their outcomes in declared tasks.\par
\noindent Declare which actions, body parts, tasks, units, frames and times the data cover. Keep commands, measured movement, sensory consequences and simulated outcomes separate. An observed association does not establish causation, a complete repertoire or the person’s unique intent. This proposed data organization differs from the tentative Enactotype individual-profile label.\par


\noindent\textit{Current registry lookup labels:} No additional label declared.\par

% Stable term ID: endohumanisation
% Canonical status: "user_requested_proposed_programme_role"
\subsection*{EndoHumanisation}
\noindent Choose a particular human-related feature, then ask what introducing or studying that feature would change in the model organism. EndoHumanisation names that proposed line of work.\par
\noindent A proposed research stage studying specified human-related genetic features within a model organism and testing their relevance to a declared biological question.\par
\noindent The roadmap places XenoTransgenesis at step 1.1. The working interpretation concerns defined genetic features; the actual manipulation and its established scientific name remain to be specified. A selected feature does not make an animal human or establish human cognition, disease fidelity or an appropriate intervention age. Biological relevance, function and animal/participant, ethical, safety and rights requirements need their own evidence and admission. This entry supplies no experimental procedure.\par


\noindent\textit{Current registry lookup labels:} No additional label declared.\par

% Stable term ID: error-range-based-on-validation-examples
\begin{samepage}
\subsection*{Error range based on validation examples}
\noindent A range placed around a prediction using the errors seen on the data used for model selection.\par
\end{samepage}
\begin{samepage}
\noindent The ranges have constant width for each output. They do not guarantee future coverage when examples overlap in time or their errors are related.\par\nopagebreak[4]

\end{samepage}


\noindent\textit{Current registry lookup labels:} Residual band; empirical residual bounds; validation residual quantile; 90\% per-output residual band\par

% Stable term ID: events-grouped-as-one-neural-feature
\begin{samepage}
\subsection*{Events grouped as one neural feature}
\noindent The source recording process groups events into features that the model can count separately.\par
\end{samepage}
\begin{samepage}
\noindent A group is not a verified lifelong identity for one neuron. It must stay separate from unassigned events and from the channel labels used in other recordings.\par\nopagebreak[4]

\end{samepage}


\noindent\textit{Current registry lookup labels:} Sorted neural unit; sorted unit; sorted cell; sorted feature\par

% Stable term ID: exohumanisation
% Canonical status: "user_requested_proposed_programme_role"
\subsection*{ExoHumanisation}
\noindent Added tissue can survive without doing the job expected of it. ExoHumanisation names a proposed study of a specified human-related graft: what survives, what connects, and what function can actually be measured?\par
\noindent A proposed research stage studying specified human-related donor cells or tissue in a model organism and testing their survival, integration and function.\par
\noindent The roadmap associates step 2.1 with XenoTransplantation. A cross-species graft retains that established meaning; cell survival, anatomical connection, activity and useful functional integration are distinct outcomes. A graft does not establish a human nervous system, a shared mind or whole-organism human equivalence. Species, tissue identity, biological purpose, ethics, safety and rights remain explicit dependencies. This label provides no transplantation procedure or intervention-age window.\par


\noindent\textit{Current registry lookup labels:} No additional label declared.\par

% Stable term ID: expected-event-count-and-its-fit-score
\begin{samepage}
\subsection*{Expected event count and its fit score}
\noindent A model estimates the average event count for given inputs, then a score checks how that estimate fits the recorded counts.\par
\end{samepage}
\begin{samepage}
\noindent The paper uses Poisson negative log likelihood for that check. A predicted average may be fractional; a good count-fit score does not establish a complete biological spike model.\par\nopagebreak[4]

\end{samepage}


\noindent\textit{Current registry lookup labels:} Conditional count mean and Poisson likelihood; expected count mean; conditional count; Poisson; negative log likelihood; NLL\par

% Stable term ID: external-movement-model
\begin{samepage}
\subsection*{External movement model}
\noindent The external part of the body model, constrained by independently measured surface movement and pose in a declared task and coordinate frame.\par
\end{samepage}
\noindent The motion-capture relay-composite retains the source views, clocks, calibration and coordinate frame for the declared task. The earlier upper/lower hand concept and the later synchronized left/right tripod arm proposal are different arrangements; each needs its own visibility, synchronization and error checks. The composite is derived. Hidden internal states, forces and intent are not uniquely measured. LiDAR and optical infrared tracking require distinct hardware, calibration and uncertainty qualification.\par


\noindent\textit{Current registry lookup labels:} No additional label declared.\par

% Stable term ID: extra-correction-and-its-scoring-rule
\begin{samepage}
\subsection*{Extra correction and its scoring rule}
\noindent A model can add a correction to an existing controller. A reward scores the result; command-change rate describes how quickly the input changes.\par
\end{samepage}
\begin{samepage}
\noindent The formal names include residual policy, reinforcement learning and actuator slew. A simulator can score tracking or task completion without proving biological learning or correct hidden neural activity.\par\nopagebreak[4]

\end{samepage}


\noindent\textit{Current registry lookup labels:} Residual policy, reward and actuator slew; residual policy; reinforcement learning; RL; verifiable reward; actuator slew; task completion\par

% Stable term ID: file-fingerprints
\begin{samepage}
\subsection*{File fingerprints}
\noindent A checksum or hash is calculated from file bytes to detect changes or match a stated reference.\par
\end{samepage}
\begin{samepage}
\noindent The named methods here include MD5 and SHA-256. A match checks bytes against that reference; it does not prove scientific truth or verify an entire file when only a range was acquired.\par\nopagebreak[4]

\end{samepage}


\noindent\textit{Current registry lookup labels:} Checksums and hashes; checksum; hash; MD5; SHA-256; source checksum; archive hash\par

% Stable term ID: force-in-the-modeled-muscle-and-tendon
\begin{samepage}
\subsection*{Force in the modeled muscle and tendon}
\noindent The force produced by the computer model for a muscle-tendon unit, or a force measured by a stated physical method.\par
\end{samepage}
\begin{samepage}
\noindent The native test reads forces from the running simulator rather than an old lookup table. Those values are still simulated rather than isolated forces recorded in the source animal.\par\nopagebreak[4]

\end{samepage}


\noindent\textit{Current registry lookup labels:} Tendon force; muscle force; isolated muscle force\par

% Stable term ID: form-function-action
\begin{samepage}
\subsection*{Form \(\rightarrow\) function \(\rightarrow\) action}
\noindent The programme narrative follows structural constraints, the context-dependent processes they support, and independently measured actions or outcomes.\par
\end{samepage}
\noindent The three words organize the explanation; every molecular and neural level has its own structure and dynamics. Form informs and constrains function but does not determine the full process alone. Feedback returns outcomes to the system; this is not a one-way causal chain or a guaranteed one-to-one mapping.\par
\noindent\textit{Organizational level:} Declared system or interface\par
\noindent\textit{Computational granularity:} Declared process/interface/model\par
\begin{samepage}
\noindent\textit{Statistical aggregation:} Validation domain stated independently of physical scale\par\nopagebreak[4]

\end{samepage}


\noindent\textit{Current registry lookup labels:} Form to function to action\par

% Stable term ID: functional-electrical-stimulation
\begin{samepage}
\subsection*{Functional electrical stimulation}
\noindent Electrical stimulation coordinated to produce a specified functional movement or task.\par
\end{samepage}
\noindent FES is the appropriate functional-actuation label for the proposed later assisted-body study. It requires separately qualified devices, control/safety limits and specialist-supervised research approval. This registry contains no stimulation settings or self-experiment instructions.\par
\noindent\textit{Organizational level:} Declared game/model/interface role\par
\noindent\textit{Computational granularity:} Declared task-specific model or interface\par
\begin{samepage}
\noindent\textit{Statistical aggregation:} Validation scope declared per experiment; no physical-scale inference\par\nopagebreak[4]

\end{samepage}
\noindent\textit{Source:} Restoring cortical control of functional movement in a human with quadriplegia\par
\noindent\textit{Source identifier:} bouton2016\par
\noindent\textit{DOI:} 10.1038/nature17435\par
\noindent\url{https://www.nature.com/articles/nature17435}\par
\noindent Research precedent linking implanted cortical recording to a custom neuromuscular stimulation sleeve; not an EEG/OpenBCI or elite-skill result\par


\noindent\textit{Current registry lookup labels:} FES\par

% Stable term ID: further-training-for-the-target-task
\begin{samepage}
\subsection*{Further training for the target task}
\noindent Continuing model training so it predicts the desired target, using the relevant inputs and labels.\par
\end{samepage}
\begin{samepage}
\noindent The paper calls this fine-tuning. Changing software weights is not evidence that an animal or person learned a biological skill.\par\nopagebreak[4]

\end{samepage}


\noindent\textit{Current registry lookup labels:} Fine-tuning; paired-output fine-tuning\par

% Stable term ID: future-testing-and-analysis-of-existing-data
\begin{samepage}
\subsection*{Future testing and analysis of existing data}
\noindent A future test checks later unseen outcomes under a fixed plan; an existing-data test analyzes outcomes already recorded.\par
\end{samepage}
\begin{samepage}
\noindent The second source file was chosen before inspection, but its recording day is earlier than the original model’s day. Choosing a file in advance does not make it future biological deployment.\par\nopagebreak[4]

\end{samepage}


\noindent\textit{Current registry lookup labels:} Prospective and retrospective evaluation; prospective; retrospective; prospective biological deployment\par

% Stable term ID: gap-between-data-portions
\begin{samepage}
\subsection*{Gap between data portions}
\noindent Leaving out an interval around a split boundary so overlapping windows do not cross into another portion.\par
\end{samepage}
\begin{samepage}
\noindent The continuous pipeline uses a one-second gap and checks the input and target support. This protects the software split, not every assumption about the original measurement.\par\nopagebreak[4]

\end{samepage}


\noindent\textit{Current registry lookup labels:} Partition embargo; embargo; boundary embargo\par

% Stable term ID: genotype
% Canonical status: "established_scientific_reference"
\subsection*{Genotype}
\noindent A genetic description can record which versions of a DNA sequence an organism carries. Genotype names that genetic make-up, whether the description covers a selected part of the DNA or a wider set. The scope must be stated; the genetic description alone does not tell us every observable trait.\par
\noindent The genetic constitution being described for an organism, with the genetic scope of the description stated.\par
\noindent Genotype is an established scientific term. Its description is not a complete prediction of the organism’s observed traits; phenotype also depends on environment and development. The three labels in the user’s organism-level example are not three obligatory sequential stages or three fixed individual types.\par


\noindent\textit{Current registry lookup labels:} No additional label declared.\par

% Stable term ID: hallucination
\begin{samepage}
\subsection*{Hallucination}
\noindent A model can predict a movement correctly while inventing the biological mechanism it claims produced it.\par
\end{samepage}
\noindent Invented or ungrounded model content presented as if it corresponded to reality. In programme claims, this includes fabricated biological structure, dynamics or ability presented as established biological fidelity.\par
\noindent Research on AlphaFold3 uses the term for generative-model errors; this programme also applies it to unsupported biological attributions. A disclosed approximation, hypothesis, counterfactual or unvalidated candidate is a different kind of statement, and missing evidence alone does not prove it false. Numerical error concerns the calculation; objective misspecification concerns what is being optimized; causal attribution error concerns the mechanism being credited. Those differences matter when deciding what a result supports. An ungrounded generated answer can happen to be correct without validating its claimed biological mechanism or reliability elsewhere. Informed predictions therefore still require independent performance evaluation. The rolling-dice analogy describes reliance on unverified claims, not literal randomness or equal odds. Calling every approximation a hallucination, or every connectome or simulation fictitious, would erase the distinction.\par
\noindent\textit{Scope:} usage: Model-generated biological content and claims of biological fidelity.\par
\begin{samepage}
\noindent\textit{Computational granularity:} Model output and its interpretation\par\nopagebreak[4]

\end{samepage}


\noindent\textit{Current registry lookup labels:} hallucination; model hallucination\par

% Stable term ID: help-during-current-performance
\subsection*{Help during current performance}
\noindent Help from a controller, replay, trainer or interface that improves performance while it is present.\par
\noindent For example, a trainee can move successfully while following trainer commands. This does not show independent learning or tissue regeneration.\par

\noindent\textit{Current registry lookup labels:} Assistance; compensation; assisted performance\par

% Stable term ID: how-a-muscle-s-pull-turns-a-joint
\begin{samepage}
\subsection*{How a muscle’s pull turns a joint}
\noindent To relate a muscle’s pulling force to joint turning, the model uses its geometry. Consistent work calculations also relate those forces to small coordinate changes.\par
\end{samepage}
\begin{samepage}
\noindent The formal terms are moment arm and virtual work. Their numerical checks do not measure biological muscle strength; the human controller uses them to account for dependent shoulder movements.\par\nopagebreak[4]

\end{samepage}


\noindent\textit{Current registry lookup labels:} Moment arm and virtual work; moment arm; virtual-work projection; dependent shoulder-girdle load\par

% Stable term ID: how-many-labeled-examples-were-supplied
\begin{samepage}
\subsection*{How many labeled examples were supplied}
\noindent The labeled examples given to train or choose a model, including examples from prior recordings.\par
\end{samepage}
\begin{samepage}
\noindent The first 5\% counts new training labels only. Validation labels, earlier-source labels and extra neural samples are additional information.\par\nopagebreak[4]

\end{samepage}


\noindent\textit{Current registry lookup labels:} Label budget and supervision; label budget; paired supervision; adaptation labels; validation labels; source-task supervision\par

% Stable term ID: how-scores-from-several-outputs-are-combined
\begin{samepage}
\subsection*{How scores from several outputs are combined}
\noindent Scores can be averaged across individual outputs or first combined within defined groups. In this entry, macro describes the averaging and weighting rule, not the size of a body part, the scope of a nervous system, or the level of detail in a computer model.\par
\end{samepage}
\begin{samepage}
\noindent The sensory result gives equal weight to five physical output groups. This avoids letting a larger group dominate but does not make their different physical units interchangeable.\par\nopagebreak[4]

\end{samepage}


\noindent\textit{Current registry lookup labels:} Per-output and group-macro scoring; per-output; group-macro; macro R²\par

% Stable term ID: how-uncertain-the-result-remains
\begin{samepage}
\subsection*{How uncertain the result remains}
\noindent The remaining doubt about a prediction, including why it could be wrong or change.\par
\end{samepage}
\begin{samepage}
\noindent Noise in the instrument, errors in the model, changing conditions and missing observations are different causes. The current ranges do not separate all of them.\par\nopagebreak[4]

\end{samepage}


\noindent\textit{Current registry lookup labels:} Prediction uncertainty; uncertainty; model uncertainty; measurement noise; drift\par

% Stable term ID: human-arm-model-used-for-playback
\begin{samepage}
\subsection*{Human arm model used for playback}
\noindent MyoArm is the human arm model used for the displayed muscle-driven replay.\par
\end{samepage}
\begin{samepage}
\noindent The replay uses 63 modeled muscles and two mapped joint goals. Those quantities are simulated, and the macaque velocity mapping does not prove neural or skill transfer into a human body.\par\nopagebreak[4]

\end{samepage}


\noindent\textit{Current registry lookup labels:} MyoArm; human MyoArm; MyoArm right\par

% Stable term ID: internal-musculoskeletal-model
\begin{samepage}
\subsection*{Internal musculoskeletal model}
\noindent The internal anatomy part of the body model: specified bones, joints and soft tissues, assembled from appropriate registered anatomical measurements.\par
\end{samepage}
\noindent Proposed complementary X-ray/CT and MRI inform different structures. Anatomy does not directly measure activation, force, material properties or neural function. Wētā FX provides visual/simulation inspiration, not a verified identical acquisition pipeline.\par


\noindent\textit{Current registry lookup labels:} No additional label declared.\par

% Stable term ID: interoception
\begin{samepage}
\subsection*{Interoception}
\noindent Information about the body’s internal condition, such as signals related to heartbeat or breathing.\par
\end{samepage}
\begin{samepage}
\noindent\textit{Scope:} Sensing the body’s internal condition differs from sensing a limb’s position and movement. Internal sensing remains a research objective; this prototype does not provide a fully measured body-wide internal signal.\par\nopagebreak[4]

\end{samepage}


\noindent\textit{Current registry lookup labels:} Sensing the body from inside; interoception; interoceptive signal; internal sensations\par

% Stable term ID: interoceptome
% Canonical status: "user_requested_proposed_data_organization"
\subsection*{Interoceptome}
\noindent A heartbeat signal and a breathing signal describe different parts of the body’s internal condition. An Interoceptome would organize the selected observations and their relationships, retaining how, where and when each was recorded.\par
\noindent A proposed organization of specified observations of internal-body sensing and the relationships being studied among them.\par
\noindent This is a data-organization proposal within Interoceptomics. Its coverage must name the organism, body systems, measurements, timing and task. Measured signals, calculated quantities and model estimates remain distinguishable. Touch and limb-position sensing are not relabelled as interoception, and the suffix does not establish a complete internal-body map or an individual biological type.\par


\noindent\textit{Current registry lookup labels:} No additional label declared.\par

% Stable term ID: interoceptomics
\begin{samepage}
\subsection*{Interoceptomics}
\noindent How does the body sense a heartbeat or its breathing, and how are the nerve connections involved arranged? Interoceptomics brings those questions together in a proposed study, relating observations of internal-body sensing to maps of the connections carrying or processing that information.\par
\end{samepage}
\noindent A proposed field combining interoception and connectomics to study internal-body sensing and its relationships with neural connectivity across scales.\par
\noindent The programme uses Interoceptomics as an umbrella for this proposed research. Giving the work a name does not establish a validated comprehensive map or a recognized external standard.\par
\noindent\textit{Scope:} macro: Whole-body, organ and network relationships relevant to internal-body sensing; micro: Local signals, cells, circuits and interfaces relevant to those relationships.\par
\begin{samepage}
\noindent\textit{Organizational level:} whole system; network; organ; local circuit; cell; signal/feature\par\nopagebreak[4]
\nopagebreak[4]
\end{samepage}


\noindent\textit{Current registry lookup labels:} interoceptomics\par

% Stable term ID: keeping-each-source-separate
\begin{samepage}
\subsection*{Keeping each source separate}
\noindent Keeping track of which recording, channel, time and physical setting a value belongs to.\par
\end{samepage}
\begin{samepage}
\noindent A channel numbered 12 in one day is not automatically the same cell as channel 12 on another day. A software label also cannot prove where supplied data physically came from.\par\nopagebreak[4]

\end{samepage}


\noindent\textit{Current registry lookup labels:} Source conditioning; source-conditioned; source-specific; recording-specific\par

% Stable term ID: keeping-each-trial-together
\begin{samepage}
\subsection*{Keeping each trial together}
\noindent Assigning an entire trial to one data portion instead of dividing it among portions.\par
\end{samepage}
\begin{samepage}
\noindent This prevents pieces of the same trial from being shared between the relevant splits. It does not remove prior inspection or other source limitations.\par\nopagebreak[4]

\end{samepage}


\noindent\textit{Current registry lookup labels:} Whole-trial separation; whole-trial-separated; successful trial; trial partition\par

% Stable term ID: keeping-up-now-and-computing-for-later
\begin{samepage}
\subsection*{Keeping up now and computing for later}
\noindent Real-time computation keeps pace with incoming data; offline computation can take longer and be played back later.\par
\end{samepage}
\begin{samepage}
\noindent Native planning takes about 4.52 seconds for a 20 ms simulated step and is offline. Smooth cached playback does not prove live planning or biological control.\par\nopagebreak[4]

\end{samepage}


\noindent\textit{Current registry lookup labels:} Real-time and offline execution; real-time; offline; paced playback; cached reference\par

% Stable term ID: kind-of-measurement
\begin{samepage}
\subsection*{Kind of measurement}
\noindent What a method records, such as electrical spikes, scalp voltage, muscle electrical activity or a brain-scan signal.\par
\end{samepage}
\begin{samepage}
\noindent These signals measure different things. They cannot be treated as interchangeable numbers merely because all of them relate to the nervous system.\par\nopagebreak[4]

\end{samepage}


\noindent\textit{Current registry lookup labels:} Signal modality; modality; observation process\par

% Stable term ID: large-movement-model
\subsection*{Large Movement Model}
\noindent The shorthand ``NeuroTranslate with MoCap MyoSim = Large Movement Model'' names the proposed combination of NeuroTranslate, MoCap ReComp and MyoSim Replay. It brings a learned mapping, measured movement and virtual mechanics into one architecture. The equals sign is not a mathematical identity or evidence of a trained model. EEG, EMG, measured movement, virtual mechanics and the trainee retain their distinct roles.\par
\noindent The proposed movement-focused learned model that brings neural structure and functional observations together with internal musculoskeletal anatomy and independently measured external movement in a declared task.\par
\noindent The programme model-family name, analogous in ambition to large language model. It does not replace accurate technical architecture descriptions or mean that every DNN is a movement model. Scans and capture alone do not train a model. Parameter count, scale, generalization and nervous-system fidelity require evidence. System 3 thinking remains the partnership framing. In the October 19 shorthand, MoCap MyoSim combines MoCap ReComp and MyoSim Replay with NeuroTranslate. The equals sign names the proposed architecture; it is not a mathematical identity or evidence that a trained model already exists. These three named groups do not remove the distinct EEG, EMG, measured-movement, virtual-mechanics and trainee roles.\par

\noindent\textit{Current registry lookup labels:} large movement models; LMM; Large movement model\par

% Stable term ID: lasting-change-in-the-nervous-system
\begin{samepage}
\subsection*{Lasting change in the nervous system}
\noindent Changes in the nervous system that can contribute to learning or adaptation.\par
\end{samepage}
\begin{samepage}
\noindent Updating software weights or controller gains is a different event. The practice interface in the computer model does not measure synaptic change in a living organism.\par\nopagebreak[4]

\end{samepage}


\noindent\textit{Current registry lookup labels:} Biological plasticity; synaptic plasticity; biological learning\par

% Stable term ID: learning-before-the-final-task
\begin{samepage}
\subsection*{Learning before the final task}
\noindent Training first on another objective, such as predicting a later neural count, before training for a paired movement target.\par
\end{samepage}
\begin{samepage}
\noindent The paper calls this pretraining. Earlier source-task labels remain part of the information the model received, even when only new-setting labels are counted in a small budget.\par\nopagebreak[4]

\end{samepage}


\noindent\textit{Current registry lookup labels:} Pretraining; neural next-bin pretraining\par

% Stable term ID: mapping-movement-onto-another-body-model
\begin{samepage}
\subsection*{Mapping movement onto another body model}
\noindent Turn a specified movement quantity, such as fingertip velocity components, into movement goals for a particular computer body. The mapping depends on that body and the chosen task.\par
\end{samepage}
\noindent Mapping specified movement quantities into goals for a particular body model, with the source quantities, target coordinates and transformation declared. This does not by itself reconstruct full posture or demonstrate biological transfer.\par
\begin{samepage}
\noindent The retained source-specific adapter takes two macaque fingertip velocity components as its source and produces human-model shoulder and elbow goals. Their signs carry direction; scalar speed would discard it. These are goals for that model and task, so the mapping establishes neither cross-species skill transfer nor verified biological function.\par\nopagebreak[4]

\end{samepage}


\noindent\textit{Current registry lookup labels:} Retargeting; body adapter; organism-specific body adapter; joint goal\par

% Stable term ID: meme-ification
\begin{samepage}
\subsection*{Meme-jacking}
\noindent Borrow the recognition; earn the explanation. A familiar image should help the audience see the scientific point before the script develops it.\par
\end{samepage}
\noindent Use a recognizable meme or cultural reference to make a scientific point immediately legible, drawing on familiarity already earned by others.\par
\begin{samepage}
\noindent Memification hijacking is the user’s expansion. References may come from any online culture, not only films. Apply one purposeful reference per teaching concept without crowding formal methods or labels. Recognition does not grant reuse rights, endorsement, proven comprehension or scientific validity. Preserve the original Meme-ification history.\par\nopagebreak[4]

\end{samepage}
\noindent Meme-ification/memeification are spelling forms; mimification is retained only as the corrected historical transcription.\par


\noindent\textit{Current registry lookup labels:} Meme-ification; meme-ification; memeification; mimification; meme-jacking; meme jacking; memification hijacking\par

% Stable term ID: mind-meld
\subsection*{Mind-Meld}
\noindent What information could one system give another, and would it help with the task? Mind-Meld names the proposed first Brain2Brain encounter, whose signal and effects would have to be measured.\par
\noindent A proposed first Brain2Brain study of a specified information link between neural systems during a declared task.\par
\noindent The roadmap places this at step 3.1 before later NeuroTransfer studies. The name does not describe fused minds, shared private experience, arbitrary thought reading or demonstrated skill transfer. Specify the source, recipient, recording and feedback interfaces, task, controls and independent outcome. Assistance during a link and learning retained without trainer information need separate tests, with any necessary compensatory bridge held constant. A reported natural connection is a motivating case, not evidence that this engineered link exists.\par

\noindent\textit{Current registry lookup labels:} No additional label declared.\par

% Stable term ID: mind-merge
\subsection*{Mind-Merge}
\noindent An agent could guide practice by giving feedback about a stated task. Mind-Merge names the proposed study of that interaction: what help does the person receive, and what can they do later without it?\par
\noindent A proposed Agent2Brain study of task-specific feedback from a software agent to a person and its measurable effects.\par
\noindent The roadmap places this at step 4.1. Ordinary sensory feedback and direct neural delivery remain different routes with different interfaces, evidence and admission. EEG and EMG record; neither supplies a neural writer. No merged identity, shared mind, retained skill or physical intervention follows from the name. Preserve the reverse observation path, matched practice and delayed/new-task retention tests, removing trainer information while keeping a necessary bridge constant. Physical spinal work remains read-only.\par

\noindent\textit{Current registry lookup labels:} No additional label declared.\par

% Stable term ID: mixed-source-animal-and-added-tissue
\begin{samepage}
\subsection*{Mixed-source animal and added tissue}
\noindent A chimera has biological contributions from different sources. A graft is added tissue; when it crosses species, it is a xenotransplant. An organoid is an organized tissue model.\par
\end{samepage}
\begin{samepage}
\noindent Survival, physical connection, activity and useful function are separate questions. The proposed comparisons do not establish a humanized whole brain or an implemented treatment.\par\nopagebreak[4]

\end{samepage}


\noindent\textit{Current registry lookup labels:} Chimera, graft and organoid; chimera; chimerism; graft; organoid; host–graft; xenotransplant; xenograft; XenoTransplantation\par

% Stable term ID: motion-capture-composite-relay
\begin{samepage}
\subsection*{MoCap ReComp}
\noindent MoCap ReComp is Motion-Capture Relay-Composite: the proposed measurement component that brings views of the same movement into a common physical frame. It should preserve which view supplied each point, so the joined map can be compared with its original observations.\par
\noindent The proposal starts with two dot maps of the same moving body: one from above, one from below. Align their time and position in a common physical coordinate frame, then join them into one motion map while retaining the source of each point.\par
\end{samepage}
\noindent A proposed method that creates two dot maps (3D point clouds) of the same moving body, one captured from above and one from below, aligns them in time and a common physical coordinate frame, then stitches them into one composite motion map while retaining each point’s source.\par
\noindent To compare the same movement from above and below, the proposal must align both streams in time and a common physical coordinate frame while retaining each observation’s source. Bidirectional denotes those opposing measurements, not neural communication or causal direction. LiDAR ranging and infrared optical motion capture are proposed alternatives; marker triangulation and LiDAR point clouds are different measurements, so the actual modality and camera/ranging geometry must be qualified. The maze geometry calls for measurements perpendicular to the floor, with viewing paths from both sides through material intended to be clear. Transmission and material effects need testing; concentrated and distributed arrays remain alternatives. Hardware, coverage, timing, calibration, spatial alignment, occlusion, motion error and latency require validation against real references. The joined map is a derived record, not a direct observation of hidden joints, tissue mechanics, muscle commands or neural function. Neither successful implementation, complete body reconstruction nor material transparency has been established by the name. The single vertical camera illustration remains schematic.\par
\noindent\textit{Geometry:} Declare the top/bottom sensor positions, measured object, coordinate frame, viewing/illumination paths and coverage. Preserve the requested perpendicular-to-maze-floor geometry; document deviations rather than inferring them from an illustration.\par
\noindent\textit{Modality:} Declare LiDAR ranging or infrared optical motion capture, including marker-based/markerless method, camera count, triangulation geometry, material path and actual calibration. No complete reconstruction or physiological writing follows from opposing views.\par
\noindent\textit{Timing:} Record source timestamps, clocks, acquisition cadence, synchronization error and any interpolation, offline processing or end-to-end latency. No real-time capability is assumed.\par
\noindent\textit{Calibration and registration:} Preserve calibration data and transformations placing both streams into a common physical frame; measure registration error and keep original observations and reconstruction steps traceable.\par
\noindent\textit{Limits:} Record overlap, occlusion, missing observations, material effects, uncertainty and failure cases. Opposing measurements alone do not prove complete 3D pose, forces or biological function.\par
\noindent\textit{Physical scale:} Declare measured distance ranges and units for an actual acquisition; no supported numeric range or physical unit assignment is presently available for this proposal.\par
\noindent\textit{Statistical aggregation:} Declare any chosen evaluation grouping/averaging rule separately from physical scale and computational granularity.\par
\noindent\textit{References and acceptance:} Use \hyperref[app:reference-comparison-protocol]{the Reference Comparison Protocol and unfilled record} to identify the real-reference series needed for each acquisition component and the assembled map. Its provenance, rights, organism-stage, physical-unit, modality, clock, calibration, occlusion, contact, independent held-out, error, uncertainty and acceptance fields remain unfilled and HOLD. It supplies no measured series or apparatus validation. The unchanged programme reference criteria and every stronger original requirement govern their respective claims; no acceptance threshold is introduced or relaxed.\par
\noindent\textit{Outputs:} Retain the top dot map, bottom dot map and joined map as separate traceable outputs. Their timestamps and common-frame transformations must be explicit; do not fill unobserved regions as though they were measured.\par
\noindent\textit{Scope:} proposed application: Complementary top/bottom motion measurements of a body moving through the maze; concentrated coordinated platforms and distributed arrays remain alternatives; measurement boundary: Observed distances or geometry remain distinct from registered/stitched composites and any model-estimated pose or musculoskeletal quantity.\par
\noindent\textit{Organizational level:} whole system; signal/feature\par
\begin{samepage}
\noindent\textit{Computational granularity:} top dot map (3D point cloud); bottom dot map (3D point cloud); time alignment and spatial registration; joined composite motion sequence\par\nopagebreak[4]

\end{samepage}
\noindent Historical spellings name this same proposed capture method. Relay here means passing captured measurements onward; the separate agent-coordination Relay protocol does not describe motion capture.\par


\noindent\textit{Current registry lookup labels:} motion capture composite relay; motion-capture composite relay; motion capture composite; bidirectional composite motion capturing using LiDAR; Motion capture composite relay; motion capture relay composite; motion-capture relay composite; bidirectional LiDAR; opposed-view infrared motion capture; bidirectional infrared; motion-capture relay-composite; Motion-Capture Relay-Composite\par

% Stable term ID: model-that-chooses-a-label
\begin{samepage}
\subsection*{Model that chooses a label}
\noindent A statistical model that assigns scores to categories, such as the ten attempted-finger labels.\par
\end{samepage}
\begin{samepage}
\noindent The selected human model uses logistic classification, but its scores do not establish validated confidence in biological truth. Its reported accuracy comes from development data that have already been seen.\par\nopagebreak[4]

\end{samepage}


\noindent\textit{Current registry lookup labels:} Logistic classifier; logistic regression; classifier; classification; logistic\par

% Stable term ID: model-weights-kept-fixed
\begin{samepage}
\subsection*{Model weights kept fixed}
\noindent Keeping the fitted numbers in a model unchanged during a defined comparison.\par
\end{samepage}
\begin{samepage}
\noindent The tested transfer trains a connector while retaining the older main model. Fixed weights make the comparison clear; they do not prove biological sameness across days.\par\nopagebreak[4]

\end{samepage}


\noindent\textit{Current registry lookup labels:} Frozen parameters; frozen trunk; frozen readout; frozen checkpoint\par

% Stable term ID: modeled-route-from-one-muscle-attachment-to-another
\begin{samepage}
\subsection*{Modeled route from one muscle attachment to another}
\noindent The route and length assigned to a muscle and tendon as the model’s joints move.\par
\end{samepage}
\begin{samepage}
\noindent The native arm retains 39 named routes. Their modeled lengths differ from direct measurements of living tissue, muscle commands and force.\par\nopagebreak[4]

\end{samepage}


\noindent\textit{Current registry lookup labels:} Muscle-tendon path; musculotendon length; modeled muscle length; geometric path\par

% Stable term ID: models-that-convert-inputs-or-predict-outputs
\begin{samepage}
\subsection*{Models that convert inputs or predict outputs}
\noindent One model can turn an input into a representation; another can use available inputs to predict a named result.\par
\end{samepage}
\begin{samepage}
\noindent The paper calls these encoders and decoders. Predicting counts from mechanics or mechanics from counts does not by itself show the biological cause of either.\par\nopagebreak[4]

\end{samepage}


\noindent\textit{Current registry lookup labels:} Encoder and decoder; encoding; decoding; readout; decode; encoder; decoder\par

% Stable term ID: movement-brain-area
\begin{samepage}
\subsection*{Movement brain area}
\noindent A brain area involved in movement. The paper uses the names motor cortex, MC and M1 for relevant source recordings.\par
\end{samepage}
\begin{samepage}
\noindent The signals help predict recorded movement-related targets. Knowing the brain area does not identify matching cells across recordings or reveal a whole-body command pathway.\par\nopagebreak[4]

\end{samepage}


\noindent\textit{Current registry lookup labels:} Motor cortex; M1; MC; primary motor cortex\par

% Stable term ID: muscle-electrical-recording
\begin{samepage}
\subsection*{Muscle electrical recording}
\noindent Electrical activity recorded from muscles; the method is called EMG.\par
\end{samepage}
\begin{samepage}
\noindent Muscle electrical activity differs from force and from the computer model’s command and activity states. The trained targets here are rescaled signal summaries.\par\nopagebreak[4]

\end{samepage}


\noindent\textit{Current registry lookup labels:} Electromyography; EMG\par

% Stable term ID: myoception
\subsection*{Myoception}
\noindent Sensing muscle-driven movement and the feedback that comes with it.\par
\noindent The programme’s term for the sensory side of musculoskeletal contraction and movement, including proprioception, tactile feedback and other task-relevant sensations linked to action.\par
\noindent This is the user’s proposed programme scope, not a claim that every relevant signal has been recorded or a new external scientific standard established. Muscle electrical activity, proprioceptive or tactile measurements, neural recordings and simulated feedback remain distinct observations. EMG does not directly record sensory afferents or the full sensory consequences of contraction. Declare each source’s organism, anatomy, modality, task, units, timing, uncertainty and measured, derived or simulated status. The wider neural-data ambition includes brain, spinal and peripheral activity and relevant visceral signals; the current focus is the motor–sensory system. No body-wide coverage, unique hidden neural state or subjective experience is established by this term.\par

\noindent\textit{Current registry lookup labels:} myoception; M-Y-O-ception\par

% Stable term ID: musculoskeletal-simulation-physics-engine
\begin{samepage}
\subsection*{MyoSim Replay}
\noindent MyoSim Replay is the proposed MusculoSkeletal Simulation Physics-Engine. Declared anatomy, mechanical parameters and control inputs would produce movement predictions to compare with independent measurements. The external MyoSim body-model collection is a separate resource. Simulation replay is not a measured biological replay.\par
\noindent The proposed digital model would calculate how muscles and the skeleton move and exert forces, with real measurements used to check its predictions.\par
\end{samepage}
\noindent The proposed physics-based simulation component that models specified muscles, bones, joints and their mechanical interactions, using declared anatomy, parameters, forces and contact rules, with predictions tested against independent movement and other applicable physical measurements.\par
\noindent In the proposed five-component arrangement, declared anatomy, muscle/control inputs and physical context would drive a simulation of movement and forces. The motion-capture relay-composite would provide traceable measured movement to compare with those predictions. That comparison needs matching timing, coordinates, units and calibration, and tests that preserve anatomical and physical constraints, uncertainty, held-out cases, baseline controls, a series of real references and independent acceptance. Movement geometry alone does not measure muscle activation, internal forces, hidden neural state or subjective experience; force, activation and causal claims need additional independent references. Visual effects (VFX) rendering shows the model but cannot validate its mechanics. This remains a proposed simulation component, not a newly demonstrated implementation or external scientific field. Approximate simulation and biological virtualization verified within declared error bounds remain distinct; this component establishes neither universal precision nor whole-body completeness.\par
\noindent\textit{Scope:} programme role: Virtual-body mechanics component; initially constrained hand and foot tasks; inputs and outputs: Declared anatomical parameters, muscle/control inputs and physical context produce simulated movement and mechanical quantities. Independently acquired motion and other necessary observables remain references, not simulator ground truth.\par
\noindent\textit{Organizational level:} whole system; organ\par
\begin{samepage}
\noindent\textit{Computational granularity:} musculoskeletal model; mechanical solver; measurement-to-prediction comparison\par\nopagebreak[4]

\end{samepage}


{\rightskip=0pt plus1fil\relax
\noindent\textit{Current registry lookup labels:} musculoskeletal simulation physics engine; musculoskeletal-simulation physics engine; musculoskeletal simulation; musculoskeletal physics model; musculoskeletal-simulation physics-engine; MusculoSkeletal Simulation Physics-Engine\par
}

% Stable term ID: names-used-for-source-subjects
\begin{samepage}
\subsection*{Names used for source subjects}
\noindent Indy, Han and Monkey L label macaque sources; N1 labels the human finger source.\par
\end{samepage}
\begin{samepage}
\noindent These labels keep different subjects separate. Similar brain-region or task names do not make them one synchronized recording.\par\nopagebreak[4]

\end{samepage}


\noindent\textit{Current registry lookup labels:} Source subject labels; Indy; Han; Monkey L; N1\par

% Stable term ID: nerves-bringing-body-information-in
\begin{samepage}
\subsection*{Nerves bringing body information in}
\noindent Afferent paths carry body-related information toward the brain and spinal cord, together called the central nervous system.\par
\end{samepage}
\begin{samepage}
\noindent For example, a limb-related sensory signal can travel inward. The source-specific cortical predictions do not measure this whole path or make missing nerve activity observed.\par\nopagebreak[4]

\end{samepage}


\noindent\textit{Current registry lookup labels:} afferent; afferent pathway; body-to-CNS; CNS; central nervous system\par

% Stable term ID: nerves-carrying-commands-out
\begin{samepage}
\subsection*{Nerves carrying commands out}
\noindent Efferent paths carry outputs from the brain and spinal cord toward the body, including muscles and organs.\par
\end{samepage}
\begin{samepage}
\noindent A prediction of movement is not a recording of that outward pathway. The prototype does not supply a synchronized voluntary sequence through spinal cord, nerves and muscles.\par\nopagebreak[4]

\end{samepage}


\noindent\textit{Current registry lookup labels:} efferent; efferent pathway; CNS-to-body\par

% Stable term ID: neural-decoding
\begin{samepage}
\subsection*{Neural decoding}
\noindent Estimating declared information such as intended action from recorded neural observations using a specified model.\par
\end{samepage}
\noindent Decode accuracy, channels, source identity, task and held-out domain must be demonstrated. Brain-only control requires exclusion of covert muscle, eye and motion signals; decoding does not write signals into the brain.\par
\noindent\textit{Organizational level:} Declared game/model/interface role\par
\noindent\textit{Computational granularity:} Declared task-specific model or interface\par
\noindent\textit{Statistical aggregation:} Validation scope declared per experiment; no physical-scale inference\par
\noindent\textit{Source:} Restoring cortical control of functional movement in a human with quadriplegia\par
\noindent\textit{Source identifier:} bouton2016\par
\noindent\textit{DOI:} 10.1038/nature17435\par
\noindent\url{https://www.nature.com/articles/nature17435}\par
\noindent Research precedent linking implanted cortical recording to a custom neuromuscular stimulation sleeve; not an EEG/OpenBCI or elite-skill result\par
\noindent\textit{Source:} A brain-computer interface that evokes tactile sensations improves robotic arm control\par
\noindent\textit{Source identifier:} flesher2021\par
\noindent\textit{DOI:} 10.1126/science.abd0380\par
\noindent\url{https://pubmed.ncbi.nlm.nih.gov/34016775/}\par
\noindent Research precedent for motor decoding plus sensory feedback through intracortical microstimulation; invasive and not a consumer EEG capability\par


\noindent\textit{Current registry lookup labels:} No additional label declared.\par

% Stable term ID: neural-encoding
\begin{samepage}
\subsection*{Neural encoding}
\noindent Mapping declared information to a representation in neural activity; in a sensory-interface design, constructing candidate stimulation patterns for specified feedback.\par
\end{samepage}
\noindent A computational encoder is distinct from a qualified stimulation device and evidence that a participant receives the intended sensation. The term does not mean arbitrary thought writing or skill upload.\par
\noindent\textit{Organizational level:} Declared game/model/interface role\par
\noindent\textit{Computational granularity:} Declared task-specific model or interface\par
\noindent\textit{Statistical aggregation:} Validation scope declared per experiment; no physical-scale inference\par
\noindent\textit{Source:} A brain-computer interface that evokes tactile sensations improves robotic arm control\par
\noindent\textit{Source identifier:} flesher2021\par
\noindent\textit{DOI:} 10.1126/science.abd0380\par
\noindent\url{https://pubmed.ncbi.nlm.nih.gov/34016775/}\par
\noindent Research precedent for motor decoding plus sensory feedback through intracortical microstimulation; invasive and not a consumer EEG capability\par


\noindent\textit{Current registry lookup labels:} No additional label declared.\par

% Stable term ID: playing-back-a-recording
\begin{samepage}
\subsection*{Neural Replay}
\noindent Revisiting a stored sequence of neural samples with its stated timing and channel order.\par
\end{samepage}
\begin{samepage}
\noindent This playback is not a live recording or instructions for stimulating a nervous system. A displayed arm movement does not show what the source animal’s muscles actually did.\par\nopagebreak[4]

\end{samepage}


\noindent\textit{Current registry lookup labels:} Neural replay; recorded replay; archived replay; cached playback; Playing back a recording\par

% Stable term ID: neuroagent-tutor
\subsection*{NeuroAgent Trainer}
\noindent A software agent would guide a trainee during a stated task.\par
\noindent A proposed software-agent trainer role in a declared Agent2Brain NeuroTransfer study.\par
\noindent Identify the actual software model and feedback route. Computed activations/events do not make a biological trainee digital. Illustrative connections do not establish direct neural writing, successful transfer or retained learning. NeuroAgent Tutor is a historical label; NeuroAgent Trainer is the current name. It describes a proposed software-agent trainer role, without selecting the agent or qualifying its feedback interface.\par
\noindent\textit{Organizational level:} whole system; signal/feature\par

\noindent\textit{Current registry lookup labels:} NeuroAgent Tutor\par

% Stable term ID: neuroarcher
\begin{samepage}
\subsection*{NeuroArcher}
\noindent A proposed archery game that studies coordinated use of both hands and arms through physical motion control, brain-signal control of a digital body, and a separately admitted later assisted-body experiment.\par
\end{samepage}
\noindent The near-term prop is a toy bow with small suction-cup tips and an over-arm arrow pouch. Target/creature choice is unresolved. Simulated performance does not establish muscle stimulation or retained biological archery skill.\par
\noindent\textit{Organizational level:} Declared game/model/interface role\par
\noindent\textit{Computational granularity:} Declared task-specific model or interface\par
\begin{samepage}
\noindent\textit{Statistical aggregation:} Validation scope declared per experiment; no physical-scale inference\par\nopagebreak[4]

\end{samepage}


\noindent\textit{Current registry lookup labels:} Neuro Archer\par

% Stable term ID: neuroboxer
\begin{samepage}
\subsection*{NeuroBoxer}
\noindent A proposed boxing game that compares arm-motion control, brain-signal control of a digital body, and a separately admitted later assisted-body experiment.\par
\end{samepage}
\noindent Prime-era Mike Tyson is the requested virtual opponent and meme reference. A game, recorded prediction, stimulation-assisted motion and retained biological skill learning need different evidence. No actual opponent, head contact, finished game or stimulation capability is claimed.\par
\noindent\textit{Organizational level:} Declared game/model/interface role\par
\noindent\textit{Computational granularity:} Declared task-specific model or interface\par
\begin{samepage}
\noindent\textit{Statistical aggregation:} Validation scope declared per experiment; no physical-scale inference\par\nopagebreak[4]

\end{samepage}


\noindent\textit{Current registry lookup labels:} Neuro Boxer\par

% Stable term ID: neuroconnectome
\begin{samepage}
\subsection*{NeuroConnectome}
\noindent A map of specified neural connections used to constrain a model of the nervous system.\par
\end{samepage}
\begin{samepage}
\noindent Connectome is the established term. The user’s NeuroConnectome label does not establish complete coverage or identify functional dynamics, state, plasticity, sensory context or unique thought from wiring alone. Declare species, anatomical coverage, resolution and provenance.\par\nopagebreak[4]

\end{samepage}


\noindent\textit{Current registry lookup labels:} Neuroconnectome\par

% Stable term ID: neurohandler
% Canonical status: "user_requested_proposed_programme_role"
\subsection*{NeuroHandler}
\noindent How well can a hand carry out the chosen task, and what information helps it do so? NeuroHandler is the proposed manual-dexterity experiment, with the hand’s measured action providing the result to explain.\par
\noindent A proposed manual-dexterity experiment studying specified neural/body observations, task control and independently measured outcomes.\par
\noindent The earlier proposed manual tasks remain: rapid hand rubbing, the requested knuckle gesture and the interlaced-hand stretch. NeuroHandler names the manual-dexterity experiment; it does not establish a finished apparatus or successful model. Define the action, available inputs, physical or simulated setting, comparison and outcome before evaluation. Keep motion control, neural prediction, assistance and retained learning separate; any physical output requires separate qualification and admission. No dexterity gain, clinical benefit or intervention procedure is established.\par


\noindent\textit{Current registry lookup labels:} No additional label declared.\par

% Stable term ID: neurojuggler
% Canonical status: "user_requested_proposed_programme_role"
\subsection*{NeuroJuggler}
\noindent Moving each limb successfully does not settle whether the limbs work together. NeuroJuggler is the proposed coordination experiment: define the joint task, then measure the relationship between the limb actions and the result.\par
\noindent A proposed multi-limb coordination experiment studying how specified limb actions are coordinated in a declared task.\par
\noindent The earlier coordination proposal remains a hot-potato transfer between the hands, followed preferably by two-ball juggling. Those are proposed tasks; the final apparatus and qualified measurement setup remain unestablished. The label supplies no successful performance or implemented neural interface. Declare the task, limbs, inputs, timing, body model or physical setting and independent outcomes. Separate coordination from isolated-limb performance, simulated success from biological function, and current assistance from retained learning. Physical output, safety, rights and held-out tests keep their own admission requirements.\par


\noindent\textit{Current registry lookup labels:} No additional label declared.\par

% Stable term ID: neuromuscular-electrical-stimulation
\begin{samepage}
\subsection*{Neuromuscular electrical stimulation}
\noindent Electrical stimulation intended to elicit muscle contraction through the relevant neuromuscular pathways.\par
\end{samepage}
\noindent NMES is a writing/actuation intervention, whereas EMG records muscle activity. EMG electrodes or an OpenBCI recording board do not by themselves provide qualified stimulation. Neither contraction nor imposed movement demonstrates retained skill.\par
\noindent\textit{Organizational level:} Declared game/model/interface role\par
\noindent\textit{Computational granularity:} Declared task-specific model or interface\par
\begin{samepage}
\noindent\textit{Statistical aggregation:} Validation scope declared per experiment; no physical-scale inference\par\nopagebreak[4]

\end{samepage}
\noindent\textit{Source:} Restoring cortical control of functional movement in a human with quadriplegia\par
\noindent\textit{Source identifier:} bouton2016\par
\noindent\textit{DOI:} 10.1038/nature17435\par
\noindent\url{https://www.nature.com/articles/nature17435}\par
\noindent Research precedent linking implanted cortical recording to a custom neuromuscular stimulation sleeve; not an EEG/OpenBCI or elite-skill result\par


\noindent\textit{Current registry lookup labels:} NMES; Electrical muscle stimulation; EMS\par

% Stable term ID: neuroreplay
\subsection*{NeuroReplay}
\noindent What changes when a trainee receives replay-based guidance? NeuroReplay names that proposed study: compare performance while guidance is available, then test what remains when the guidance ends.\par
\noindent A proposed Stage 0 research direction for studying replay-guided sensorimotor feedback and its effects on current performance and later learning.\par
\noindent This is the new roadmap label for steps 0.2–0.5, with a derived working scope; it does not rename Neural Replay, the existing stored-sample playback entry playing-back-a-recording. A stored sequence can supply a comparison, but playback alone establishes neither delivery to a nervous system nor biological learning. Fix the task, signal, feedback route, timing and equal-practice comparisons; distinguish matched, stale or unmatched replay, no-link and inactive or sham controls. Test delayed retention and new tasks after removing trainer information while keeping any necessary compensatory bridge constant. The questioned antechinus age ranges are not established developmental or intervention windows. No physical writer, procedure or present result is supplied by this label.\par

\noindent\textit{Current registry lookup labels:} No additional label declared.\par

% Stable term ID: neuroscience-recording-files
\begin{samepage}
\subsection*{Neuroscience recording files}
\noindent NWB stores neuroscience data and descriptions. In the paired source files here, SET holds recording descriptions and FDT holds binary samples.\par
\end{samepage}
\begin{samepage}
\noindent A stale unit in the NWB muscle source must not override verified processing. SET headers match source hashes; FDT previews have checked local byte ranges rather than full verification of each upstream file.\par\nopagebreak[4]

\end{samepage}


\noindent\textit{Current registry lookup labels:} Neurodata Without Borders file; NWB; Recording header and sample file; SET and FDT files; SET; FDT; EEGLAB header; byte-range preview\par

% Stable term ID: neurotransfer
\begin{samepage}
\subsection*{NeuroTransfer}
\noindent Specify what would move between which systems, then test what the recipient can do with and without assistance.\par
\end{samepage}
\noindent The proposed transfer of specified task-relevant information or a mapping between declared systems, with assistance and retained learning evaluated separately.\par
\begin{samepage}
\noindent The label does not establish a successful biological transfer or retained skill. A weight copy, recorded replay, guided movement and independently retained learning answer different questions. Each proposed interface needs its own source, target, task, controls, rights and qualified outcome.\par\nopagebreak[4]

\end{samepage}


\noindent\textit{Current registry lookup labels:} neurotransfer\par

% Stable term ID: neurotranslate
\begin{samepage}
\subsection*{NeuroTranslate}
\noindent Brain2Body Neural Activity is the long-form component label paired with NeuroTranslate in this proposed architecture. Recorded activity is an observation. NeuroTranslate names the mapping operation that could relate specified observations to movement, or reconstruct specified recorded features in the reverse direction. The observation, learned relationship and predicted value remain different things.\par
\noindent State the observations, the target and the task, then test how well the mapping works.\par
\end{samepage}
\noindent The proposed operation of mapping specified neural and body observations between declared inputs and targets in a stated task.\par
\begin{samepage}
\noindent A model could predict movement from specified observations, or reconstruct specified recorded features in the reverse direction. Each direction needs its own test in the stated task. Neither establishes a unique hidden state, sensory delivery or biological skill transfer. NeuroTranslate names this proposed mapping operation; NeuroTranslation names the field and NeuroTranslator the implemented software.\par\nopagebreak[4]

\end{samepage}


\noindent\textit{Current registry lookup labels:} neurotranslate; Brain2Body Neural Activity\par

% Stable term ID: neurotranslation
\begin{samepage}
\subsection*{NeuroTranslation}
\noindent How could a particular neural recording be related to a defined movement task? NeuroTranslation is the proposed study of mapping nerve signals and the patterns that describe them between clearly defined situations like these. Every mapping must state where its inputs come from, what its outputs mean, and how it is tested.\par
\end{samepage}
\noindent A proposed field for translating or mapping neural signals and representations between explicitly defined source and target domains, with stated context, scale, tests and limitations.\par
\noindent Each connection must be assessed against its named source and stated test. The project has not demonstrated a universal brain language or an uploaded skill.\par
\noindent\textit{Scope:} macro: Relationships between neural systems, body systems and defined tasks or organisms; micro: Mappings involving identified local signals, channels, features and circuit-level representations.\par
\begin{samepage}
\noindent\textit{Organizational level:} whole system; network; organ; local circuit; signal/feature\par\nopagebreak[4]
\nopagebreak[4]
\end{samepage}


\noindent\textit{Current registry lookup labels:} neurotranslation; neural translation; Connecting neural recordings to defined tasks; Neurotranslation\par

% Stable term ID: source-conditioned-neurotranslation-architecture
\begin{samepage}
\subsection*{NeuroTranslation Network}
\noindent A plan for connecting more neural and body components while keeping the origin and meaning of each value clear.\par
\end{samepage}
\begin{samepage}
\noindent The proposed network would connect neural and body components while carrying the origin and meaning of each value through the mapping. Each link could then be examined in its own recording context. Body-wide translation and biological teaching remain unproved.\par\nopagebreak[4]

\end{samepage}


\noindent\textit{Current registry lookup labels:} Proposed larger design; Source-Conditioned Neurotranslation Architecture; SCNA; Neurotranslation Network\par

% Stable term ID: neurotranslation-engine
\begin{samepage}
\subsection*{NeuroTranslator}
\noindent The software that joins particular recordings, prediction models, playback and computer-body checks.\par
\end{samepage}
\begin{samepage}
\noindent Predictions and playback can be checked against the particular recordings they use. That is the scope of the implemented NeuroTranslator; neither a universal brain language nor skill upload has been demonstrated.\par\nopagebreak[4]

\end{samepage}


\noindent\textit{Current registry lookup labels:} Working translation software; Neurotranslation Engine; NTE; Neurotranslator\par

% Stable term ID: neurotwin-tutor
\subsection*{NeuroTwin Trainer}
\noindent A biological trainer would guide a trainee during a stated task.\par
\noindent A proposed biological trainer role in a declared Brain2Brain NeuroTransfer study.\par
\noindent The requested name alone establishes neither genetic twins, a validated digital twin, communication nor retained learning. Biological electrical activity is electrochemical and can include graded voltages and spikes; identify any actual measurement, otherwise keep the activity illustrative. NeuroTwin Tutor is a historical label; NeuroTwin Trainer is the current name. These names belong to a proposed programme view. They establish neither an implemented interface nor a demonstrated learning result.\par
\noindent\textit{Organizational level:} whole system; signal/feature\par

\noindent\textit{Current registry lookup labels:} NeuroTwin Tutor\par

% Stable term ID: numerical-arrays-saved-arrays-and-prediction
\begin{samepage}
\subsection*{Numerical arrays, saved arrays and prediction}
\noindent NumPy provides tools for numerical arrays. NPZ saves arrays; inference means running a trained model to produce an output.\par
\end{samepage}
\begin{samepage}
\noindent Portable execution still needs the correct processing and source identity. Matching saved-model outputs does not establish success on new organisms or biological skills.\par\nopagebreak[4]

\end{samepage}


\noindent\textit{Current registry lookup labels:} NumPy inference and NPZ archive; NumPy; NPZ; portable inference; inference\par

% Stable term ID: numerical-checks-on-allowed-movement
\begin{samepage}
\subsection*{Numerical checks on allowed movement}
\noindent Checks that the computer model stays defined and within its stated joint limits.\par
\end{samepage}
\begin{samepage}
\noindent Joint-limit penetration means going past an angular bound, not entering tissue. The 0.001 rad tolerance is a numerical rule, not a certificate of correct neural commands.\par\nopagebreak[4]

\end{samepage}


\noindent\textit{Current registry lookup labels:} Mechanical feasibility and anatomical tolerance; anatomical limit; limit penetration; finite state; joint range; numerical tolerance\par

% Stable term ID: numerical-software-and-its-data-file
\begin{samepage}
\subsection*{Numerical software and its data file}
\noindent MATLAB is numerical-computing software, and MAT is a format for its structured data.\par
\end{samepage}
\begin{samepage}
\noindent The second Indy recording uses the original MAT source. The format does not prove matching cells, past-only source processing or complete observation.\par\nopagebreak[4]

\end{samepage}


\noindent\textit{Current registry lookup labels:} MATLAB and MAT file; MATLAB; MAT\par

% Stable term ID: observe-it-connect-it-keep-it-useful
\begin{samepage}
\begin{samepage}
\subsection*{Observe it, connect it, keep it useful}
\noindent Observation burden, prediction on excluded cases and lasting task benefit need separate measurements. A clear picture does not answer those experimental questions.\par\nopagebreak[4]

\end{samepage}
\end{samepage}


\noindent\textit{Current registry lookup labels:} Access it, Ground it, Keep it; access; grounding; stability; signal stability\par

% Stable term ID: one-multiplier-for-modeled-lengths
\begin{samepage}
\subsection*{One multiplier for modeled lengths}
\noindent A single number chosen to reduce differences between the model’s muscle lengths and the source-derived length references.\par
\end{samepage}
\begin{samepage}
\noindent The number is fitted on training poses and checked on distinct observed poses. It is not complete body calibration; no movement-control experiment is run on the scaled candidate.\par\nopagebreak[4]

\end{samepage}


\noindent\textit{Current registry lookup labels:} Scalar length fit; global length factor; uniform factor; a*; scalar bridge; recording length-fit candidate; SDK-scaled candidate\par

% Stable term ID: one-recording-period
\begin{samepage}
\subsection*{One recording period}
\noindent One experiment or recording period with its own subject, channels, timing and task.\par
\end{samepage}
\begin{samepage}
\noindent Two recording days, or recordings from different animals, do not form one observed nervous system simply because the files are loaded together.\par\nopagebreak[4]

\end{samepage}


\noindent\textit{Current registry lookup labels:} Recording and session; source recording; session; recording family\par

% Stable term ID: open-body-model-collection
\begin{samepage}
\subsection*{Open body-model collection}
\noindent MyoSim supplies computer-body models with muscles and bones, including models used in this project.\par
\end{samepage}
\begin{samepage}
\noindent The collection is earlier work, not a new project invention. Availability of an arm model does not establish matching macaque and human muscles or neural commands.\par\nopagebreak[4]

\end{samepage}


\noindent\textit{Current registry lookup labels:} MyoSim, MyoSuite and MyoArm; MyoSim\par

% Stable term ID: open-muscle-control-task-tools
\begin{samepage}
\subsection*{Open muscle-control task tools}
\noindent MyoSuite provides tasks and tools for experiments with muscle-driven computer-body control.\par
\end{samepage}
\begin{samepage}
\noindent It is earlier software reused here. Success on a simulated task is not proof of biological learning or a recovered neural command.\par\nopagebreak[4]

\end{samepage}


\noindent\textit{Current registry lookup labels:} MyoSuite\par

% Stable term ID: other-published-models-of-neural-recordings
\begin{samepage}
\subsection*{Other published models of neural recordings}
\noindent NDT2 learns patterns across recordings made inside the cortex, the brain’s outer layer. POYO handles differing sets of recorded neural features by representing events and learning relationships between them.\par
\end{samepage}
\begin{samepage}
\noindent The paper cites these as earlier work. It does not claim to have implemented them or use their existence as proof of universal transfer or matching cells.\par\nopagebreak[4]

\end{samepage}


\noindent\textit{Current registry lookup labels:} NDT2; large-scale neural pretraining; Published model for varying recorded populations; POYO; neural-event tokenization; cross-attention bottleneck\par

% Stable term ID: passing-information-and-gaining-a-lasting-skill
\begin{samepage}
\subsection*{Passing information and gaining a lasting skill}
\noindent Receiving useful task information differs from acquiring a durable ability that works independently in new conditions.\par
\end{samepage}
\begin{samepage}
\noindent A short software label budget is not measured biological practice time. Sewing, cube solving and breakdancing remain future examples without matched skill recordings here.\par\nopagebreak[4]

\end{samepage}


\noindent\textit{Current registry lookup labels:} Task information transfer and skill acquisition; task transfer; skill transfer; skill upload; compressed acquisition; learning speed; learning; physical skill learning; skill acquisition\par

% Stable term ID: past-only-inputs-and-processing-that-uses-later-data
\begin{samepage}
\subsection*{Past-only inputs and processing that uses later data}
\noindent The local model uses completed past inputs. Some earlier source processing may still have used later samples.\par
\end{samepage}
\begin{samepage}
\noindent Past-only indexing cannot reverse earlier filtering or rescaling that used the whole recording. It does not certify the full measurement-to-output delay in a living system.\par\nopagebreak[4]

\end{samepage}


\noindent\textit{Current registry lookup labels:} Causal indexing and acausal preprocessing; causal; algorithmic support; acausal filtering; offline interpolation; upstream filtering delay\par

% Stable term ID: peripheral-nerve-electrical-recording
\begin{samepage}
\subsection*{Peripheral nerve electrical recording}
\noindent A method for recording electrical signals associated with peripheral nerves, called ENG.\par
\end{samepage}
\begin{samepage}
\noindent The human preview contains aggregate electrical observations. It does not provide separately identified nerve-cell spike trains or a complete voluntary command sequence.\par\nopagebreak[4]

\end{samepage}


\noindent\textit{Current registry lookup labels:} Electroneurography; ENG; surface ENG; peripheral recording\par

% Stable term ID: phenotype
% Canonical status: "established_scientific_reference"
\subsection*{Phenotype}
\noindent The trait is what must be observed: for example, a specified body characteristic or response in the setting being studied. Phenotype names these observable traits, rather than the genetic description alone.\par
\noindent The observable traits of an organism arising through genotype, environment and development.\par
\noindent State the actual trait, measurement method, developmental and environmental context. A phenotype is not a universal function that can be read uniquely from structure or sequence. Retain the distinction between a measured trait and a model’s estimate.\par


\noindent\textit{Current registry lookup labels:} No additional label declared.\par

% Stable term ID: placing-body-parts-from-joint-positions
\begin{samepage}
\subsection*{Placing body parts from joint positions}
\noindent Calculating where body parts and muscle routes lie when joint positions are given.\par
\end{samepage}
\begin{samepage}
\noindent The formal name is forward kinematics, or FK. It places geometry; it does not calculate muscle-driven motion through time or recover muscle commands.\par\nopagebreak[4]

\end{samepage}


\noindent\textit{Current registry lookup labels:} Forward kinematics; FK; kinematic replay\par

% Stable term ID: planned-commands-corrections-and-constant-input
\begin{samepage}
\subsection*{Planned commands, corrections and constant input}
\noindent A planned command is feedforward; adjusting it from current model information is feedback; a constant muscle input is tonic control.\par
\end{samepage}
\begin{samepage}
\noindent The human correction can change the bounded command goal. Recorded movement stays a scoring reference, and the native correction candidates remain rejected against constant-input comparisons.\par\nopagebreak[4]

\end{samepage}


\noindent\textit{Current registry lookup labels:} Feedforward, feedback and tonic control; feedforward; feedback; corrected controller; tonic excitation; constant tonic\par

% Stable term ID: prediction-score-compared-with-a-mean-value-guess
\begin{samepage}
\subsection*{Prediction score compared with a mean-value guess}
\noindent R² compares squared prediction errors with a constant mean-value reference. Higher values usually indicate better prediction on the stated data.\par
\end{samepage}
\begin{samepage}
\noindent The score can be negative. Different recordings, target processing and ways of combining outputs prevent treating all R² values as directly comparable.\par\nopagebreak[4]

\end{samepage}


\noindent\textit{Current registry lookup labels:} Coefficient of determination; R²; R2; R\textasciicircum{}2; variance explained\par

% Stable term ID: processing-within-spinal-and-peripheral-pathways
\begin{samepage}
\subsection*{Processing within spinal and peripheral pathways}
\noindent Body-related information can be processed locally in spinal and peripheral structures rather than passing through one simple unchanged wire.\par
\end{samepage}
\begin{samepage}
\noindent A cortex-only prediction does not observe those local processes. The separate recordings and computer-body checks do not establish a complete measured brain-to-body causal chain.\par\nopagebreak[4]

\end{samepage}


\noindent\textit{Current registry lookup labels:} local spinal processing; local peripheral processing; spinal/peripheral processing\par

% Stable term ID: programmable-neurobiology
\begin{samepage}
\subsection*{Programmable neurobiology}
\noindent The proposed ability to specify a bounded nervous-system or sensorimotor outcome, select a qualified intervention, predict its effects, measure what happens and adapt within a validated closed loop.\par
\end{samepage}
\noindent\textit{Scope:} Brain, spinal, peripheral, muscle, mechanical and sensory feedback roles need separate evidence. Bidirectional prediction is different from safe sensory/action writing. An intermediary may improve a mapping without making it one-to-one. Arbitrary skill upload or universal disease/disability repair is not demonstrated.\par
\noindent\textit{Organizational level:} Declared system or interface\par
\noindent\textit{Computational granularity:} Declared process/interface/model\par
\noindent\textit{Statistical aggregation:} Validation domain stated independently of physical scale\par


\noindent\textit{Current registry lookup labels:} No additional label declared.\par

% Stable term ID: programmable-synthetic-biology
\begin{samepage}
\subsection*{Programmable synthetic biology}
\noindent Engineering specified biological behaviours by combining design, context-dependent molecular prediction, experimental measurement and iterative correction.\par
\end{samepage}
\noindent DNA/regulation, RNA/cell state and protein/interactions are explanatory layers, not a sufficient literal three-model recipe. Metabolism, feedback, environment and measured function matter. AlphaGenome and AlphaFold do not jointly certify whole-cell programmability; AlphaTranscriptome is a desired capability placeholder here.\par
\noindent\textit{Organizational level:} Declared system or interface\par
\noindent\textit{Computational granularity:} Declared process/interface/model\par
\begin{samepage}
\noindent\textit{Statistical aggregation:} Validation domain stated independently of physical scale\par\nopagebreak[4]

\end{samepage}


\noindent\textit{Current registry lookup labels:} Programmable molecular biology\par

% Stable term ID: public-reaching-recordings
\begin{samepage}
\subsection*{Public reaching recordings}
\noindent Public files pairing cortical electrical events with reaching behavior. The sources use names such as RTT and MC\_RTT.\par
\end{samepage}
\begin{samepage}
\noindent The original model’s release and the separately acquired Indy day have different processing. Their dates and target definitions must remain explicit.\par\nopagebreak[4]

\end{samepage}


\noindent\textit{Current registry lookup labels:} Reaching source archive; RTT; MC\_RTT; MC\textbackslash{}\_RTT; O'Doherty reaching archive; Sabes reaching corpus\par

% Stable term ID: public-source-records-and-document-links
\begin{samepage}
\subsection*{Public source records and document links}
\noindent Repositories such as DANDI and OpenNeuro identify source versions. PDF is a document format, and a URL gives a resource location.\par
\end{samepage}
\begin{samepage}
\noindent Source licenses and credit must be retained. A local paper draft or public input dataset does not mean the project’s new code is published, licensed or submitted to arXiv or a journal.\par\nopagebreak[4]

\end{samepage}


\noindent\textit{Current registry lookup labels:} Public records, document files and links; DANDI; OpenNeuro; PDF; URL; arXiv; public deposition; license; attribution\par

% Stable term ID: published-models-for-brain-scans-and-scalp-recordings
\begin{samepage}
\subsection*{Published models for brain scans and scalp recordings}
\noindent BrainLM learns from fMRI brain-scan data. Neuro-GPT combines an EEG input model with a model of changes over time and is discussed for motor imagery.\par
\end{samepage}
\begin{samepage}
\noindent These are earlier methods for different measurements. The project does not claim to have trained them, make scans and electrical spikes interchangeable, or upload skills.\par\nopagebreak[4]

\end{samepage}


\noindent\textit{Current registry lookup labels:} BrainLM; Published model of scalp signals; Neuro-GPT; NeuroGPT\par

% Stable term ID: published-neural-model-comparison-tasks
\begin{samepage}
\subsection*{Published neural-model comparison tasks}
\noindent The Neural Latents Benchmark, or NLB, provides shared tasks for comparing models of neural recordings.\par
\end{samepage}
\begin{samepage}
\noindent It informs the choice of comparisons. The paper’s own local scores are not official NLB benchmark scores.\par\nopagebreak[4]

\end{samepage}


\noindent\textit{Current registry lookup labels:} Neural Latents Benchmark; NLB\par

% Stable term ID: published-tests-of-changing-recordings
\begin{samepage}
\subsection*{Published tests of changing recordings}
\noindent FALCON is a published benchmark concerned with neural decoding when recordings change and recalibration is limited.\par
\end{samepage}
\begin{samepage}
\noindent The paper uses related source recordings and its own data splits. Its local muscle-signal scores are not official FALCON scores.\par\nopagebreak[4]

\end{samepage}


\noindent\textit{Current registry lookup labels:} FALCON benchmark; FALCON; FALCON M1; FALCON M1-A\par

% Stable term ID: ranges-that-need-statistical-assumptions
\begin{samepage}
\subsection*{Ranges that need statistical assumptions}
\noindent A confidence interval gives a range for a quantity being estimated. A conformal prediction range tries to contain a future actual result at a stated frequency when its required assumptions hold.\par
\end{samepage}
\begin{samepage}
\noindent The transfer figure shows one fixed run. The paper’s ranges are based on earlier prediction errors, with related errors across time; they do not guarantee how often future true values will fall inside.\par\nopagebreak[4]

\end{samepage}


\noindent\textit{Current registry lookup labels:} Confidence interval and conformal coverage; confidence interval; repeated-seed interval; Prediction range with tested coverage; conformal; conformal coverage; prospective-conformal; distribution-free coverage\par

% Stable term ID: recent-signal-model
\begin{samepage}
\subsection*{Recent-signal model}
\noindent A model component that keeps a learned memory of recent inputs as it predicts the next result.\par
\end{samepage}
\begin{samepage}
\noindent Its formal name is gated recurrent unit, or GRU. This established method uses completed input history; its memory does not establish a universal neural language.\par\nopagebreak[4]

\end{samepage}


\noindent\textit{Current registry lookup labels:} Gated recurrent unit; GRU; compact recurrent model; recurrent network\par

% Stable term ID: record-of-what-was-observed
\begin{samepage}
\subsection*{Record of what was observed}
\noindent A record showing which channels or intervals were actually measured and usable.\par
\end{samepage}
\begin{samepage}
\noindent Missing observations stay missing; they do not become zeros. The second Indy source does not supply separate observation intervals for every recorded unit.\par\nopagebreak[4]

\end{samepage}


\noindent\textit{Current registry lookup labels:} Observation mask; observation interval; invalid support; missing observation\par

% Stable term ID: recorded-neural-events
\begin{samepage}
\subsection*{Recorded neural events}
\noindent A spike is an electrical event identified in a recording. A spike train is the sequence of those events over time.\par
\end{samepage}
\begin{samepage}
\noindent A recorded event count is a whole number tied to one channel and time interval. A model’s average expected count can be a fraction and is a different kind of value.\par\nopagebreak[4]

\end{samepage}


\noindent\textit{Current registry lookup labels:} Spike and spike train; neural event; spike count; spike train\par

% Stable term ID: records-linking-files-checks-and-claims
\begin{samepage}
\subsection*{Records linking files, checks and claims}
\noindent A manifest lists specified files and identities; a receipt records checks; an evidence matrix links a claim to its supporting record.\par
\end{samepage}
\begin{samepage}
\noindent These records also keep failures and withdrawn results visible. A catalog listing does not show that every dataset was acquired or used to train a model.\par\nopagebreak[4]

\end{samepage}


\noindent\textit{Current registry lookup labels:} Manifest, receipt and evidence matrix; manifest; source manifest; audit receipt; processing receipt; evidence matrix; withdrawal record\par

% Stable term ID: relay
\begin{samepage}
\subsection*{Relay}
\noindent Pass focused work between agents, then check their results before accepting it.\par
\end{samepage}
\noindent A manual coordination protocol that passes bounded work and verified evidence between specialist sessions, with one integration owner and an independent acceptance reviewer. Each session follows the currently authorized model, effort and speed routing.\par
\noindent A specialist can work on a bounded question while an independent reviewer checks its evidence before the integration owner accepts it. This is a project coordination workflow, distinct from an established external standard or native Ultra mode. Efficiency and quality claims apply only to tested tasks; universal performance equivalence and a 50 percent weekly saving remain undemonstrated.\par
\begin{samepage}
\noindent\textit{Computational granularity:} finite delegated tasks and evidence artifacts\par\nopagebreak[4]

\end{samepage}


\noindent\textit{Current registry lookup labels:} relay; manual Relay protocol; Relay coordination\par

% Stable term ID: rescaled-muscle-signal
\begin{samepage}
\subsection*{Rescaled muscle signal}
\noindent A processed summary of muscle electrical activity put onto a relative scale.\par
\end{samepage}
\begin{samepage}
\noindent The 16 targets have no physical voltage unit. A stale millivolt label in the source does not make them measured voltage, force or muscle activation.\par\nopagebreak[4]

\end{samepage}


\noindent\textit{Current registry lookup labels:} Normalized EMG envelope; EMG envelope; quantile normalization; dimensionless EMG\par

% Stable term ID: rescaling-using-training-data-only
\begin{samepage}
\subsection*{Rescaling using training data only}
\noindent Learning a rescaling rule from the training portion, then applying the same rule to later values.\par
\end{samepage}
\begin{samepage}
\noindent This avoids using local test values to fit that rule. It cannot undo source processing that already used future samples or the whole recording.\par\nopagebreak[4]

\end{samepage}


\noindent\textit{Current registry lookup labels:} Training-only normalization; scaler; target statistics; normalization\par

% Stable term ID: response-after-an-electrical-stimulus
\begin{samepage}
\subsection*{Response after an electrical stimulus}
\noindent An electrical recording made while observing the response to a stimulus.\par
\end{samepage}
\begin{samepage}
\noindent The human preview is a stimulus-response experiment, not voluntary movement. The stimulus itself can also leave an unwanted signal in the recording; that is an artifact.\par\nopagebreak[4]

\end{samepage}


\noindent\textit{Current registry lookup labels:} Evoked potential and stimulus artifact; evoked stimulation; evoked recording\par

% Stable term ID: results-already-seen-during-development
\begin{samepage}
\subsection*{Results already seen during development}
\noindent Results that have been inspected and can therefore influence later decisions.\par
\end{samepage}
\begin{samepage}
\noindent The human 90\% result and local motor test windows have this limitation. They can be reported, but their later reuse is not a new independent test.\par\nopagebreak[4]

\end{samepage}


\noindent\textit{Current registry lookup labels:} Development evidence; development holdout; pilot; descriptive local test; inspected test\par

% Stable term ID: reusing-learning-in-another-recording
\begin{samepage}
\subsection*{Reusing learning in another recording}
\noindent Trying to use information learned in one setting to improve performance in another.\par
\end{samepage}
\begin{samepage}
\noindent The comparison must state the new-setting labels and the controls. The reported frozen motor transfer and unadjusted human cross-day result do not establish successful biological or cross-species transfer.\par\nopagebreak[4]

\end{samepage}


\noindent\textit{Current registry lookup labels:} Cross-recording transfer; transfer; cross-day transfer; cross-species transfer\par

% Stable term ID: routing-a-muscle-around-a-modeled-obstacle
\begin{samepage}
\subsection*{Routing a muscle around a modeled obstacle}
\noindent A rule bends a modeled muscle route around an obstacle. One obstacle shape, a torus, is shaped like a doughnut.\par
\end{samepage}
\begin{samepage}
\noindent The correction restricts three elbow wrapping rules to a stated attachment segment. It removes reproduced route errors without deleting a muscle or claiming a new anatomical measurement.\par\nopagebreak[4]

\end{samepage}


\noindent\textit{Current registry lookup labels:} Muscle wrapping and torus; wrap; wrapping route; torus; ELBOW\_ULNA; ELBOW\textbackslash{}\_ULNA; off-torus contact\par

% Stable term ID: saved-model-and-matching-output-check
\begin{samepage}
\subsection*{Saved model and matching-output check}
\noindent A checkpoint saves a model’s state. A parity check asks whether another execution produces the specified matching output.\par
\end{samepage}
\begin{samepage}
\noindent Repeating a result needs the same support, ordering, conversions, data splits, settings and body-model version. Passing software checks does not certify biological validity.\par\nopagebreak[4]

\end{samepage}


\noindent\textit{Current registry lookup labels:} Checkpoint and parity; checkpoint; replay/checkpoint parity; model parity; reproduction\par

% Stable term ID: scalp-electrical-recording
\begin{samepage}
\subsection*{Scalp electrical recording}
\noindent Electrodes on the scalp record electrical potential; the method is called EEG.\par
\end{samepage}
\begin{samepage}
\noindent This is an aggregate electrical signal, not a list of individual cortical spikes. The project’s paired preview follows electrical stimulation rather than voluntary motor training.\par\nopagebreak[4]

\end{samepage}


\noindent\textit{Current registry lookup labels:} Electroencephalography; EEG; scalp potentials\par

% Stable term ID: settings-needed-to-repeat-a-run
\begin{samepage}
\subsection*{Settings needed to repeat a run}
\noindent The numbers and choices fixing one training run’s starting point and procedure.\par
\end{samepage}
\begin{samepage}
\noindent One fixed run can be reproduced. It does not show how results vary across repeated starts, animals or future recording days.\par\nopagebreak[4]

\end{samepage}


\noindent\textit{Current registry lookup labels:} Random seed and configuration; seed; single-seed; fixed protocol\par

% Stable term ID: shared-internal-model
\begin{samepage}
\subsection*{Shared internal model}
\noindent A model can learn an internal pattern and reuse that pattern when making several predictions.\par
\end{samepage}
\begin{samepage}
\noindent The paper calls the pattern a latent representation and the reusable model a backbone or trunk. The tested shared model and frozen-transfer candidates were rejected.\par\nopagebreak[4]

\end{samepage}


\noindent\textit{Current registry lookup labels:} Latent representation and backbone; latent model; shared latent; backbone; trunk; shared parameters\par

% Stable term ID: simple-weighted-prediction-model
\begin{samepage}
\subsection*{Simple weighted prediction model}
\noindent A linear model that combines inputs using fitted weights while discouraging excessively large weights.\par
\end{samepage}
\begin{samepage}
\noindent The formal method is ridge regression. It provides a meaningful comparison: a more complex model must add value rather than merely produce a plausible display.\par\nopagebreak[4]

\end{samepage}


\noindent\textit{Current registry lookup labels:} Ridge regression; ridge; linear control; linear model\par

% Stable term ID: simplicity-principle
\begin{samepage}
\subsection*{Simplicity principle}
\noindent Introduce one concept or element at a time, using separate slides or successive reveals according to what makes the relationship easiest to understand.\par
\end{samepage}
\noindent Introduce each new element on its own, then show how the elements connect. Use separate slides or successive reveals to make the logical relationship visible. Simplicity cannot remove scientific qualifications, evidence, source rights or required methods. Actual comprehension and exported readability still need review.\par


\noindent\textit{Current registry lookup labels:} No additional label declared.\par

% Stable term ID: skin-prick-suit
\subsection*{Skin-prick suit}
\noindent A suit shows proposed interface contacts across the body.\par
\noindent The proposed body-side (peripheral) interface, drawn as a suit.\par
\noindent A contextual display label in the user's Interoceptomics visual instructions, not a general clinical synonym. The name establishes no acquired neural data, registration, tissue access, function, safety, clinical use, species or physical dimensions.\par

\noindent\textit{Current registry lookup labels:} skin prick suit; peripheral suit; peripheral prosthesis; peripheral interface\par

% Stable term ID: small-connector-model
\begin{samepage}
\subsection*{Small connector model}
\noindent A smaller fitted mapping that connects one recording to an existing model or interface.\par
\end{samepage}
\begin{samepage}
\noindent The connector can change while the older model stays fixed. That does not establish matching neurons or useful transfer without new-setting calibration.\par\nopagebreak[4]

\end{samepage}


\noindent\textit{Current registry lookup labels:} Adapter; linear adapter; recording adapter\par

% Stable term ID: software-that-calculates-a-computer-body-s-motion
\begin{samepage}
\subsection*{Software that calculates a computer body’s motion}
\noindent MuJoCo calculates how a modeled body moves under forces and control inputs.\par
\end{samepage}
\begin{samepage}
\noindent It is established software reused here. Its availability does not prove that the chosen body model or mapping matches a particular animal’s physiology.\par\nopagebreak[4]

\end{samepage}


\noindent\textit{Current registry lookup labels:} MuJoCo; dynamics engine\par

% Stable term ID: spinal-strip
\subsection*{Spinal strip}
\noindent A strip shows the proposed interface along the modeled spine.\par
\noindent The proposed spinal interface, drawn as a strip along the modeled spine.\par
\noindent A contextual display label in the user's Interoceptomics visual instructions, not a general clinical synonym. The name establishes no acquired neural data, registration, tissue access, function, safety, clinical use, species or physical dimensions. It does not validate a modeled dorsal-surface or tailbud overlay as biological spinal-cord or conus anatomy.\par

\noindent\textit{Current registry lookup labels:} spinal strip; spinal strap; spinal prosthesis; spinal interface\par

% Stable term ID: spine2body
\begin{samepage}
\subsection*{Spine2Body}
\noindent A muscle signal becomes a movement through tendons, joints, contact and load. Measure those steps together, and a failed reach becomes a specific question about activation, coordination or mechanics.\par
\end{samepage}
\noindent The pathway through which spinal output reaches peripheral motor nerves and muscles, producing action through the body's mechanics, while sensory signals return from the body.\par
\noindent Surface EMG measures muscle electrical activity, not peripheral nerve firing or force. Spinal plasticity, momentary reflex modulation and peripheral adaptation remain distinct. Simulated movement is a prediction to compare with independent measurements.\par
\begin{samepage}
\noindent\textit{Organizational level:} whole system; network; signal/feature\par\nopagebreak[4]

\end{samepage}


\noindent\textit{Current registry lookup labels:} spine2body; spine to body; spine-to-body\par

% Stable term ID: stopping-training-under-a-stated-rule
\begin{samepage}
\subsection*{Stopping training under a stated rule}
\noindent Ending training when the chosen development rule says to stop, rather than always using every allowed round.\par
\end{samepage}
\begin{samepage}
\noindent Compared methods have equal round limits but can perform different numbers of actual updates. This is not measured biological learning time.\par\nopagebreak[4]

\end{samepage}


\noindent\textit{Current registry lookup labels:} Early stopping; epoch cap; training update\par

% Stable term ID: structured-task-text-and-supported-instructions
\begin{samepage}
\subsection*{Structured task text and supported instructions}
\noindent JSON stores structured text values. The task compiler converts a supported instruction into a defined program.\par
\end{samepage}
\begin{samepage}
\noindent The compiler chooses available recorded building blocks and limited computer-controller practice. It does not program arbitrary biological skills or supply missing task evidence.\par\nopagebreak[4]

\end{samepage}


\noindent\textit{Current registry lookup labels:} JSON and task compiler; JSON; task JSON; natural-language compiler; supported primitive\par

% Stable term ID: surface-electrical-recording-over-the-spine
\begin{samepage}
\subsection*{Surface electrical recording over the spine}
\noindent Surface electrodes placed over spinal regions record electrical potentials; the method is called ESG.\par
\end{samepage}
\begin{samepage}
\noindent These measurements combine electrical contributions. They are not separately identified spinal-neuron spike trains.\par\nopagebreak[4]

\end{samepage}


\noindent\textit{Current registry lookup labels:} Electrospinography; ESG; surface spinal recording; surface ESG\par

% Stable term ID: system-3-thinking
\begin{samepage}
\subsection*{System 3 thinking}
\noindent System 3 thinking is the partnership; a large movement model is the movement-focused model we aim to build within it.\par
\end{samepage}
\noindent The programme’s human–AI partnership framing in which an external supervisory model helps interpret, train or calibrate task-specific neural/body mappings; the movement-focused model being built is called a large movement model.\par
\noindent Expert footage provides movement examples, but it does not supply the depicted experts’ neural or muscle recordings. An external supervisor would help interpret, train or calibrate task-specific neural/body mappings while the large movement model remains the movement-focused model being built within that partnership. Shaw and Nave use System3 for artificial cognition outside the brain in a different research account. This programme’s task-specific framing establishes no new anatomical brain layer, biological neo-neocortex, third anatomical thinking system or achieved universal nervous-system model.\par
\noindent\textit{Organizational level:} Declared game/model/interface role\par
\noindent\textit{Computational granularity:} Declared task-specific model or interface\par
\begin{samepage}
\noindent\textit{Statistical aggregation:} Validation scope declared per experiment; no physical-scale inference\par\nopagebreak[4]

\end{samepage}
\noindent\textit{Source:} Thinking—Fast, Slow, and Artificial\par
\noindent\textit{Source identifier:} shaw-nave2026\par
\noindent\textit{DOI:} 10.2139/ssrn.6097646\par
\noindent\url{https://papers.ssrn.com/sol3/papers.cfm?abstract_id=6097646}\par
\noindent Prior proposed System3 usage for artificial cognition outside the brain; not an established third anatomical brain system\par
\noindent\textit{Source:} Thinking Fast, Slow, Artificially: AI and Your Brain\par
\noindent\textit{Source identifier:} wharton-authors2026\par
\noindent\href{https://executiveeducation.wharton.upenn.edu/thought-leadership/wharton-at-work/2026/05/thinking-fast-slow-and-artificially/}{\nolinkurl{https://executiveeducation.wharton.upenn.edu/}}\\
\href{https://executiveeducation.wharton.upenn.edu/thought-leadership/wharton-at-work/2026/05/thinking-fast-slow-and-artificially/}{\nolinkurl{thought-leadership/wharton-at-work/2026/05/}}\\
\href{https://executiveeducation.wharton.upenn.edu/thought-leadership/wharton-at-work/2026/05/thinking-fast-slow-and-artificially/}{\nolinkurl{thinking-fast-slow-and-artificially/}}\par
\noindent Primary institutional account of Shaw and Nave System3 research\par


\noindent\textit{Current registry lookup labels:} System three thinking; System3 thinking\par

% Stable term ID: testing-on-an-excluded-day-animal-or-species
\begin{samepage}
\subsection*{Testing on an excluded day, animal or species}
\noindent Checking a model on complete identities that were kept out of the relevant development steps.\par
\end{samepage}
\begin{samepage}
\noindent The paper uses formal names such as leave-session-out or leave-species-out. A good score within one session does not answer these broader tests.\par\nopagebreak[4]

\end{samepage}


\noindent\textit{Current registry lookup labels:} Leave-out generalization; leave-session-out; leave-animal-out; leave-species-out; unseen-organism support\par

% Stable term ID: testing-possible-next-actions-in-the-model
\begin{samepage}
\subsection*{Testing possible next actions in the model}
\noindent Simulating near-future actions before choosing the next command.\par
\end{samepage}
\begin{samepage}
\noindent The copied test states must not change the live model. The native result remains rejected and too slow for real-time use; cached playback is a separate function.\par\nopagebreak[4]

\end{samepage}


\noindent\textit{Current registry lookup labels:} Short-horizon planning and lookahead; horizon controller; lookahead; copied-state rollout; model predictive planning\par

% Stable term ID: the-comparison-needed-to-judge-improvement
\begin{samepage}
\subsection*{The comparison needed to judge improvement}
\noindent A defined comparison that shows whether the proposed method adds useful information or performance.\par
\end{samepage}
\begin{samepage}
\noindent A simple predictor, random fixed model and constant muscle command answer different questions. Producing a plausible animation is not a substitute for beating the relevant comparison.\par\nopagebreak[4]

\end{samepage}


\noindent\textit{Current registry lookup labels:} Baseline and control; simple baseline; tonic baseline; random frozen control; behavior-only predictor\par

% Stable term ID: thought2action
\begin{samepage}
\subsection*{Thought2Action}
\noindent Begin with one intention and one action: reach for the cup. Measure the neural, muscular and mechanical steps between them. The aim is to make that intent-to-action relationship as close to one-to-one as possible within the declared task, person, state, context and tolerated error.\par
\end{samepage}
\noindent A task-conditioned mapping from a participant's stated or inferred intent to a specified action, evaluated through independently measured outcomes.\par
\noindent Intent inferred from EEG is not direct access to private thought. Movement constrains competing explanations; it does not certify an achieved universal bijection or the best possible mapping.\par
\begin{samepage}
\noindent\textit{Organizational level:} whole system; network; signal/feature\par\nopagebreak[4]

\end{samepage}


\noindent\textit{Current registry lookup labels:} thought2action; thought to action; thought-to-action\par

% Stable term ID: thought2world
\begin{samepage}
\subsection*{Thought2World}
\noindent Follow the intention all the way to its consequence, then follow sensation and perception back to the next intention. Measuring the middle makes errors easier to locate and the next experiment easier to define.\par
\end{samepage}
\noindent The composed pathway Thought2Action followed by Action2World: intent is translated into an action, and that action produces an outcome that can be independently checked.\par
\noindent This is a proposed programme pathway, not a claim of unrestricted thought control. It is not silently equivalent to the existing Body2World lookup label, Body2Action or Brain2Motion. External touch and movement tasks alone do not solve internal-body sensing or the full Interoceptomics programme.\par
\begin{samepage}
\noindent\textit{Organizational level:} whole system; network; signal/feature\par\nopagebreak[4]

\end{samepage}


\noindent\textit{Current registry lookup labels:} thought2world; thought to world; thought-to-world\par

% Stable term ID: timing-relationships-between-spikes
\begin{samepage}
\subsection*{Timing relationships between spikes}
\noindent How closely spikes can repeat and how spikes from different recorded features occur together.\par
\end{samepage}
\begin{samepage}
\noindent The project checks predicted average counts. It does not establish these finer timing relationships or a complete account of biological spike noise.\par\nopagebreak[4]

\end{samepage}


\noindent\textit{Current registry lookup labels:} Refractoriness and spike synchrony; refractoriness; joint spike synchrony; biological count-noise model\par

% Stable term ID: touch-brain-area
\begin{samepage}
\subsection*{Touch brain area}
\noindent A brain area that processes body-related sensory information. The source study records from somatosensory cortex, including Area 2.\par
\end{samepage}
\begin{samepage}
\noindent The formal shorthand is S1. The selected model estimates an earlier limb state; it does not recover voluntary motor commands or spinal and nerve recordings.\par\nopagebreak[4]

\end{samepage}


\noindent\textit{Current registry lookup labels:} Somatosensory cortex and Area 2; S1; Area 2; Area2\_Bump; somatosensory\par

% Stable term ID: guide-and-learner
\subsection*{Trainer and trainee}
\noindent The trainer supplies guidance; the trainee’s later control is tested. An experienced trainer helping a novice asks a different question from two novices practicing together.\par
\noindent The proposed comparisons include experienced-to-novice, novice-to-novice, individual practice and matched no-link conditions. Later trainee-alone testing removes trainer commands while retaining any bridge required to express movement.\par

\noindent\textit{Current registry lookup labels:} Tutor and pupil; biological tutor; artificial tutor; live tutoring; tutor withdrawal; experienced-to-novice tutor; novice-to-novice; individual practice; Trainer and learner; biological trainer; artificial trainer; live training; trainer withdrawal; experienced-to-novice trainer; Guide and learner; Guide and trainee\par

% Stable term ID: training-a-fresh-model
\begin{samepage}
\subsection*{Training a fresh model}
\noindent Training a new model for the current recording without importing an older model’s learned main weights.\par
\end{samepage}
\begin{samepage}
\noindent This is the scratch specialist. It performs better than the frozen-transfer candidate in both reported label budgets.\par\nopagebreak[4]

\end{samepage}


\noindent\textit{Current registry lookup labels:} Scratch training; scratch specialist; new-recording scratch specialist\par

% Stable term ID: training-with-neural-samples-without-movement-labels
\begin{samepage}
\subsection*{Training with neural samples without movement labels}
\noindent Using neural samples for a training objective that does not require paired movement labels.\par
\end{samepage}
\begin{samepage}
\noindent These samples still provide information and require computation. The first 5\% label condition permits full training-neural support, so it is not only 28.5 seconds of total observation.\par\nopagebreak[4]

\end{samepage}


\noindent\textit{Current registry lookup labels:} Neural-only training support; neural-only pretraining; unlabeled neural support\par

% Stable term ID: trying-to-move-and-actually-moving
\begin{samepage}
\subsection*{Trying to move and actually moving}
\noindent Trying or being cued to move is different from measuring the movement that occurred.\par
\end{samepage}
\begin{samepage}
\noindent The human labels describe attempted finger movement. They do not independently separate cue from intent or provide executed finger trajectories.\par\nopagebreak[4]

\end{samepage}


\noindent\textit{Current registry lookup labels:} Attempted versus executed movement; attempted finger identity; cued attempted movement; executed kinematics\par

% Stable term ID: typical-error-on-the-original-scale
\begin{samepage}
\subsection*{Typical error on the original scale}
\noindent RMS or RMSE squares the deviations, averages them, and takes a square root so the result has the original unit.\par
\end{samepage}
\begin{samepage}
\noindent An error per joint differs from a combined seven-joint error. An error against a changed command goal also differs from an error against the same recorded reference.\par\nopagebreak[4]

\end{samepage}


\noindent\textit{Current registry lookup labels:} Root mean square error; RMSE; RMS; root-mean-square; tracking error\par

% Stable term ID: unassigned-event-channel
\begin{samepage}
\subsection*{Unassigned event channel}
\noindent A channel collecting events not assigned to the sorted features used by the model.\par
\end{samepage}
\begin{samepage}
\noindent These events are excluded from the 171 sorted features in the second Indy recording. They cannot silently become another identified cell.\par\nopagebreak[4]

\end{samepage}


\noindent\textit{Current registry lookup labels:} Residual or hash cell; residual/hash; hash slot; threshold crossing\par

% Stable term ID: units-for-distance-speed-and-angle
\begin{samepage}
\subsection*{Units for distance, speed and angle}
\noindent Metres, centimetres, millimetres and micrometres measure distance. Metres per second measures speed; degrees and radians measure angle. Radians per second (rad/s) measures the rate of change of angle.\par
\end{samepage}
\begin{samepage}
\noindent One metre is 100 cm or 1000 mm; one mm is 1000 µm. Right-angle and varied-angle maze layouts differ in geometry, not in a validated difficulty multiplier. Unit conversion or a small model error does not prove anatomical or decoded-pose accuracy.\par\nopagebreak[4]

\end{samepage}


\noindent\textit{Current registry lookup labels:} Length, velocity and angle units; m; metre; cm; centimetre; mm; millimetre; µm; micrometre; um; m/s; rad; radian; degrees; °; rad/s; orthogonal; nonorthogonal; varied-angle maze; maze layout\par

% Stable term ID: units-for-pushes-turns-and-electrical-potential
\begin{samepage}
\subsection*{Units for pushes, turns and electrical potential}
\noindent Newtons measure force; newton metres measure turning moment; millivolts measure electrical potential. A dimensionless value has no physical unit.\par
\end{samepage}
\begin{samepage}
\noindent The handle record mixes force and moment units. The rescaled muscle targets are dimensionless despite a stale mV label and cannot be read as force or activation.\par\nopagebreak[4]
\noindent\textnormal{Lookup labels:} Force, moment and voltage units; N; newton; N m; N·m; newton metre; mV; millivolt; dimensionless\par
\end{samepage}


\noindent\textit{Current registry lookup labels:} Force, moment and voltage units; N; newton; N m; N·m; newton metre; mV; millivolt; dimensionless\par

% Stable term ID: units-for-time-and-sample-rate
\begin{samepage}
\subsection*{Units for time and sample rate}
\noindent Seconds measure time; milliseconds are thousandths of a second. Hertz counts repetitions or samples per second.\par
\end{samepage}
\begin{samepage}
\noindent A 250 Hz source sample interval is 4 ms; a 500 Hz simulation interval is 2 ms. These rates differ from playback pace, biological travel time and time spent computing.\par\nopagebreak[4]
\noindent\textnormal{Lookup labels:} Time and frequency units; s; second; seconds; ms; millisecond; milliseconds; Hz; hertz; sample rate; 250 Hz; 500 Hz\par
\end{samepage}


\noindent\textit{Current registry lookup labels:} Time and frequency units; s; second; seconds; ms; millisecond; milliseconds; Hz; hertz; sample rate; 250 Hz; 500 Hz\par

% Stable term ID: unwanted-stimulus-signal-in-a-recording
\begin{samepage}
\subsection*{Unwanted stimulus signal in a recording}
\noindent Part of a recording caused by the stimulus or measuring setup rather than the intended response alone.\par
\end{samepage}
\begin{samepage}
\noindent Stimulus artifacts remain in the human previews. Their presence and the evoked task limit interpretation as voluntary motor evidence.\par\nopagebreak[4]

\end{samepage}


\noindent\textit{Current registry lookup labels:} stimulus artifact; artifact\par

% Stable term ID: use-in-a-new-setting-without-labeled-adjustment
\begin{samepage}
\subsection*{Use in a new setting without labeled adjustment}
\noindent Applying a learned model to a new setting without using its labels to adjust the model.\par
\end{samepage}
\begin{samepage}
\noindent The second M1 model uses explicit new-recording labels. Its test score is a result after calibration, not zero-shot success.\par\nopagebreak[4]

\end{samepage}


\noindent\textit{Current registry lookup labels:} Zero-shot transfer; zero-shot\par

% Stable term ID: speed-estimated-from-earlier-positions
\begin{samepage}
\subsection*{Velocity estimated from earlier positions}
\noindent Two recorded fingertip positions can give an estimate of how fast the fingertip moved and in which direction between the samples. The signed components describe velocity; speed is its non-negative magnitude.\par
\end{samepage}
\noindent A position-derived velocity estimate formed from the difference between current and earlier recorded positions divided by their elapsed time; individual components retain sign. Speed is the magnitude of the velocity vector.\par
\begin{samepage}
\noindent In the separate calibrated motor recording, each interval uses the last available fingertip sample and the actual elapsed time to form a velocity target in m/s. This makes the calculation causal with respect to available position samples, not proof that the tracker supplied them without filtering or delay. Upstream filtering and latency remain unestablished; the derived target is not an original, directly supplied velocity trace.\par\nopagebreak[4]

\end{samepage}


\noindent\textit{Current registry lookup labels:} Backward finite difference; backward difference; finite difference; causally derived velocity; 250 Hz position support; Speed estimated from earlier positions\par

% Stable term ID: speed-of-each-model-movement
\begin{samepage}
\subsection*{Velocity of each model movement}
\noindent The model’s stored velocity for its movement degrees of freedom, including direction.\par
\end{samepage}
\begin{samepage}
\noindent The formal variable name is qvel. Its signed generalized velocities depend on the model’s degrees of freedom; speed is a magnitude. Component units and frames depend on the declared coordinate and model. In MuJoCo, quaternion orientations can make qpos and qvel different lengths, so qvel is not an entry-by-entry derivative of qpos.\par\nopagebreak[4]

\end{samepage}


\noindent\textit{Current registry lookup labels:} qvel; generalized velocity; joint speed; Speed of each model movement\par

% Stable term ID: what-a-model-estimated
\begin{samepage}
\subsection*{What a model estimated}
\noindent A model’s estimate of a defined result, such as speed, from the inputs it was given.\par
\end{samepage}
\begin{samepage}
\noindent An estimate remains an estimate even when it matches the test data well. It does not become a new measurement of the spinal cord, nerves or muscles.\par\nopagebreak[4]

\end{samepage}


\noindent\textit{Current registry lookup labels:} Decoded or inferred prediction; decoded; inferred; reconstructed; prediction\par

% Stable term ID: what-an-instrument-recorded
\begin{samepage}
\subsection*{What an instrument recorded}
\noindent A value produced by a stated measuring method, such as an electrode or a movement tracker.\par
\end{samepage}
\begin{samepage}
\noindent Recorded data can contain noise, filtering or artifacts. Calculating muscle length from tracked markers does not turn that length into a direct tissue measurement.\par\nopagebreak[4]

\end{samepage}


\noindent\textit{Current registry lookup labels:} Recorded measurement; measured; observed; recorded signal\par

% Stable term ID: what-the-computer-body-produced
\begin{samepage}
\subsection*{What the computer body produced}
\noindent A value produced while running a computer model, such as a modeled muscle’s activity.\par
\end{samepage}
\begin{samepage}
\noindent Computer-body movement and forces can be checked for numerical errors. They remain separate from measurements of a living body and from evidence that an organism learned.\par\nopagebreak[4]

\end{samepage}


\noindent\textit{Current registry lookup labels:} Simulated quantity; simulated; numerical model state\par

% Stable term ID: what-was-calculated-from-tracking
\begin{samepage}
\subsection*{What was calculated from tracking}
\noindent A value calculated from measurements together with a body model, such as an estimated joint angle.\par
\end{samepage}
\begin{samepage}
\noindent The study’s joint angles and muscle lengths use tracked movement and a model. They are not direct recordings of muscle commands, tissue length or peripheral nerve spikes.\par\nopagebreak[4]

\end{samepage}


\noindent\textit{Current registry lookup labels:} Model-derived quantity; tracking-derived; source-derived; biomechanical covariate\par

% Stable term ID: where-a-value-came-from
\begin{samepage}
\subsection*{Where a value came from}
\noindent The record of a value’s source and the steps used to produce it.\par
\end{samepage}
\begin{samepage}
\noindent For example, a measured hand position, an angle calculated from that position and a computer-generated muscle state have different origins. The software must preserve those differences.\par\nopagebreak[4]

\end{samepage}


\noindent\textit{Current registry lookup labels:} Provenance; lineage; source integrity\par

% Stable term ID: which-recorded-channel-is-which
\begin{samepage}
\subsection*{Which recorded channel is which}
\noindent The labels and order used to identify neural features within one recording.\par
\end{samepage}
\begin{samepage}
\noindent The labels do not identify the same neurons across days. The project did not accept unverified matches between cells in different recordings.\par\nopagebreak[4]

\end{samepage}


\noindent\textit{Current registry lookup labels:} Source-local unit identity; ordered channel identity; cell correspondence; unit order; neuron index\par

% Stable term ID: which-samples-a-value-depends-on
\begin{samepage}
\subsection*{Which samples a value depends on}
\noindent The input history and exact time interval needed to produce a value or target.\par
\end{samepage}
\begin{samepage}
\noindent A single timestamp does not show all of this support. Overlapping histories and invalid samples near data splits need explicit checks.\par\nopagebreak[4]

\end{samepage}

\begin{samepage}
\subsection*{Composed pathways}
\noindent The proposed intermediary has a positive purpose: understand and measure the processes between intent and action so the task-specific mapping can be made as close to one-to-one as possible. The task, participant, state, context, target and tolerated error must be specified and independently tested. This is an aim for identifiability and control, not an achieved universal bijection or a demonstrated optimum.\par
\end{samepage}
\noindent Form \texorpdfstring{\ensuremath{\rightarrow}}{->} function \texorpdfstring{\ensuremath{\rightarrow}}{->} action and intent \texorpdfstring{\ensuremath{\rightarrow}}{->} action describe connected questions. The following proposed pathways separate the links that would need independent tests before their combined prediction could be trusted.\par
\noindent\textit{Biological pathway:} Brain2Spine \texorpdfstring{\ensuremath{\rightarrow}}{->} Spine2Body; resulting in Brain2Body\par
\noindent Relate brain activity to spinal processing, and spinal/peripheral processing to the body; compose the tested links into an intent-to-body-action pathway.\par
\noindent\textit{Functional pathway:} Thought2Action \texorpdfstring{\ensuremath{\rightarrow}}{->} Action2World; resulting in Thought2World\par
\noindent Relate a declared intention to a measured action, and that action to its measured consequence in the task environment.\par
\noindent\textit{Molecular analogy:} DNA \texorpdfstring{\ensuremath{\rightarrow}}{->} RNA; RNA \texorpdfstring{\ensuremath{\rightarrow}}{->} protein; resulting in DNA \texorpdfstring{\ensuremath{\rightarrow}}{->} protein\par
\noindent The intermediate processes explain how an initial specification reaches its downstream product. This analogy illustrates the need for a measured middle; it does not assert identical biology or a context-free one-to-one relationship.\par
\noindent\textit{System 3 engineering route:} Neural activity \texorpdfstring{\ensuremath{\rightarrow}}{->} Musculoskeletal simulation \texorpdfstring{\ensuremath{\rightarrow}}{->} Programmable model\par
\noindent The proposed engineering route to NeuroTranslation, coordinated by the System3 human–AI supervisory model.\par
\noindent Paired source-bound brain/spinal/peripheral/muscle and sensory/outcome observations constrain mechanics and a task-specific learner. A musculoskeletal simulation represents the biomechanical middle. Model read/write capabilities must be qualified individually.\par
\noindent Intent \texorpdfstring{\ensuremath{\rightarrow}}{->} action \texorpdfstring{\ensuremath{\rightarrow}}{->} outcome returns through sensory feedback to perception and subsequent intent; actual observed/inferred/simulated links remain distinguishable.\par
\noindent Interoceptomics combines interoception and connectomics. Touch, proprioception and external sensorimotor tasks can support this programme; they do not by themselves demonstrate internal-body sensing. Combining the proposed pathways advances a research aim rather than establishing a comprehensive solution.\par
\noindent Brain2Body names the wider programme. Body2World remains a contextual lookup label for Body2Action; Brain2Motion and Body2Action retain their own scopes. These related names do not make their measurement targets equivalent to those of the five composed paths.\par
\noindent The physical spinal pathway remains read-only in this proposal. Each claimed recording or writing capability needs its own independent qualification.\par

\noindent\textit{Current registry lookup labels:} Temporal support and context; support; history window; context; next-target support; overlapping histories\par

% Stable term ID: whole-body-interface
% Canonical status: "user_requested_proposed_programme_role"
\subsection*{Whole Body Interface}
\noindent The hand is one part of the question. Whole Body Interface names the wider ambition to connect declared body systems, while showing which pathways have actually been observed and which remain missing.\par
\noindent A proposed extension of brain–body interfacing toward specified observations and outputs across the brain, spinal, peripheral and body systems.\par
\noindent Whole body states an ambition, not achieved coverage. An interface exchanges specified observations or outputs; the existing model spanning the body is a related modelling ambition, not the same object. State the recorded systems, missing pathways, units, clocks, frames and output routes. Separate datasets can support qualified learning and transfer, but they do not become simultaneous measurements of one body. Physiological input/output direction, software input/output and ordinary sensory feedback remain distinct; no spinal writer or physical actuation is authorized.\par


\noindent\textit{Current registry lookup labels:} No additional label declared.\par

% Stable term ID: xenotransgenicism
% Canonical status: "user_requested_working_label_scope_unsettled"
\subsection*{XenoTransgenesis}
\noindent Which genetic feature would be studied, in which organism, and what would count as the relevant change? XenoTransgenesis is the user’s working label for that proposed question within EndoHumanisation.\par
\noindent A working roadmap name for proposed research involving specified genetic components from another species in a model organism.\par
\begin{samepage}
\noindent XenoTransgenesis is the preferred spelling; XenoTransgenicism remains a historical lookup label. This is not an established replacement name for every form of transgenesis, genome editing or humanisation; use the established term appropriate to any actual manipulation. Neither a genetic change nor its name establishes the intended functional effect or human equivalence. No procedure, approved organism, developmental schedule or experimental permission is defined here.\par\nopagebreak[4]

\end{samepage}


\noindent\textit{Current registry lookup labels:} XenoTransgenicism\par

\subsection*{A cross-species graft}

\noindent A chimera has biological contributions from different sources. A graft is added tissue; when it crosses species, it is a xenotransplant. An organoid is an organized tissue model.\par

\noindent Survival, physical connection, activity and useful function are separate questions. The proposed comparisons do not establish a humanized whole brain or an implemented treatment.\par

\begin{samepage}

\noindent Cross-species-graft component only; not a strict synonym for every chimera, graft or organoid, or achieved ExoHumanisation.\par\nopagebreak[4]

\noindent\textit{Additional lookup label:} XenoTransplantation\par

\end{samepage}

\subsection*{Proposed names and explanatory comparisons}

% Proposal ID: neo-neocortex-inspiration

\subsubsection*{Neo-neocortex}

\noindent A metaphor for an external machine contribution to human cognition; no anatomical addition is claimed.\par

% Proposal ID: alpha-transcriptome-capability-placeholder

\subsubsection*{AlphaTranscriptome}

\noindent AlphaTranscriptome is a proposed name for a desired RNA/cell-state prediction capability. It is not an adopted general-purpose model. AlphaGenome already predicts some expression and splicing~\citep{alphagenomeRetained}; other RNA/cell-state models have bounded domains. These precedents leave whole-cell programmability to be independently demonstrated.\par

% Proposal ID: alpha-cell-supplied-passage-concept

\subsubsection*{AlphaCell}

\noindent A label for integrating cellular layers appears in the supplied passage; no implementation, novelty or product entitlement follows.\par

% Proposal ID: connetype-individual-profile-proposal

\subsubsection*{Connetype}

\noindent A connection profile would describe the selected observations for one individual. Its anatomical coverage, resolution, method and uncertainty must remain visible. The profile does not identify that person’s unique function, thoughts or a fixed biological type.\par

\noindent An individual’s specified neural-connection profile under a declared measurement and representation.\par

\noindent The original three-name line ends in a question mark; all three remain tentative. No formal empirical definition or repeatability threshold is fixed.\par

\noindent The tentative name is Connetype. Connectype remains an unapproved editorial alternative; it is neither a lookup alias nor an adopted spelling.\par

% Proposal ID: interotype-individual-profile-proposal

\subsubsection*{Interotype}

\noindent A heartbeat-related observation can vary with state and measurement conditions. A proposed individual profile must state what was observed and how repeatability would be assessed; it cannot treat touch or proprioception as interchangeable with interoception.\par

\noindent An individual’s profile of specified internal-body sensing observations under declared conditions.\par

\noindent The original three-name line ends in a question mark; all three remain tentative. No formal empirical definition or repeatability threshold is fixed.\par

\noindent Interotype is a tentative name, rather than an established scientific type.\par

% Proposal ID: enactotype-individual-profile-proposal

\subsubsection*{Enactotype}

\noindent A profile of selected actions is bounded by the tasks, observations and conditions used to construct it. Repeatability, state dependence and generalization remain empirical questions; a stable innate action type or inferred unique intent is not established.\par

\noindent An individual’s profile of specified actions and action–outcome observations in declared tasks.\par

\noindent The original three-name line ends in a question mark; all three remain tentative. No formal empirical definition or repeatability threshold is fixed.\par

\noindent Enactotype is a tentative name, rather than an established scientific type.\par

% Proposal ID: form-function-feature-explanatory-lens

\subsubsection*{Form, Function, Feature}

\noindent Start with a structure, ask what it does in the stated setting, then name the observable feature being compared. A hand’s form, its grasping function and a measured property of a particular grasp answer related but different questions.\par

\noindent Feature must identify the actual observed or defined trait in its context; it is not a new universal third biological layer or a substitute for function or action. The existing Form \(\rightarrow\) function \(\rightarrow\) action account and intent–action distinction remain unchanged. Structure constrains function; context, history, selection and co-option prevent a guaranteed one-to-one or teleological chain.\par

% Proposal ID: semiosis-sensorimotor-explanatory-analogy

\subsubsection*{Semiosis as a sensorimotor analogy}

\noindent A drawing can resemble a hand, a pointing finger can indicate an object, and an instruction can follow a learned convention. Such signs can guide action through different relations; they do not assign perception, sensation and action to three exclusive boxes.\par

\noindent Icon, index and symbol describe sign–object relations, while sensation, perception and action name processes. A sign can combine relations. No one-to-one biological correspondence, neural code, private-experience identity or strict isomorphism is established. Outside-in and inside-out perspectives must state whether they concern physiology, modelling or interpretation.\par

% Proposal ID: agent-vs-neurology-stage5-proposal

\subsubsection*{Agent vs Neurology}

\noindent The roadmap ends by asking what an agent could contribute to a defined neurological problem. Agent vs Neurology retains that proposed destination; the specific problem, task and comparison still have to be chosen.\par

\noindent An unresolved Stage 5 research question, not a selected disorder, intervention, diagnostic system or clinical trial.\par

\noindent No specific disease, patient population, endpoint, treatment mechanism, universal cure or successful clinical benefit is inferred. Any future study must define those matters and its ordinary feedback or separately qualified physical interface under the original scientific, participant, rights, safety and independent gates. A prior software, metadata or model result cannot supply this evidence.\par

\noindent Agent vs Neurology names Stage 5 of the proposed roadmap.\par

% Source-key boundary measured separately in the retained technical and explained tails.
% Preserve natbib/hyperref item handling, order, font, and positive 7pt bibsep.
\begingroup
\let\NTlayoutOriginalBibitem\bibitem
\def\NTlayoutBalanceKey{tsai2019}
\renewcommand{\bibitem}[2][]{%
  \clubpenalty=10000\widowpenalty=10000\displaywidowpenalty=10000
  \def\NTlayoutCurrentKey{#2}%
  \ifx\NTlayoutCurrentKey\NTlayoutBalanceKey
    \par\newpage
  \fi
  \NTlayoutOriginalBibitem[#1]{#2}%
}
\bibliographystyle{unsrtnat}\bibliography{references,myohand-numerical-citations-common}
\endgroup
\clearpage
\appendix
% Source-only shared appendix. No data, apparatus, reference-series or output acceptance.
\section{Reference Comparison Protocol: Unfilled Record}
\label{app:reference-comparison-protocol}

A moving hand gives the instruments something in common to answer to. Scalp
EEG would observe electrical activity at declared scalp sites; surface EMG
would observe muscle electrical activity at declared arm sites. Neither tells
us, by itself, where the hand went. The proposed LEFT and RIGHT tripod views
would provide complementary movement observations. Each acquisition needs its
own identifiable series of real references, followed by an independent
comparison of the assembly. Two streams can agree because they share the same
mistaken calibration.

\textbf{This is an unfilled protocol and record.} No observations, admitted
reference series, measured errors or acceptance decisions have been entered.
The apparatus and its anatomical and functional claims remain unvalidated.
The older above/below maze illustration retains its separate schematic
context, perpendicular-to-floor geometry and material-path questions. Its
directions are not synonyms for the current LEFT/RIGHT tripod arrangement.
Actual positions, viewing paths, coverage and transmission would need evidence.

\subsection*{Identify the body and the observation}

A photograph can constrain a visible proportion; it cannot supply an EEG
channel or a calibrated trajectory. Record each reference's provider, stable
identifier, version or acquisition context, checksum if materialized, license,
access terms and supporting claim. Identify the organism, developmental stage,
participant or specimen, session, body part, task and view or time interval.
Record physical dimensions and units, posture, support and contact, so that the
comparison concerns the same body and movement. Keep those dimensions separate
from computational detail and from the rule used to average errors. Existing
datasets, atlases and photographs retain their separate study identities;
none is silently admitted for this apparatus.

EEG needs sites, reference, channel order and electrical units. Surface EMG
needs named muscles, sleeve and contact identities, reference and electrical
units; it supplies neither incoming sensory-nerve traffic nor direct force.
Each tripod view needs its selected modality and camera or ranging geometry.
Optical triangulation and LiDAR ranging measure different things. Preserve raw
observations, derived features and model estimates separately. These recordings
do not identify hidden spinal activity, individual neurons or a unique brain state.

\subsection*{Make time and position comparable}

A reach can follow the measured path and still arrive late. Record source
timestamps, clock identities, cadence and units, synchronization events,
offsets, drift, dropped samples and interpolation. Retain calibration
observations, frame transformations, their uncertainty and checks against an
independent reference. Report offline processing and end-to-end latency
explicitly. An alignment error must not disappear inside a movement score,
and no real-time capability follows from this protocol.

Compare scalp EEG, each arm's surface EMG, each tripod acquisition, and the
time/frame calibration and any actual material path against their respective
series of real references. Keep overlap, occlusion, dropout and failed views
alongside anatomy, proportion, support and contact. Preserve both motion
observations and the joined map. A missing finger remains missing unless its
replacement is labelled as an estimate. The assembled map then needs its own
independent movement comparison; agreement among its inputs cannot establish
the accuracy of their common reconstruction.

\subsection*{Test what the assembly claims}

Independently measured movement can test a virtual body's path. Contact and
force claims need appropriate measurements of contact and force; a lifelike
image cannot establish valid mechanics. A learner's prediction needs a
withheld measured target, declared input availability, baselines, negative
controls, independent held-out participants, sessions, tasks or species where
applicable, and leakage checks. Keep the reserved recordings closed and the
original tolerances unchanged. This record admits no new fitting, biological
teaching, participant acquisition, stimulation, physical intervention or
model promotion. Physical spinal scope remains read-only.

For each component and the assembled apparatus, retain matched reference
views or time intervals, the target observable, result, error with units,
repeatability, uncertainty and failures. Declare the applicable original
acceptance rule, its source, domain, comparator and tolerances before evaluation.
The programme's unchanged real-reference contract and every stronger original
scientific, anatomical, functional, physical, held-out, rights, visual, runtime
and export criterion still govern their respective claims. No universal
numeric threshold is introduced or relaxed. Biological function and clinical
readiness need evidence beyond a matched motion record.

\subsection*{Unfilled comparison record}

Use one record per reference and comparison, retaining a series for every
component and for the assembled map. All values below are \texttt{null};
reference-series lists are empty and every acceptance state is \texttt{HOLD}.

\begin{description}
\item[Source and rights.] Record ID, component or assembly, provider, source
identifier, version or acquisition context, local checksum if materialized,
license, access and supporting claim.
\item[Body and task.] Organism, developmental stage, participant or specimen,
session, body part, task, view or time interval, physical dimensions and units,
posture, support, contact, actual sensor geometry and material path.
\item[Acquisition.] Modality, measured quantity and units, sites, montage,
reference, channel and contact identities, source frame and clock, timestamps,
cadence, synchronization events, offsets, drift, calibration observations and
transforms, calibration uncertainty, interpolation, missingness, occlusion,
offline processing and end-to-end latency.
\item[Independent comparison.]\begin{samepage} Component-series and assembly-series identifiers,
independent reference and its acquisition, aligned interval, raw/derived/simulated
provenance, input availability, held-out split and custodian, baseline, negative
controls, leakage checks and target observable.
\item[Result and decision.] Result, error and units, repeatability, uncertainty,
failure cases, applicable original acceptance rule and source, domain,
comparator, tolerances, independent reviewer, evidence and verdict.
\par
\end{samepage}
\end{description}

No measured series or scientific validation is supplied by this empty record.
Rights, physical, clinical, runtime, native, render, export and publication
acceptance remain separate and open; no hosted release is claimed.

\end{document}
